% DER_XVIII-b: Grothendieck's text with the commentary set in, then the overview, inside a group that applies the chunk's own definitions
\dgrpart{1}{XVIII-b}{Part 2 of 5, manuscript pages 127--228}{XVIII-b\enspace\textperiodcentered\enspace Chapter XVIII, part 2 of 5\enspace\textperiodcentered\enspace manuscript pages 127--228}{Chapter XVIII \textperiodcentered\ \textup{XVIII-b}}
\dgrchunkbegin
% DER_XVIII-b: the definitions of the chunk's preamble, applied inside its group
\let\printedpage\relax\newcommand{\printedpage}[1]{\marginpar{\footnotesize p.~#1}}
\let\msp\relax\newcommand{\msp}[1]{\marginpar{\footnotesize [#1]}}
\let\Filt\relax\newcommand{\Filt}{\operatorname{Filt}}
\let\card\relax\newcommand{\card}{\operatorname{card}}
\let\Fl\relax\newcommand{\Fl}{\operatorname{Fl}}
\let\Ob\relax\newcommand{\Ob}{\operatorname{Ob}}
\let\Ind\relax\newcommand{\Ind}{\operatorname{Ind}}
\let\cardred\relax\newcommand{\cardred}{\mathop{\mathrm{card}\,\textup{réd}}\nolimits}
\let\circled\relax\newcommand{\circled}[1]{\tikz[baseline=(c.base)]\node[draw,circle,inner sep=1pt](c){#1};}
\let\Mc\relax\newcommand{\Mc}{\mathcal{M}}
\let\U\relax\newcommand{\U}{\mathfrak{U}}
\let\Hom\relax\newcommand{\Hom}{\operatorname{Hom}}
\let\End\relax\newcommand{\End}{\operatorname{End}}
\let\Kar\relax\newcommand{\Kar}{\operatorname{Kar}}
\let\red\relax\newcommand{\red}{\operatorname{\text{réd}}}
\let\uHom\relax\newcommand{\uHom}{\underline{\operatorname{Hom}}}
\let\Ens\relax\newcommand{\Ens}{\mathrm{Ens}}
\let\cM\relax\newcommand{\cM}{\mathcal{M}}
\let\cD\relax\newcommand{\cD}{\mathcal{D}}
\let\defeq\relax\newcommand{\defeq}{\overset{\text{déf}}{=}}
\let\rbox\relax\newcommand{\rbox}[1]{\tikz[baseline=(X.base)]\node[draw,rounded corners=4pt,inner sep=2pt](X){\textup{#1}};}

% Reader layout additions; no mathematical source readings changed here.

\clearpage

% ---- from der-ch18_p041-045.tex
\setcounter{footnote}{49}

\printedpage{39}
% head: A. Grothendieck : Dérivateurs, Ch. XVIII, Édition des Universités de Montpellier II et Paris VII

\begin{itemize}
\item[\textit{1)}] $\Filt'^{\pi}$ \textit{est stable dans $\mathcal{M}$ par} $\varinjlim$ \textit{de cardinal $\leqslant \pi$.}
\item[\textit{1$'$)}] $\Filt'^{\pi}$ \textit{est stable par conoyaux de double-flèches.}
\item[\textit{2)}] \textit{Tout objet $x$ de $\mathcal{M}$ est de la forme $\varinjlim\limits_{I} x_i$, où $x_i$ dans} $\Filt'^{\pi}$\textit{, $I$ grande devant $\pi$.}
\item[\textit{3)}] \textit{La condition d) de 1°) implique $P(\pi)$ (donc lui est équivalente), i.e. on a}
\[ \mathcal{M}_{\pi} \ = \ \Filt'^{\pi}. \]
\end{itemize}

\msp{page 127}
\textit{Ces conditions sont satisfaites dans les trois cas suivants :}
\begin{itemize}
\item[\textit{A)}] \textit{$\pi^{\pi_0} = \pi$ (pour mémoire, c'est contenu dans 2°).}
\item[\textit{B)}] \textit{Les familles épimorphiques strictes dans $\mathcal{M}$ sont épimorphiques strictes universelles.}
\item[\textit{C)}] \textit{Les limites inductives filtrantes dans $\mathcal{M}$ commutent au changement de base (ce qui est le cas notamment si les} $\varinjlim$ \textit{filtrantes sont exactes à gauche)} (\footnote{Je suppose que cette condition revient à la même condition sur $C$, pour les $\varinjlim$ filtrantes de cardinal $\leqslant \pi$.}).
\end{itemize}

\textbf{Remarque.} La filtration de $\mathcal{M}$ pour les $\Filt'^{\pi}$ n'est autre que celle envisagée dans SGA 4 I 9.13, où on fait de plus l'hypothèse B) sur $\mathcal{M}$, qui après-coup n'était plus tellement de mon goût. Toujours est-il que d'après (4.5.13.1), cette filtration se rapproche plus de la ``bonne'' filtration par les $\mathcal{M}_{\pi}$, que celle que j'avais eu encore naguère tendance à lui préférer, celle par les $\Filt^{\pi}$. Elle a d'ailleurs sur cette dernière l'avantage, dans le cas envisagé B), d'être une filtration cardinale (ce qui n'est pas nécessairement le cas pour $\Filt^{\pi}$, du moins s'il existe des cardinaux infinis $\pi \geqslant \pi_0$ avec $\pi^{\pi_0} \neq \pi$). Je soupçonne maintenant qu'en général (et s'il existe des cardinaux infinis $\pi > \pi_0$ tels que $\pi^{\pi_0} \neq \pi$), en dehors des conditions A) B) C), la filtration par les $(\Filt'^{\pi})$ n'est \textit{pas} une
\msp{page 128}
filtration cardinale. En tous cas, 4.5.6 et la présente proposition, 3°, montrent que si soit $(\Filt^{\pi})_{\pi\geqslant\pi_2}$, soit $(\Filt'^{\pi})_{\pi\geqslant\pi_2}$ est une filtration cardinale ($\pi_2$ étant un cardinal donné $\geqslant \pi_1$), alors elle coïncide avec la filtration par les $\mathcal{M}_{\pi}$. Donc apparemment, dans tous les cas où je sais construire une filtration cardinale, sur un paratopos du moins (i.e. sur une catégorie $\mathcal{M}$ stable par petites $\varinjlim$ , et dont tous les objets sont accessibles -- l'existence d'une filtration cardinale impliquant de plus celle d'une petite sous-catégorie strictement génératrice \dots), celle-ci n'est autre que celle par les $\mathcal{M}_{\pi}$ -- chose dont je ne m'étais par aperçu en rédigeant SGA 4 I 9. % [sic] m'étais par aperçu -> m'étais pas aperçu, "par" for the negation "pas"

\textsc{Démonstration de 4.5.13.}

\textit{Prouvons 1°.} On a les implications tautologiques
\[
\begin{tikzcd}[row sep=large]
\mathrm{b}'' \arrow[r, Rightarrow] \arrow[d, Rightarrow] & \mathrm{b}' \arrow[r, Rightarrow] \arrow[d, Rightarrow] & \mathrm{b} \arrow[d, Rightarrow] \\
\mathrm{a}'' \arrow[r, Rightarrow] & \mathrm{a}' \arrow[r, Rightarrow] & \mathrm{a} \text{ ,}
\end{tikzcd}
\qquad\qquad \text{(\footnotemark) ,}
\]
\footnotetext{compte tenu que $C$ est strictement génératrice dans $\mathcal{M}$, donc $x \simeq \varinjlim\limits_{i\in C/x} x_i$.}
et pour prouver l'équivalence de ces six conditions, il suffit de prouver
\[ \mathrm{a} \ \Longrightarrow \ \mathrm{b}'' \ . \]
% foot: Version 17 décembre 2025

\printedpage{40}
% head: A. Grothendieck : Dérivateurs, Ch. XVIII, Édition des Universités de Montpellier II et Paris VII

\msp{page 129}
Si on a
\[ x \ \simeq \ \varinjlim_{I} x_i \ , \qquad\quad \card \Fl I \ \leqslant \ \pi \ , \]
soit $I'$ l'ensemble des \textit{sous-diagrammes} $J \subset I$ finis de $I$, ordonné par inclusion, donc $I'$ est un ensemble ordonné filtrant, et du fait que $C$ est stable par $\varinjlim$ finies, on trouve un foncteur
\[ J \longmapsto x_J \ = \ \varinjlim_{i\in J} x_i \ : \ I' \longrightarrow C \ , \]
et on aura
\[ x \ \simeq \ \varinjlim_{J\in I'} x_J \ . \]
D'autre part, on a
\[
\begin{aligned}
\card I' \ &= \ \card \mathfrak{P}_f( \ \underbrace{\Fl(I)}_{\text{de card $\leqslant \pi$}} \ ) \\
&\leqslant \ \pi + \pi^2 + \ \cdots \ + \pi^n + \ \dots \\
&\leqslant \ \pi\aleph_0 \ = \ \pi \ ,
\end{aligned}
\]
q.e.d. Pour établir 1°, il reste à établir
\[
\begin{tikzcd}[column sep=small]
& P(\pi) \arrow[dl, Rightarrow] \arrow[dr, Rightarrow] & \\
\mathrm{c} & & \mathrm{d} \text{ ,}
\end{tikzcd}
\]
et l'implication inverse $\mathrm{c} \Longrightarrow P(\pi)$ quand $\pi \geqslant \pi_1$. Mais $P(\pi) \Longrightarrow \mathrm{c}$ est une tautologie, et $P(\pi) \Longrightarrow \mathrm{d}$ puisque $\mathcal{M}_{\pi}$ est stable par $\varinjlim$ de cardinal $\leqslant \pi$. Reste à prouver que $\mathrm{c} \Longrightarrow P(\pi)$ si $\pi \geqslant \pi_1$.

\dgrblock{dgrb998}{The intrinsic size of an object}
\begin{dcomment}{exposition}{p128}{The intrinsic size of an object}
Grothendieck is trying to replace a convenient measure of the size of a presentation by a measure belonging to the object itself. This is the reason for the repeated comparison of $\Filt^{\pi}$, $\Filt'^{\pi}$ and $\mathcal M_\pi$. The first two remember how an object has been assembled from a chosen small category $C$; the third asks which diagrams a represented functor can see. A change of presentation should not alter the last answer.

There is a shift of convention to make before using modern references. In these pages ``grand devant $\pi$'' means that every diagram having at most $\pi$ arrows admits a cocone. Thus, with $\lambda=\pi^+$, this is $\lambda$-filteredness, and an object called $\pi$-accessible here is $\lambda$-presentable in the usual terminology. No regularity of $\pi$ is being assumed. The successor $\pi^+$ is regular, and it is this regular cardinal which enters the standard interchange theorem for sets. We shall retain the manuscript's indices while making this translation whenever a modern formulation is needed.

The essential calculation is
\[
 \Hom\left(\varinjlim_{j\in J}X_j,\varinjlim_{i\in I}Y_i\right)
 =\varprojlim_{j\in J}\Hom\left(X_j,\varinjlim_{i\in I}Y_i\right).
\]
If $|\Fl J|\leq\pi$, $I$ is $\pi^+$-filtered and every $X_j$ is $\pi$-accessible, the right side becomes
$\varprojlim_j\varinjlim_i\Hom(X_j,Y_i)$ and then
$\varinjlim_i\varprojlim_j\Hom(X_j,Y_i)$. The last interchange holds in sets: at most $\pi$ representatives can be put at one stage, and at most $\pi$ required equalities can be made true at a later stage. This proves closure of $\mathcal M_\pi$ under the indicated colimits. It also explains why relations, and not merely generators, enter a size bound.

The historical comparison with \citet[Expos\'e~I, \S9]{SGA4I1972} concerns this same passage from bounded presentations to bounded represented functors. The manuscript's provisional claim that one of the presentation filtrations is automatically the intrinsic filtration must therefore be read with its hypotheses. The following pages repeatedly repair precisely that identification. The relevant modern recognition theorem is the characterization by a small strong generator of presentable objects, recorded as Theorem~1.20 in \citet{AdamekRosicky1994}; the word ``strong'' carries mathematical content.
\end{dcomment}


\msp{page 130}
Supposons donc qu'on ait une famille épimorphique stricte
\[ x_{\alpha} \longrightarrow x \qquad\quad (\alpha \in A,\ \card A \leqslant \pi) \ , \]
cela signifie donc que $x$ est limite inductive d'un système inductif indexé par la catégorie $B$ ayant comme ensemble d'objets $A \amalg (A\times A)$, les flèches non identiques étant représentées par les flèches canoniques
\[
\begin{tikzcd}[column sep=small]
& x_{\alpha} \times_x x_{\beta} \arrow[dl, "p_{\alpha,\beta}"'] \arrow[dr, "q_{\alpha,\beta}"] & \\
x_{\alpha} & & x_{\beta} \text{ .}
\end{tikzcd}
\]

\dgrblock{dgrb999}{Finite pieces of a diagram}
\begin{dcomment}{exposition}{p129}{Finite pieces of a diagram}
Grothendieck wants to make an arbitrary bounded colimit into an ordinary filtered colimit without enlarging its cardinal. His argument works if a finite subdiagram is a finite set of constraints. It would not work if it meant a finite subcategory closed under composition: one endomorphism can generate infinitely many arrows.

Here is a precise construction which also makes the counting transparent. For a finite set $V$ of objects of $I$ and a finite set $E$ of arrows whose endpoints lie in $V$, put
\[
 X_{V,E}=\operatorname{coeq}\left(
 \coprod_{f\in E}X_{s(f)}\rightrightarrows
 \coprod_{v\in V}X_v\right).
\]
The two maps use respectively the inclusion of the source summand and the given arrow followed by the inclusion of its target. These are finite colimits and therefore belong to $C$. Enlarging $(V,E)$ induces a canonical map of coequalizers. The ordered set of these pairs is filtered, has cardinal at most $\pi$ when $|\Fl I|\leq\pi$, and its colimit is $\varinjlim_I X_i$. Indeed, a compatible family of maps from all $X_{V,E}$ to an object $Y$ is exactly a family of maps $X_i\to Y$ satisfying every arrow relation of $I$. This proves the asserted universal property, including diagrams with loops.

The calculation illustrates an advantage of the primed filtration. It allows arbitrary bounded presentations to be flattened into a single bounded filtered presentation when finite colimits in $C$ are available. It does not replace ordinary filteredness by $\pi_0^+$-filteredness. The latter requires simultaneous solutions of much larger families of constraints and is the source of the cardinal exponent appearing later.
\end{dcomment}


(\textbf{N.B.} $B$ n'est par une catégorie ordonnée, car pour l'objet $(\alpha,\alpha)$, représenté par $x_{\alpha}\times_x x_{\alpha}$, % [sic] n'est par une -> n'est pas une, "par" for the negation "pas"
il y a \textit{deux} flèches dans l'objet $\alpha$, représenté par $x_{\alpha}$. Mais on a évidemment $\card \Fl(B) \leqslant \pi$.) Si on savait que $C$ est stable par produits binaires et par passage à des sous-objets (ici les $x_{\alpha} \times_x x_{\beta} \ \subset \ x_{\alpha} \times x_{\beta}$), on aurait prouvé la condition a). À défaut, utilisons encore le fait que $C$ est génératrice par épimorphismes (même stricts) et choisissons pour tout couple $(\alpha,\beta)$ un ensemble épimorphique de flèches :
\[ \underbrace{x_{\alpha,\beta,\gamma}}_{\in \Ob C_0} \xrightarrow{u_{\alpha,\beta,\gamma}} x_{\alpha} \times_x x_{\beta} \ , \]
% foot: Version 17 décembre 2025

\printedpage{41}
% head: A. Grothendieck : Dérivateurs, Ch. XVIII, Édition des Universités de Montpellier II et Paris VII
où $C_0$ est une sous-catégorie réduite exhaustive de $C$. La flèche $u_{\alpha,\beta,\gamma}$ est connue quand on connaît les composés
\msp{page 131}
\[
\begin{tikzcd}
& x_{\alpha} \\
x_{\alpha,\beta,\gamma} \arrow[ur, "v_{\alpha,\beta,\gamma}"] \arrow[dr, "w_{\alpha,\beta,\gamma}"'] & \\
& x_{\beta} \text{ ,}
\end{tikzcd}
\]
donc pour $x_{\alpha,\beta,\gamma} \in \Ob C_0$ fixé, les choix possibles pour $u_{\alpha,\beta,\gamma}$ sont de cardinal majoré par $\pi_1\pi_1 = \pi_1$, de plus on a $\card \Ob C_0 \leqslant \pi_1$, donc le nombre total des flèches $v_{\alpha,\beta,\gamma}$, $w_{\alpha,\beta,\gamma}$ (pour $\alpha$, $\beta$ fixés) est majoré par $\pi_1 \cdot \pi_1 = \pi_1$. Donc le nombre total des flèches $u_{\alpha,\beta,\gamma}$, pour $\alpha$, $\beta$, $\gamma$ variables, est majoré par $\pi \cdot \pi_1 = \pi$ (si $\pi \geqslant \pi_1$). Donc $x$ est représenté par une catégorie $I$ ayant comme objets l'ensemble somme de $A$ et de l'ensemble des indices triples $\alpha,\beta,\gamma$, et pour chaque tel indice, exactement deux flèches $v_{\alpha,\beta,\gamma}$, $w_{\alpha,\beta,\gamma}$ non identiques de source cet objet, et allant respectivement dans $x_{\alpha}$, $x_{\beta}$. On a donc $\card \Fl I \leqslant \pi$, et on a un foncteur $I \longrightarrow C$, $\alpha \longmapsto x_{\alpha}$, $(\alpha,\beta,\gamma) \longmapsto x_{\alpha,\beta,\gamma}$, tel que $x \ = \ \varinjlim\limits_{i\in I} x_i$. Cela achève la démonstration de 1°).

\msp{page 132}
\textit{Prouvons 2°.} La première inclusion (4.5.13.1) est tautologique, la deuxième ne fait qu'exprimer $P(\pi) \Longrightarrow \mathrm{d}$. Le fait que $\pi^{\pi_0} = \pi$ implique l'égalité des membres extrêmes est un rappel de 4.5.6. Reste à prouver que $\Filt'^{\pi}$ est stable dans $\mathcal{M}$ par sommes de cardinal $\leqslant \pi$. Mais si $(x_{\alpha})_{\alpha\in A}$ est une famille d'objets de $\Filt'^{\pi}$, avec $\card A \leqslant \pi$, et si $x = \coprod_{\alpha\in A} x_{\alpha}$, prenant pour tout $x_{\alpha}$ un foncteur $\varphi_{\alpha} : I_{\alpha} \longrightarrow C$, avec $\card \Fl I_{\alpha} \leqslant \pi$, tel que $x_{\alpha}$ soit la $\varinjlim$ dans $\mathcal{M}$ de ce foncteur, on trouve que $x$ est $\varinjlim$ du foncteur $\varphi : I \longrightarrow C$ défini par les $\varphi_{\alpha}$, où $I = \coprod_{\alpha} I_{\alpha}$ (somme dans Cat). Il est évident que $\card \Fl(I) \leqslant \pi\cdot\pi = \pi$, ce qui achève de prouver 2°.

\textit{Prouvons 3°.} On a les implications tautologiques $1 \Longrightarrow 1'$, et $1' \Longrightarrow 1$, puisque $\Filt'^{\pi}$ est stable par sommes de cardinal $\leqslant \pi$. Donc $1 \Longleftrightarrow 1'$, et implique que pour tout $x$ dans $\mathcal{M}$,
\msp{page 133}
$I \ = \ \Filt'^{\pi}/x$ est stable par $\varinjlim$ de cardinal $\leqslant \pi$, donc est grand devant $\pi$. Comme $x \simeq \varinjlim\limits_{i\in I} x_i$, cela prouve que $1 \Longrightarrow 2$. L'implication $2 \Longrightarrow 3$ est standard, et résulte de $\Filt'^{\pi} \subset \mathcal{M}_{\pi}$, je vais expliciter dans un

\textbf{Lemme 4.5.14.} \textit{Soient $\mathcal{M}$ une catégorie stable par petites} $\varinjlim$ \textit{, $\pi$ un cardinal infini, $\mathcal{M}_0$ une sous-catégorie strictement pleine de $\mathcal{M}$. Considérons les conditions}
\begin{itemize}
\item[\textit{a)}] \textit{$\mathcal{M}_0$ est une sous-catégorie strictement génératrice, et stable dans $\mathcal{M}$ par} $\varinjlim$ \textit{de cardinal $\leqslant \pi$.}
\item[\textit{b)}] \textit{Pour tout objet de $\mathcal{M}$, on a un isomorphisme}
\[ x \ \simeq \ \varinjlim_{i\in I} x_i \ , \]
\textit{les $x_i$ dans $\mathcal{M}_0$, $I$ grande devant $\pi$.}
% foot: Version 17 décembre 2025

\printedpage{42}
% head: A. Grothendieck : Dérivateurs, Ch. XVIII, Édition des Universités de Montpellier II et Paris VII
\item[\textit{b$'$)}] \textit{Comme b), avec $I$ catégorie} ordonnée.
\item[\textit{c)}] \textit{$\mathcal{M}_{\pi} \subset \mathcal{M}_0$ (où $\mathcal{M}_{\pi}$ est la sous-catégorie strictement pleine de $\mathcal{M}$ formée des objets $\pi$-accessibles).}
\end{itemize}
\textit{On a donc les implications}
\[ a \ \Longrightarrow \ b \ \Longleftrightarrow \ b' \ \Longrightarrow \ c \ , \]
\textit{et si on suppose $\mathcal{M}_0 \subset \mathcal{M}_{\pi}$, et $\mathcal{M}_{\pi}$ sous-catégorie strictement génératrice de $\mathcal{M}$, alors les quatre conditions $a$, $b$, $b'$, $c$ sont équivalentes, et équivalent à $\mathcal{M}_0 = \mathcal{M}_{\pi}$} (\footnote{cf. démonstration de proposition 4.5.6, $4 \Longrightarrow 5$, page 117.}).

\msp{page 134}
Enfin (revenant à la démonstration de 4.5.13, 3°), l'implication $3) \Longrightarrow 1)$ provient du fait que $\mathcal{M}_{\pi}$ est stable par $\varinjlim$ de cardinal $\leqslant \pi$.

\dgrblock{dgrb1000}{The retract which a filtered presentation supplies}
\begin{dcomment}{development}{p133}{The retract which a filtered presentation supplies}
Grothendieck is isolating the argument which should identify the small stages with the intrinsically accessible objects. The argument does identify them when the stages are closed under the stated colimits. The standalone lemma goes further than its proof permits: presentability gives a retract, and strict fullness does not put that retract back into the subcategory.

\textbf{Corrected implication.}\quad Suppose every object of $\mathcal M$ is a $\pi^+$-filtered colimit of objects of $\mathcal M_0$. Then every object of $\mathcal M_\pi$ is a retract of an object of $\mathcal M_0$. Consequently the printed implication $b\Rightarrow c$ holds if $\mathcal M_0$ is closed under retracts.

To prove this, write $X=\varinjlim_i X_i$ with $X_i\in\mathcal M_0$. If $X$ is $\pi$-accessible, its identity is represented by some $s:X\to X_i$. Writing $r:X_i\to X$ for the colimit map gives $rs=1_X$. Nothing in this argument says that $X$ itself belongs to an arbitrary strictly full subcategory. If $\mathcal M_0$ is closed under the colimits in condition~$a$, the missing conclusion does follow: $r$ is a coequalizer of $sr$ and $1_{X_i}$, so $X\in\mathcal M_0$. This is why the later use of $a\Rightarrow c$ in Proposition~4.5.17 remains sound.

\textbf{Counterexample to the unrestricted lemma.}\ours\quad In $\Ens$, let $\mathcal M_0$ contain exactly the sets of cardinal at most $\pi$ other than singletons. This is a strictly full subcategory of $\mathcal M_\pi$. Every set with at least two elements is the $\pi^+$-filtered union of its subsets of cardinal at most $\pi$ whose cardinal is not one: the union of at most $\pi$ such subsets is another such subset, and each element belongs to a two-element subset. The empty set has its constant presentation.

For the singleton, index by the ordinal $\lambda=\pi^+$, put $X_\alpha=\{0,1\}$, and give every strictly increasing transition the idempotent map $e$ with $e(0)=e(1)=0$. The identity transitions are identities, and $e^2=e$ makes this a functor. Its colimit is a singleton, since every element becomes $0$ at a later stage. The ordinal is $\pi^+$-filtered, but the singleton is absent from $\mathcal M_0$. Thus even the additional assumptions $\mathcal M_0\subset\mathcal M_\pi$ and density of $\mathcal M_\pi$ do not repair the printed equivalence.

The retract principle is standard in the theory of presentable objects; the same argument appears in the higher categorical formulation of \citet[Lemma~5.4.2.4]{Lurie2009HTT}. We include the explicit counterexample to audit this particular formulation; it is not a claim of priority. Our literature comparison concerns the valid retract statement, not an alleged later proof of the erroneous implication.
\end{dcomment}


Reste à prouver que ces conditions, i.e. la stabilité de $\Filt'^{\pi}$ par conoyaux de double-flèches, est satisfaite dans les cas A) B) C), le cas A) étant d'ailleurs déjà connu. Le cas B) résulte immédiatement de la transitivité des familles épimorphiques strictes universelles, et c'est d'ailleurs aussi contenu explicitement (stabilité de $\Filt'^{\pi}$ par $\varinjlim$ de cardinal $\leqslant \pi$) dans SGA 4 I 9.13 (sinon dans l'énoncé, du moins c'est signalé dans la démonstration). En rédigeant SGA 4 I 9, je savais donc qu'on a $\mathcal{M}_{\pi} \subset \Filt'^{\pi}$ pour tout $\pi$, mais l'inclusion, plus évidente pourtant dans le cas d'espèce, $\Filt'^{\pi} \subset \mathcal{M}_{\pi}$ (quand on suppose les objets de $C$ $\pi_1$-accessibles) m'avait apparemment échappé\,?

Reste à traiter le cas C). $\Big($J'ai l'impression qu'il est un peu plus général que le cas B), i.e. que la condition B) sur $\mathcal{M}$ implique C), mais non inversement. Ainsi, la condition C) est vérifiée dès que $\mathcal{M}$ correspond à une espèce de structure algébrique définie en
\msp{page 135}
termes de $\varprojlim$ \textit{finies}, i.e. quand $\mathcal{M}$ est de la forme $\Ind(C)$, $C$ stable par $\varinjlim$ \textit{finies}, alors que je doute que la condition B) soit toujours satisfaite pour ces catégories. Par exemple dans la catégorie Cat (qui est un paratopos de cette espèce), les $\varinjlim$ filtrantes sont exactes à gauche, mais je ne crois pas que les familles épimorphiques stricts soient universelles. % [sic] épimorphiques stricts -> épimorphiques strictes, agreement with "familles"
Prenons p.ex. deux homomorphismes de groupes $G_0 \xrightarrow{u_0} G$, $G_1 \xrightarrow{u_1} G$ tels que $G_0 * G_1 \longrightarrow G$ soit surjectif, alors cette famille à deux éléments est épimorphique stricte dans Gr et dans Cat (\footnote{!, c'est \textit{faux} en général, mais vrai dans le cas particulier plus bas.}), mais n'est pas forcément épimorphique stricte universelle, même si p.ex. $G_0 = G_1 = H$, $G = H \times H$, et $u_0, u_1$ sont des injections $H \rightrightarrows H \times H$, données par $u_0(h) = (h,1)$, et $u_1(h) = (1,h)$. Si je prends l'application diagonale $H \longrightarrow H \times H$, les images inverses de $u_0$, $u_1$ sont le même objet de $\mathcal{M}/H$, savoir $e \longrightarrow H$, et cette famille n'est épimorphique stricte que si $H$ est le groupe unité.$\Big)$

Le cas C) va résulter du résultat plus précis :

\textbf{Corollaire 4.5.15.} \textit{Sous les conditions générales de 4.5.13, supposons que dans $\mathcal{M}$ les} $\varinjlim$
\msp{page 136}
filtrantes \textit{commutent au changement de base (p.ex. les} $\varinjlim$ \textit{filtrantes sont exactes). Alors pour tout $\pi \geqslant \pi_1$,} $\Filt'^{\pi}$ \textit{est formée des $x \in \Ob \mathcal{M}$ satisfaisant la condition suivante.}
% foot: Version 17 décembre 2025

\dgrblock{dgrb1001}{A strict covering family which loses strictness under pullback}
\begin{dcomment}{examples}{p135}{A strict covering family which loses strictness under pullback}
Grothendieck wants an elementary reason that exactness of filtered colimits is weaker than universality of every strict epimorphic family. Groups do provide that reason, but the direct-product example printed here does not: its initial family is already non-strict. The corrective footnote withdraws the general assertion and mistakenly exempts this special case.

Take $H=C_2$, and regard its two copies as the coordinate subgroups of $C_2\times C_2$. Their intersection is trivial. Map the first generator to $(12)\in S_3$ and the second to $(23)\in S_3$. These homomorphisms are compatible on every overlap test: if two maps into the two coordinate subgroups have the same composite into the direct product, both composites are trivial. A common extension would send the two commuting coordinate generators to two noncommuting transpositions. Hence no extension exists. The family is not strictly epimorphic in groups. The same argument applies to one-object group categories, so it is not strictly epimorphic in $\mathrm{Cat}$ either.

Replace the direct product by
\[
 G=C_2*C_2=\langle a,b\mid a^2=b^2=1\rangle.
\]
The two free-factor inclusions are strictly epimorphic in groups. Compatibility on the intersection imposes no relation, and the universal property of the free product gives the unique extension. They are strictly epimorphic as functors of one-object categories as well. Compatibility forces the chosen target objects to coincide; the two homomorphisms into the automorphism group of that object then extend uniquely to $G$.

Now let $K=\langle ab\rangle$. The reduced-word description shows that $ab$ has infinite order, so $K\simeq\mathbb Z$, and that $K$ meets either free factor trivially. Pullback along $K\hookrightarrow G$ gives the family consisting of two copies of $1\to\mathbb Z$. This family is not even epimorphic: the zero endomorphism and the identity of $\mathbb Z$ agree after both maps. Thus strict epimorphic families in groups, and in categories, need not be universal. Filtered colimits in both examples are nevertheless left exact; their objects, arrows or algebraic operations are governed by finite equations in sets.

The repair preserves the intended distinction while changing the ambient group. It also warns against replacing a strict epimorphic family by the assertion that its images generate a target. Generation detects uniqueness of an extension; strictness requires existence for every compatible family.
\end{dcomment}


\printedpage{43}
% head: A. Grothendieck : Dérivateurs, Ch. XVIII, Édition des Universités de Montpellier II et Paris VII
\begin{itemize}
\item[\textit{c$'$)}] \textit{Il existe une famille de flèches $(u_{\alpha})_{\alpha\in A}$, $u_{\alpha} : x_{\alpha} \longrightarrow x$, dans $\mathcal{M}$, qui est} épimorphique stricte universelle \textit{(i.e. couvrante pour la topologie canonique de $\mathcal{M}$), avec $x_{\alpha} \in \Ob C$ $\forall\ \alpha \in A$, et} $\card A \leqslant \pi$.
\end{itemize}
\textsc{Démonstration} : il faut prouver que cette condition équivaut à la condition c) de 4.5.13 1°). Comme on a $\mathrm{c}' \Longrightarrow \mathrm{c}$ trivialement, il reste à montrer que $\mathrm{c} \Longrightarrow \mathrm{c}'$. Mais on a vu dans 4.5.13 1°) que c implique a$''$, i.e. qu'on a un isomorphisme
\[ x \ = \ \varinjlim_{i\in I} x_i \ , \]
les $x_i$ dans $C$, et $I$ ensemble \textit{ordonné filtrant}, avec $\card I \leqslant \pi$. Il reste à voir que si $x$ dans $\mathcal{M}$ est une limite inductive \textit{filtrante} d'objets $x_i$, alors l'ensemble des flèches $x_i \longrightarrow x$ est épimorphique stricte \textit{universelle.} Or cela est immédiat, car pour $x' \longrightarrow x$, on a $x' = \varinjlim\limits_{i} x'_i$ (par hypothèse C), où $x'_i = x' \times_x x_i$, donc la famille $x'_i \longrightarrow x'$ est épimorphique stricte, q.e.d.

Prouvons maintenant que dans le cas C, la condition de 3°, 1 est satisfaite, i.e. que $\Filt'^{\pi}$ est stable par $\varinjlim$ de cardinal $\leqslant \pi$. Cela
\msp{page 137}
résulte aussitôt du

\textbf{Corollaire 4.5.16.} (de 4.5.15). \textit{Sous les conditions de 4.5.15 (i.e. la condition C) de 4.5.13, 3°), soit $x$ un objet de $\mathcal{M}$, et supposons qu'il existe une famille} épimorphique \underline{universelle} $(x_{\alpha} \longrightarrow x)_{\alpha\in A}$ \textit{de flèches de $\mathcal{M}$, avec $x_{\alpha} \in$} $\Ob \Filt'^{\pi}$ \textit{pour tout $\alpha \in A$, et} $\card A \leqslant \pi$. \textit{Alors $x \in$} $\Ob \Filt'^{\pi}$.

En effet, l'hypothèse sur $x_{\alpha}$ équivaut, en vertu de 4.5.15, à l'existence d'une famille épimorphique stricte universelle de flèches $(x_{\alpha\beta} \longrightarrow x_{\alpha})_{\beta\in B_{\alpha}}$, avec $\card B_{\alpha} \leqslant \pi$ et les $x_{\alpha\beta}$ dans $C$. Soit $B = \coprod_{\alpha\in A} B_{\alpha}$, on a alors $\card B \leqslant \pi^2 = \pi$, et d'autre part, la famille des flèches composées $x_{\alpha\beta} \longrightarrow x_{\alpha} \longrightarrow x$ indexée par $B$ est épimorphique stricte universelle (par transitivité de ce type de familles), donc $x \in \Ob \Filt'^{\pi}$, q.e.d.

Avec tout ça, on n'a toujours pas donné de caractérisation de $\mathcal{M}_{\pi}$ lui-même, qui est meilleure que $\Filt^{\pi}$ et $\Filt'^{\pi}$ du point de vue catégorique, car seuls les $\mathcal{M}_{\pi}$ définissent une filtration cardinale, et de plus ils ont les propriétés de stabilité désirables dans ce
\msp{page 138}
contexte, pour les $\varinjlim$ de cardinal $\leqslant \pi$. Le seul critère non tautologique que j'arrive à dégager dans le cas général (pour tout $\pi \geqslant \pi_0$, et sans hypothèse d'exactitude supplémentaire sur $\mathcal{M}$) est le suivant, beaucoup moins ``explicite'' ou ``calculatoire'' que les critères très concrets pour $\Filt^{\pi}$ et pour $\Filt'^{\pi}$ :

\dgrblock{dgrb1002}{What the covering argument actually proves}
\begin{dcomment}{status}{p137}{What the covering argument actually proves}
Grothendieck is using descent along a bounded covering family to keep a presentation bounded. Corollary~4.5.15 makes the necessary descent condition explicit: the families are universally \emph{strictly} epimorphic. In Corollary~4.5.16 the word ``strictement'' disappears from the hypothesis, although the proof still invokes transitivity of precisely that stronger notion.

With the stronger hypothesis restored, the proof is complete. Replace each $X_\alpha$ by a universal strict cover with at most $\pi$ sources in $C$. The composite is again a universal strict cover, and its index set has cardinal at most $\pi\cdot\pi=\pi$. Corollary~4.5.15 applies. The weaker printed hypothesis does not justify that invocation of transitivity. We have not established a counterexample satisfying all the standing hypotheses, in particular $\pi\geq\pi_1$, and therefore record a gap in this proof rather than declaring its conclusion false.

There is a second point in the preceding ``cela r\'esulte aussit\^ot''. Stability under universal strict covers is not, by itself, stability under arbitrary colimits. An arbitrary colimit cocone is strictly epimorphic but need not remain epimorphic after pullback, as the free-product example shows even when filtered colimits are exact. Thus the proposed deduction of the full case~C of Proposition~4.5.13 requires an additional argument. The audit keeps that issue separate from the valid universal-cover calculation and from the general intrinsic closure statement which follows.
\end{dcomment}


\textbf{Proposition 4.5.17.} \textit{Soit $\pi$ un cardinal $\geqslant \pi_0$. Alors $\mathcal{M}_{\pi}$ est égal à la plus petite sous-catégorie strictement pleine de $\mathcal{M}$, contenant $C$, et stable par} $\varinjlim$ \textit{de cardinal $\leqslant \pi$.} (Il est important de noter qu'il \textit{ne suffit pas} de dire : stable par $\varinjlim\limits_{I}$ , avec \textit{$I$ grande devant} $\pi_0$, et $\card I \leqslant \pi$. On trouve alors $\Filt^{\pi}$, qui risque d'être une catégorie strictement plus petite que $\mathcal{M}_{\pi}$. Ce n'est pas non plus la catégorie des $\varinjlim\limits_{I} x_i$, avec $x_i \in C$ $\forall\ i \in \Ob I$, et $\card \Fl I \leqslant \pi$ -- puisque cette catégorie-là est $\Filt'^{\pi}$.)

\textsc{Démonstration}. Soit $\mathcal{M}'$ la catégorie de l'énoncé, on a donc $\mathcal{M}' \subset \mathcal{M}_{\pi}$, puisque $\mathcal{M}_{\pi}$ contient $C$ et est stable par les $\varinjlim$ de cardinal $\leqslant \pi$. Inversement, le lemme 4.5.14, $\mathrm{a} \Longrightarrow \mathrm{c}$.
% foot: Version 17 décembre 2025

% ---- from der-ch18_p046-050.tex
\printedpage{44}
% head: A. Grothendieck : Dérivateurs, Ch. XVIII, Édition des Universités de Montpellier II et Paris VII
% foot: Version 17 décembre 2025

\msp{page 139}

\dgrblock{dgrb1003}{Why closing under relations gives the right category}
\begin{dcomment}{exposition}{p138}{Why closing under relations gives the right category}
Grothendieck now replaces a single bounded presentation by a closure operation. This is a less explicit description, but it has the advantage of being intrinsic once the initial generating category has been fixed. The distinction is between allowing one layer of relations and allowing every further relation needed by the construction.

Let $D$ be the smallest strictly full subcategory containing $C$ and closed under colimits indexed by categories of cardinal at most $\pi$. The represented-functor calculation already gives $D\subset\mathcal M_\pi$. To obtain the reverse inclusion, use the canonical presentation of an object $X$ by $D/X$. Density supplies its colimit, and bounded-colimit closure makes the indexing category $\pi^+$-filtered. If $X$ is $\pi$-accessible, its identity factors through one stage. The resulting retract is a coequalizer of an idempotent and an identity, and therefore still lies in $D$. This is the complete version of the short appeal to Lemma~4.5.14.

For sets, the construction is particularly concrete. Starting with a singleton and closing under at most $\pi$-sized coproducts and coequalizers gives every set of cardinal at most $\pi$, and these are exactly the $\pi^+$-presentable sets. A map from such a set into a $\pi^+$-filtered union is represented at one stage because its at most $\pi$ elements can be treated simultaneously; equality of two maps is treated in the same way. Conversely, a larger set is the $\pi^+$-filtered union of its subsets of cardinal at most $\pi$, and its identity cannot factor through a proper subset. This simple model explains the intended statement without identifying an arbitrary one-step presentation with the completed closure.
\end{dcomment}


\section*{4.6\quad Caractérisation des catégories paratopos (démonstration de 4.3.6)}

Il est temps de revenir à la démonstration du théorème 4.3.6 (p. 70-71), celle que j'avais donnée étant incomplète (et erronée – j'affirmais que $(\mathrm{Filt}^{\pi})_{\pi\geqslant\pi_0}$ est une filtration cardinale, ce qui est faux si on n'admet que $\pi^{\pi_0}=\pi$ pour tout $\pi\geqslant\pi_1$).

On avait vu l'équivalence des conditions a, a$'$, b, soit PTA. Je vais achever la démonstration par le schéma logique
\[
\text{PTA} \ \Longrightarrow\ \text{d)} \ \Longrightarrow\ \text{c)} \ \Longrightarrow\ \text{PTA}\ . \tag*{$(*)$}
\]

Prouvons PTA $\Longrightarrow$ d), i.e. que si $\mathcal{M}$ est un pseudo-topos accessible, alors existe un cardinal infini $\pi_0$, et une petite catégorie $C_0$ stable par $\varinjlim$ de cardinal $\leqslant\pi_0$, et une équivalence de catégories % [sic] alors existe -> alors il existe, missing subject "il"
\[
\mathcal{M} \xrightarrow{\ \sim\ } \Ind'_{\pi_0}(C_0)\ . \tag*{(4.6.1)}
\]

Plus précisément, on a ceci :

\textbf{Proposition 4.6.1.} \textit{Soit $\mathcal{M}$ un paratopos, i.e. une $\mathfrak{U}$-catégorie $\mathcal{M}$ stable par petites $\varinjlim$ , admettant une sous-catégorie pleine génératrice $S$ équivalent à une petite catégorie, formée d'objets accessibles de $\mathcal{M}$. Soit $\pi_0$ un cardinal tel que les objets de $S$ soient $\pi_0$-accessibles (i.e. $S\subset\mathcal{M}_{\pi_0}$), et soit $C$ la plus petite sous-catégorie strictement pleine de $\mathcal{M}$, contenant $S$ et stable par $\varinjlim$ de cardinal $\leqslant\pi_0$. Alors} % [sic] équivalent -> équivalente, agreement with the feminine "sous-catégorie"
\begin{itemize}[leftmargin=2.5em]
\item[\textit{1°)}] \textit{$C$ est équivalente à une petite catégorie.}

\msp{page 140}

\item[\textit{2°)}] \textit{$C=\mathcal{M}_{\pi_0}$, sous-catégorie strictement pleine de $\mathcal{M}$ formée des objets $\pi_0$-accessibles.}
\item[\textit{3°)}] \textit{Le foncteur canonique}
\[
\Ind'_{\pi_0}(C) \xrightarrow[\sim]{\texttt{[}\varphi\texttt{]}} \mathcal{M} % ?? uncertain: font of the brackets around \varphi above the arrow -- typewriter (editorial) brackets or ordinary math brackets
\]
\textit{(défini à isomorphisme unique près par la condition que son composé avec l'inclusion $C\hookrightarrow\Ind'_{\pi_0}(C)$ soit isomorphe à l'inclusion $C\longrightarrow\mathcal{M}$, et que $\varphi$ commute aux petites $\varinjlim$ , ou encore aux $\varinjlim$ filtrantes grandes devant $\pi_0$ – cf. 4.4.16, page 101) est une équivalence de catégories} (\footnote[54]{donc on a aussi $\Ind'_{\pi_0}(C_0)\xrightarrow{\ \sim\ }\mathcal{M}$, où $C_0$ est une \textit{petite} catégorie équivalente à $C$.}).
\end{itemize}

\textsc{Démonstration}. 1°) On va construire $C$ par induction transfinie, en posant
\[
\begin{aligned}
C_0 &\quad=\quad && \text{sous-catégorie strictement pleine de $\mathcal{M}$ engendrée par $S$.}\\
C_{\alpha+1} &\quad=\quad && \text{sous-catégorie strictement pleine de $\mathcal{M}$ engendrée par les $\varinjlim\limits_{I}$ dans}\\
& && \text{$\mathcal{M}$ d'objets de $C_\alpha$, avec $I$ catégorie telle que $\operatorname{card}\operatorname{Fl}I\leqslant\pi_0$.}\\
C_\alpha &\quad=\quad && \textstyle\bigcup_{\alpha'<\alpha}C_{\alpha'}\ \text{ si $\alpha$ est un ordinal limite.}
\end{aligned}
\]

Cela définit de façon unique les $C_\alpha$ (pour tout petit cardinal $\alpha$). Par la définition de $C$, on voit par induction transfinie que les $C_\alpha$ sont contenues dans $C$. Soit $\alpha_0$
\printedpage{45}
% head: A. Grothendieck : Dérivateurs, Ch. XVIII, Édition des Universités de Montpellier II et Paris VII
% foot: Version 17 décembre 2025
\msp{page 141}
% ?? uncertain: a hand-drawn mark (shaped roughly like a "2" or a bracket) stands in the left margin beside lines 3--5 of this page -- marginal annotation reproduced in the image, not typeset text
le premier ordinal tel que $\operatorname{card}\alpha_0>\pi_0$, donc l'ensemble $I_{\alpha_0}$ des cardinaux $\alpha<\alpha_0$ est grand devant $\pi_0$. Je dis qu'on a $C_{\alpha_0}=C$. Pour ceci, il suffit de voir que $C_{\alpha_0}$ est stable par $\varinjlim$ de cardinal $\leqslant\pi$. Mais si $I$ est une catégorie telle que $\operatorname{card}\operatorname{Fl}I\leqslant\pi_0$, et $\varphi : I\longrightarrow C_{\alpha_0}=\bigcup_{\alpha<\alpha_0}C_\alpha$ un foncteur, alors la définition de $\alpha_0$ implique que $\varphi$ se factorise en $\varphi_\alpha : I\longrightarrow C_\alpha$, avec $\alpha<\alpha_0$ (\footnote[55]{Un peu court. Il faut invoquer le fait que $I_{\alpha_0}$ est grande devant $\pi_0$.}). Donc $\varinjlim\varphi\in\operatorname{Ob}C_{\alpha+1}\subset\operatorname{Ob}C_{\alpha_0}$, ce qui montre la stabilité requise. % [sic] de cardinal \leqslant\pi -> de cardinal \leqslant\pi_0, the stability to be proved concerns cardinal \leqslant\pi_0

\dgrblock{dgrb1004}{The small category behind a cocomplete category}
\begin{dcomment}{exposition}{p140}{The small category behind a cocomplete category}
Grothendieck's purpose is to recover a small category which contains all the operations needed to reconstruct $\mathcal M$. The transfinite argument does not require a precise estimate for every Hom set. It requires only that each stage have a set of isomorphism classes, that its Hom sets be small, and that the closure operation stabilize at a specified regular ordinal.

Put $\lambda=\pi_0^+$. Choose a small skeleton at each stage. There is a set of categories with at most $\pi_0$ arrows, a set of functors from each of them to that skeleton, and hence a set of chosen colimits. Local smallness then makes the full category on those chosen objects small. At limit stages take the union of the chosen sets of objects and then the full category they span. All these operations remain in the fixed universe.

If a diagram with at most $\pi_0$ arrows takes values in $C_\lambda$, its objects occur at at most $\pi_0$ earlier stages. Regularity of $\lambda$ bounds those stages by a single $\alpha<\lambda$. Its colimit therefore occurs in $C_{\alpha+1}$. This proves $C_\lambda$ is closed under the required colimits, and proves stabilization. The source's references to ``cardinaux'' in the indexing sequence must here mean ordinals; it is their cofinality which the proof uses.

The remaining reconstruction has a different, categorical, content. For $X_i,Y_j\in C$ and the relevant filtered indices,
\[
 \Hom_{\mathcal M}(\varinjlim_iX_i,\varinjlim_jY_j)
 \simeq\varprojlim_i\varinjlim_j\Hom_C(X_i,Y_j).
\]
This is the Hom formula for the filtered completion. It proves full faithfulness. The category $C/X$ has the bounded colimits inherited from $C$ and is therefore $\lambda$-filtered; density makes its colimit $X$. This proves essential surjectivity. These two arguments establish Proposition~4.6.1 independently of the sharper numerical estimate criticized in the next block.
\end{dcomment}


Pour voir que $C$ est équivalente à une petite catégorie, il suffit de le voir pour les $C_\alpha$, en procédant par induction transfinie. Cela résulte du lemme suivant, immédiat :

\textbf{Lemme 4.6.1.1.}
\begin{itemize}[leftmargin=2.5em]
\item[\textit{1°)}] \textit{La limite inductive $C$ d'un système inductive filtrant $(C_i)_{i\in I}$ de catégories $C_i$ équivalentes à des petites catégories est équivalente à une petite catégorie. Si $\cardred(C_i)$ $\leqslant\pi$ pour tout $i$, et $\operatorname{card}I\leqslant\pi$, alors $\cardred(C)\leqslant\pi$ ($\pi$ étant un} % [sic] système inductive -> système inductif, agreement with the masculine "système"

\msp{page 142}

\textit{cardinal infini). Ici, pour toute catégorie $\mathcal{A}$, je pose $\operatorname{card}\mathcal{A}=\operatorname{card}\operatorname{Fl}\mathcal{A}$, $\cardred(\mathcal{A})$ $=\operatorname{card}(\mathcal{A}_0)$, où $\mathcal{A}_0$ est une catégorie réduite équivalente à $\mathcal{A}$.}
\item[\textit{2°)}] \textit{Soit $C_0\subset\mathcal{M}$ une sous-catégorie pleine d'une catégorie $\mathcal{M}$, $\pi_0$ un cardinal infini, et soit $C'_0$ la sous-catégorie strictement pleine de $\mathcal{M}$ engendrée par les limites inductives dans $\mathcal{M}$ (représentables dans $\mathcal{M}$) de foncteurs $I\longrightarrow C_0$, avec $I$ catégorie telle que $\operatorname{card}I\leqslant\pi_0$. Si $C_0$ est équivalente à une petite catégorie, il en est de même de $C'_0$. Plus précisément, si $\cardred(C_0)\leqslant\pi_1$, $\pi_1$ un cardinal infini, alors $\cardred(C'_0)\leqslant$ $\pi_1^{\pi_0}$} (\footnote[56]{Cette majoration n'est par prouvée et probablement fausse si on ne fait pas des hypothèses un peu plus strictes.}). % [sic] n'est par prouvée -> n'est pas prouvée, misprint in footnote 56
\end{itemize}

\textbf{Corollaire 4.6.1.2.} \textit{$\mathcal{M}$ et $C_0$ étant comme dans 2°) ci-dessus, soit $C$ la plus petite sous-catégorie strictement pleine de $\mathcal{M}$ contenant $C_0$, et stable par $\varinjlim$ de cardinal $\leqslant\pi_0$ (celles qui sont représentables dans $\mathcal{M}$). Si $C_0$ est équivalente à une petite catégorie, il en est de même de $C$. Plus précisément, si $\cardred(C_0)\leqslant\pi_1$, où $\pi_1$ est un cardinal infini, alors}

\msp{page 143}

\textit{on a $\cardred(C)\leqslant\pi_1^{\pi_0}$} (\footnote[57]{majoration non prouvée, probablement fausse, cf. annotation ci-dessus.}).

Le 3°) se prouve par la démonstration transfinie précédente, on prouve par récurrence transfinie que
\[
\cardred C_\alpha \ \leqslant\ \pi_1^{\pi_0}
\]
pour tout $\alpha$. C'est OK pour $\alpha=0$ par hypothèse (puisque $\pi_1\leqslant\pi_1^{\pi_0}$). Si c'est vrai pour $C_\alpha$, alors 2°) assure que
\[
\cardred C_\alpha \ \leqslant\ (\pi_1^{\pi_0})^{\pi_0} \ =\ \pi_1^{\pi_0\cdot\pi_0} \ =\ \pi_1^{\pi_0}\ , % [sic] C_\alpha -> C_{\alpha+1}, this is the successor step of the induction
\]
enfin, si $\alpha$ est un ordinal limite $<\alpha_0$, on aura $C_\alpha=\bigcup_{\alpha'<\alpha}C_{\alpha'}$, donc
\[
\cardred C_\alpha \ \longrightarrow\ \pi_0\cdot\pi_1^{\pi_0} \ =\ \pi_1^{\pi_0} % [sic] \longrightarrow -> \leqslant, an inequality is intended
\]
(vu que $\pi_1^{\pi_0}\geqslant\pi_0$), en vertu du 1° ci-dessus, et compte tenu du fait que $\operatorname{card}I_\alpha\leqslant\pi_0$. Soit $\pi'_0$ le cardinal successeur de $\pi_0$, alors la même majoration nous donne
\[
\cardred C_{\alpha_0} \ \leqslant\ \pi'_0\pi_1^{\pi_0}\ .
\]
\printedpage{46}
% head: A. Grothendieck : Dérivateurs, Ch. XVIII, Édition des Universités de Montpellier II et Paris VII
% foot: Version 17 décembre 2025
Mais on a $\pi_1^{\pi_0}\geqslant2^{\pi_0}>\pi_0$, donc $\pi_1^{\pi_0}\geqslant\pi'_0$, donc $\pi_1^{\pi_0}\geqslant\pi'_0$, donc $\pi'_0\pi_1^{\pi_0}=\pi_1^{\pi_0}$, d'où
\[
\cardred C_{\alpha_0} \ \leqslant\ \pi_1^{\pi_0}\ ,
\]
ce qui est la majoration annoncée.

\msp{page 144}

\dgrblock{dgrb1005}{Counting presentations does not count all new arrows}
\begin{dcomment}{development}{p143}{Counting presentations does not count all new arrows}
Grothendieck has correctly counted the possible small diagrams, but then uses that count as a bound on the morphisms between their colimits. The warning in footnotes~56 and~57 identifies the problem accurately. A colimit can acquire many maps \emph{from} an old object, and its universal property controls maps out of the colimit, not these maps into it.

\textbf{Counterexample to the stated cardinal estimate.}\ours\quad Let $\kappa>2^{\aleph_0}$, and consider algebras with a binary operation $m_\xi$ for every $\xi<\kappa$, subject only to the equations
\[
 m_\xi(x,x)=x \qquad(\xi<\kappa).
\]
This is a finitary equational category. Its free algebra $F(1)$ on one generator has one element: induction on terms reduces every term in that generator to the generator itself. Thus the full subcategory $C_0$ with sole object $F(1)$ has one arrow.

The binary coproduct $F(1)\amalg F(1)$ is $F(2)$. Its elements include the terms $m_\xi(x,y)$, and these are pairwise distinct. To verify the last assertion without assuming a normal-form theorem, fix $\xi\neq\eta$ and interpret the operations on $\{0,1\}$, taking $m_\xi$ to be the first projection and $m_\eta$ the second projection. Assign either projection to every remaining operation. All the idempotence equations hold, but the two specified terms have different values at $(0,1)$. They cannot therefore be equal in the free algebra. It follows that
\[
 |\Hom(F(1),F(2))|=|F(2)|\geq\kappa>2^{\aleph_0}.
\]
Take $\pi_0=\pi_1=\aleph_0$. The hypothesis
$\cardred C_0\leq\pi_1$ holds, but the full closure under even one binary coproduct has more than $\pi_1^{\pi_0}$ arrows in a skeleton. This disproves the quantitative assertions of Lemma~4.6.1.1(2) and Corollary~4.6.1.2 as stated. Their qualitative essential-smallness assertions remain true by the preceding argument.

The construction is an application of the standard free-algebra method, used here to test this particular estimate. We checked the primary representation discussions in \citet{AdamekRosicky1994} and the manuscript's two warnings; we did not locate this exact counterexample there. The marking records our supplied argument, without a priority claim.

A separate edge case affects Corollary~4.6.2(1). In $\Ens$ the one-object full subcategory $S=\{1\}$ is dense, so the printed definition gives $\pi_1=\cardred S=1$. The proposed upper bound is then $1^\pi=1$, although $\mathcal M_\pi$ contains many nonisomorphic objects. One must first replace $\pi_1$ by an infinite bound, such as $\max(\aleph_0,\cardred S)$. Moreover the equality $\pi_1^\pi=2^\pi$ used there requires $2\leq\pi_1\leq\pi$, exactly as the following proof shows. These corrections remove the immediate numerical contradictions; they do not retroactively validate an estimate obtained from the false unrestricted lemma. We retain the small-category reconstruction and distinguish it from the unproved quantitative refinement.
\end{dcomment}


Prouvons 4.6.1 2°, i.e. $C=\mathcal{M}_{\pi_0}$. On a $C\subset\mathcal{M}_{\pi_0}$, puisque par hypothèse $S\subset\mathcal{M}_{\pi_0}$, et on sait que $\mathcal{M}_{\pi_0}$ est stable par $\varinjlim$ de cardinal $\leqslant\pi_0$. D'autre part, on a $\mathcal{M}_{\pi_0}\subset C$, en vertu de 4.5.19 a $\Longrightarrow$ c appliqué à $C$, qui est strictement génératrice dans $\mathcal{M}$, $S$ l'étant, et $C\supset S$, et stable par $\varinjlim$ de cardinal $\leqslant\pi_0$ par construction.

\textit{Prouvons enfin 3°}. Cela équivaut à ceci :

\textbf{Lemme 4.6.1.3.} \textit{Soit $\mathcal{M}$ une $\mathfrak{U}$-catégorie stable par} petites \textit{$\varinjlim$ , $\pi_0$ un cardinal, $\mathcal{M}_0$ une sous-catégorie de $\mathcal{M}$} strictement génératrice, \textit{stable par $\varinjlim$ de cardinal $\leqslant\pi_0$, et formée d'objets $\pi_0$-accessibles. (En vertu de 4.6.1 2° déjà prouvée, appliqué à $S=\mathcal{M}_0$, si $\mathcal{M}_0$ est strictement pleine, on a donc $\mathcal{M}_0=\mathcal{M}_{\pi_0}$.) Alors le foncteur canonique}
\[
\Ind'_{\pi_0}(\mathcal{M}_0) \ \longrightarrow\ \mathcal{M}
\]
\textit{est une équivalence de catégories.}

Le fait que ce facteur est pleinement fidèle résulte aussitôt de l'hypothèse que les objets de $\mathcal{M}_0$ soient $\pi_0$-accessibles : % [sic] facteur -> foncteur, misprint
\[
\begin{aligned}
\operatorname{Hom}_{\mathcal{M}}(\varinjlim_{i\in I}\underbrace{X_i}_{\in\mathcal{M}_0},\ \varinjlim_{j\in J}\underbrace{Y_j}_{\in\mathcal{M}_0}\ ) &\ \simeq\ \varprojlim_{i\in I}\ \underbrace{\operatorname{Hom}_{\mathcal{M}}(X_i,\ \varinjlim_{j\in J}Y_j)}_{\varinjlim\limits_{j\in J}\operatorname{Hom}(X_i,Y_j)}\\
&\ \simeq\ \underline{\operatorname{Hom}}_{\Ind_{\pi_0}(\mathcal{M})}((X_i),(Y_j))\ .
\end{aligned}
\]

\msp{page 145}

Reste à voir que ce foncteur est essentiellement surjectif. Comme $\mathcal{M}_0$ est strictement génératrice, il s'ensuit que tout $x\in\mathcal{M}$ s'écrit
\[
x \ =\ \varinjlim_{i\in I}x_i\ ,
\]
où $I=\mathcal{M}_0/x$, et les $x_i\in\mathcal{M}_0$. L'hypothèse que $\mathcal{M}_0$ est stable par $\varinjlim$ de cardinal $\leqslant\pi_0$ implique la même propriété pour $I=\mathcal{M}_0/x$, donc $I$ est grande devant $\pi_0$. Donc $(x_i)_{i\in I}$ définit un objet de $\Ind_{\pi_0}\mathcal{M}_0$, dont l'image dans $\mathcal{M}$ est isomorphe à $x$, q.e.d. (Comparer SGA 4 I 9.18.)

\textbf{Corollaire 4.6.2.} \textit{Sous les conditions de 4.6.1, prenons pour $\pi_0$} le plus petit \textit{cardinal tel que les objets de $S$ soient $\pi_0$-accessibles. Pour tout $\pi\geqslant\pi_0$, soit $\mathcal{M}_\pi$ la sous-catégorie strictement pleine de $\mathcal{M}$ formée des objets $\pi$-accessibles, qui est aussi (par 4.6.1 2°) l'enveloppe de $C$ dans $\mathcal{M}$ stable par $\varinjlim$ de cardinal $\leqslant\pi$.}
\begin{itemize}[leftmargin=2.5em]
\item[\textit{1°)}] \textit{Soit $\pi_1=\cardred S$. On a alors}
\[
\cardred\mathcal{M}_\pi \ \leqslant\ \pi_1^{\pi} \qquad (=2^{\pi}\ \textit{si}\ \pi\geqslant\pi_0). \tag*{(4.6.2.1)}
\]
\printedpage{47}
% head: A. Grothendieck : Dérivateurs, Ch. XVIII, Édition des Universités de Montpellier II et Paris VII
% foot: Version 17 décembre 2025
\msp{page 146}
\item[\textit{2°)}] \textit{$(\mathcal{M}_\pi)_{\pi\geqslant\pi_1}$ est une filtration cardinale de $\mathcal{M}$, où $\pi_1=\sup(\pi_0,\operatorname{card}S)$.}
\end{itemize}

\textsc{Démonstration}. La majoration (4.6.2.1) est contenue dans 4.6.1.2, où $\pi_0$ est devenu $\pi$. Si $\pi\geqslant\pi_1$, on a
\[
\pi_1^{\pi} \ =\ 2^{\pi} \tag*{$(*)$}
\]
car on a $\pi_1\geqslant2$, d'où $\pi_1^{\pi}\geqslant2^{\pi}$, d'autre part, on a $\pi_1\leqslant2^{\pi_1}$, d'où $\pi_1^{\pi}\leqslant(2^{\pi_1})^{\pi}=2^{\pi_1\pi}=2^{\pi}$, puisque $\pi_1\pi=\pi$ pour $\pi\geqslant\pi_1$ ; d'où $(*)$.

L'assertion 2° est mise pour mémoire ; en vertu de $\mathcal{M}\simeq\Ind'_{\pi_0}\mathcal{M}_{\pi_0}$, on peut supposer que l'on a $\mathcal{M}=\Ind'_{\pi_0}(\mathcal{M}_{\pi_0})$, et on peut alors appliquer 4.5.9. On aurait d'ailleurs envie de dire que $(\mathcal{M}_\pi)_{\pi\geqslant\pi_0}$ est une filtration cardinale, sans se borner aux $\pi\geqslant\pi_1$. C'est vrai sauf pour la majoration de cardinaux dans la condition c) de la définition des filtrations cardinales (SGA 4 I 9.12) :
\[
\underbrace{x}_{\in\operatorname{Ob}\mathcal{M}_{\pi'}} \ =\ \varinjlim_{i\in I}\underbrace{x_i}_{\in\operatorname{Ob}\mathcal{M}_{\pi}}\ , \qquad\qquad \text{$I$ grande devant $\pi$ et $\operatorname{card}I\leqslant\pi'^{\pi}$}
\]
pour $\pi'\geqslant\pi$. Il faut supposer, pour cette majoration, que $\pi\geqslant\pi_1$.

\msp{page 147}

On peut à présent achever la démonstration de 4.3.6, en prouvant les deux dernières implications $(*)$ (p. 139). On a d $\Longrightarrow$ c, i.e. si $\mathcal{M}\simeq\Ind'_{\pi_0}(C)$, $C$ stable par limites inductives de cardinal $\leqslant\pi_0$, et $C$ équivalente à une petite catégorie, alors $\mathcal{M}$ est stable par petites $\varinjlim$ (4.4.13, p. 95-96), ses éléments sont accessibles, et elle admet une filtration cardinale (cf. 4.5.9). Mais il faut vérifier quand même que les éléments de $C$ sont $\pi_0$-accessibles. C'est contenu dans SGA 4 I 9.18, qui affirme que la pleine fidélité du foncteur $\Ind_{\pi_0}(C)\longrightarrow\mathcal{M}$ \textit{équivaut} au fait que les éléments de $C$ sont $\pi_0$-accessibles dans $\mathcal{M}$. La démonstration dans loc. cit. n'était pas claire, le “il suffit” a été vu bas de la page 144. Le “il faut” ne me paraît pas clair, avec la généralité donnée dans loc. cit. – et à présent je vois que c'est même faux, comme on voit en posant pour $C$ la catégorie unité $e$, donc $\Ind'_{\pi_0}C\simeq e$, et l'assertion affirmant que si $e$ est immergée pleinement dans $\mathcal{M}$ (donc par un objet \textit{rigide} de $\mathcal{M}$, i.e. un objet $a$ tel que $\operatorname{End}a=\{\mathrm{id}_a\}$), alors $a$ est $\aleph_0$-accessible. C'est faux dans la catégorie des anneaux, où $\mathbf{R}$ est un objet rigide non $\aleph_0$-accessible !

\msp{page 148}

\dgrblock{dgrb1006}{A rigid object need not be small}
\begin{dcomment}{examples}{p147}{A rigid object need not be small}
Grothendieck's objection to the converse in SGA~4 is decisive. Full faithfulness of a completion functor can be vacuous when the chosen subcategory generates very little of the ambient category. A rigid object makes the point in the smallest possible example.

In the category of unital rings, every endomorphism $f:\mathbb R\to\mathbb R$ is the identity. It is injective because its kernel is an ideal in a field and $f(1)=1$. A positive real number is the square of a nonzero real number, so its image is positive. Thus $f$ preserves the order, and it fixes $\mathbb Q$. If $f(x)\neq x$, a rational number strictly between them contradicts order preservation. Hence $\End(\mathbb R)=\{1\}$.

Nevertheless $\mathbb R$ is not $\aleph_0$-accessible in the manuscript's convention. Let $I$ be the ordered set of its countable unital subrings. The subring generated by a countable union of countable subrings is countable, so $I$ is $\aleph_1$-filtered. Its union is $\mathbb R$, since every individual real belongs to a countable subring. If $\Hom(\mathbb R,-)$ preserved this colimit, the identity would factor through one countable subring $A\subset\mathbb R$. The composite of that factorization with the inclusion would be the identity, forcing $A=\mathbb R$, a contradiction.

The full subcategory on $\mathbb R$ is nevertheless equivalent to the terminal category, and its filtered completion is again terminal. Its inclusion into rings is fully faithful. This proves that the necessity assertion of \citet[Expos\'e~I, Corollary~9.18(a)]{SGA4I1972}, taken at the stated generality, is false. For the intended application $\mathcal M=\Ind_{\pi_0}(C)$ there is no difficulty: evaluation at $c\in C$ computes $\Hom(c,-)$, and the relevant filtered colimits of presheaves are pointwise. It is this calculation, rather than full faithfulness alone, which gives presentability.
\end{dcomment}


Dans le cas qui nous occupe, $\mathcal{M}=\Ind'_{\pi_0}C$, il faut visiblement utiliser le calcul explicite des $\varinjlim$ grandes devant $\pi_0$ dans $\Ind'_{\pi_0}C\subset C^{\wedge}$, elles se calculent en effet “argument par argument”, cf. 4.4.15.

Cela achève la démonstration de c $\Longrightarrow$ d dans $(*)$, page 139. Enfin l'implication c $\Longrightarrow$ PTA est tautologique, l'existence d'une filtration cardinale impliquant celle d'une petite sous-catégorie strictement génératrice. Cela achève la démonstration de 4.3.6.

Je vais redonner une formulation légèrement plus précise de 4.3.6., en précisant les cardinaux, et en y rajoutant une condition supplémentaire, inspirée par SGA 4 I 9.11 (revue et corrigée, cf. une annotation à l'exemplaire \texttt{[de Grothendieck]} de SGA 4 I (\footnote[58]{N.~Éd.\ \ Cette annotation dit “cf. résultat plus simple ci-dessus, si $C$ est \textit{strictement} génératrice” et
% footnote 58 continues at the foot of printed page 48
le résultat en question, également en annotation marginale sur l'exemplaire de Grothendieck, est :\par
\textbf{Corollaire.} \textit{Soit $F$ admettant} \texttt{[une]} \textit{sous-catégorie génératrice par épimorphismes stricts, formée d'objets accessibles, alors} tout \textit{objet de $F$ est accessible (utiliser 7.2. (i) b)).}})).

\printedpage{48}
% head: A. Grothendieck : Dérivateurs, Ch. XVIII, Édition des Universités de Montpellier II et Paris VII
% foot: Version 17 décembre 2025

\textbf{Théorème 4.6.3.} \textit{Soit $\mathcal{M}$ une $\mathfrak{U}$-catégorie.}
\begin{itemize}[leftmargin=2.5em]
\item[\circled{1\textsuperscript{\textit{o}}}] \textit{Les conditions suivantes sont équivalentes pour un cardinal infini $\pi_0$ :}
\begin{itemize}[leftmargin=3em]
\item[$a_{\pi_0}$\textit{)}] \textit{$\mathcal{M}$ est une catégorie $\mathfrak{U}$-$\pi_0$-paratopos, i.e. $\mathcal{M}$ est stable par petites $\varinjlim$ , et contient une petite sous-catégorie pleine} \texttt{[strictement]} \textit{génératrice $S$, formée d'objets $\pi_0$-accessibles de $\mathcal{M}$.}

\msp{page 149}

\item[$b_{\pi_0}$\textit{)}] \textit{Il existe une petite catégorie $S$, et un foncteur pleinement fidèle et $\pi_0$-accessible $\mathcal{M}\xrightarrow{\ i\ }S^{\wedge}$, admettant un adjoint à gauche.}
\item[$c_{\pi_0}$\textit{)}] \textit{Il existe une petite catégorie $C$, stable par $\varinjlim$ de cardinal $\leqslant\pi_0$, telle que $\mathcal{M}\simeq\Ind'_{\pi_0}(C)$.}
\item[$c'_{\pi_0}$\textit{)}] \textit{Comme c), avec $C$ seulement équivalente à une petite catégorie.}
\end{itemize}
\item[\textit{2°)}] \textit{Ces conditions impliquent que $\mathcal{M}$ est stable par petites $\varinjlim$ et petites $\varprojlim$ , et la condition :}
\begin{itemize}[leftmargin=3em]
\item[$e_{\pi_0}$\textit{)}] \textit{Pour toute catégorie $I$ grande devant $\pi_0$, le foncteur}
\[
L \ =\ \varinjlim_{I}\ :\ \underline{\operatorname{Hom}}(I,\mathcal{M}) \ \longrightarrow\ \mathcal{M}
\]
\textit{est exact à gauche (donc exact).}
\end{itemize}
\textit{Cette condition implique notamment ceci :}
\begin{itemize}[leftmargin=3em]
\item[$\alpha(\pi_0)$] \textit{Sur la catégorie des double-flèches de $\mathcal{M}$, le foncteur} Ker \textit{(à valeurs dans $\mathcal{M}$) est $\pi_0$-accessible.}
\item[$\beta(\pi_0)$] \textit{Pour toute flèche $y\leqslant x$ dans $\mathcal{M}$, le foncteur changement de base $\mathcal{M}/x\longrightarrow\mathcal{M}/y$ est $\pi_0$-accessible.}
\end{itemize}
(\footnotemark[59]). % the text of footnote 59 is not printed on pages 46--50 of the PDF

\msp{page 150}

\dgrblock{dgrb1007}{Paratopoi and locally presentable categories}
\begin{dcomment}{later}{p149}{Paratopoi and locally presentable categories}
Grothendieck is now giving two descriptions of the same cocomplete category: a presentation by small objects and a reflection from a presheaf category. In modern language his $\pi_0$-paratopos is a locally $\pi_0^+$-presentable category. The connection with the work of \citet{GabrielUlmer1971} is structural; the publisher's contents place presentable objects, generators, and locally presentable categories together in the central part of that monograph. The equivalences used here are recorded more explicitly as Theorems~1.20 and~1.46 in the publisher's chapter introduction to \citet{AdamekRosicky1994}.

The adjunction explains the equivalence with the presheaf formulation. If
$L:\widehat S\rightleftarrows\mathcal M:N$ is a reflection and $N$ preserves $\pi_0^+$-filtered colimits, then
\[
 \Hom_{\mathcal M}(Lh_s,X)\simeq N(X)(s).
\]
Thus each reflected representable $Lh_s$ is $\pi_0$-accessible. Since a presheaf is the colimit of its representables and $L$ preserves colimits, these objects generate $\mathcal M$. Conversely, restricted Yoneda on a suitable small category of presentables is fully faithful, and its left adjoint sends a presheaf to the corresponding weighted colimit in $\mathcal M$.

There is a missing hypothesis in the later variable-cardinal statement, part~3(c). It says only that $\mathcal M\simeq\Ind'_{\pi_0}(S)$ for small $S$, omitting that $S$ admits colimits of cardinal at most $\pi_0$. The proof on manuscript page~152 restores the hypothesis. Without it, take $S$ to be a discrete category on two objects. Every filtered diagram in $S$ is constant, so its filtered completion is still $S$, which has no coproduct of its two distinct objects. The corrected statement is the cocomplete recognition theorem; allowing arbitrary $S$ leads instead to the accessible categories of the next section.
\end{dcomment}


\item[\textit{3°)}] \textit{Les conditions suivantes sont équivalentes.}
\begin{itemize}[leftmargin=3em]
\item[\textit{a)}] \textit{$\mathcal{M}$ est une catégorie $\mathfrak{U}$-paratopos, i.e. est stable par petites $\varinjlim$ , et contient une petite sous-catégorie pleine $S$ génératrice par épimorphismes strictes, et formée d'objets accessibles.} % [sic] épimorphismes strictes -> épimorphismes stricts, agreement with the masculine "épimorphismes"
\item[\textit{a$'$)}] \textit{Comme a), mais en exigeant que $\mathcal{M}$ soit faiblement accessible, i.e. que} tout \textit{objet de $\mathcal{M}$ soit accessible.}
\item[\textit{b)}] \textit{Il existe une petite catégorie $S$ et un foncteur pleinement fidèle et accessible $\mathcal{M}\hookrightarrow S^{\wedge}$ admettant un adjoint à gauche.}
\item[\circled{c}] \textit{Il existe un cardinal infini $\pi_0$, une petite catégorie $S$, et une équivalence $\mathcal{M}\simeq$ $\Ind'_{\pi_0}(S)$.}
\item[\textit{c$'$)}] \textit{Comme c, en supposant seulement $S$ équivalente à une petite catégorie.}
\item[\textit{d)}] \textit{$\mathcal{M}$ est stable par petites $\varinjlim$ , est faiblement accessible, i.e. tous ses objets sont accessibles, et admet une filtration cardinale.}
\end{itemize}
\end{itemize}

% ---- from der-ch18_p051-055.tex
\printedpage{49}
% head: A. Grothendieck : Dérivateurs, Ch. XVIII, Édition des Universités de Montpellier II et Paris VII
% foot: Version 17 décembre 2025
% The mark of footnote 59 is on the preceding page of the edition; its text is printed at the foot of this page.
\footnotetext[59]{\textbf{Question.} Condition $e_{\pi_0}$ (voire, le couple $\alpha(\pi_0) + \beta(\pi_0)$ plus l'existence d'une petite sous-catégorie strictement génératrice (et la stabilité de $\Mc$ par petites $\varinjlim$ ), implique-t-il que $\Mc$ est un $\pi_0$-paratopos ? Je présume que non.)}

\begin{itemize}[leftmargin=4em]
\item[\circled{e}] (\footnote[60]{C'est là la caractérisation intrinsèque la plus jolie des paratopos par la modalité “économie”, l'autre, l'énoncé structurel, est \circled{c} .}). \textit{$\Mc$ est stable par petites $\varinjlim$ , contient une petite sous-catégorie pleine strictement génératrice, et il existe un cardinal infini $\pi_0$ tel que la condition $\mathrm{e}_{\pi_0}$ de 2°) soit satisfaite.}
\item[\circled{e$'$}] \textit{Comme e, mais en exigeant, au lieu de $\mathrm{e}_{\pi_0}$, la condition plus faible conjonction de $\alpha(\pi_0)$ et $\beta(\pi_0)$ dans 2°.}
\end{itemize}

\msp{page 151}

\textsc{Démonstration}. La partie 1° a déjà été démontrée. La première assertions dans 2°), savoir la stabilité de $\Mc$ par petites $\varprojlim$ , est bien connue (résulte de stabilité par petites $\varinjlim$ , et de l'existence d'une petite famille génératrice, cf. SGA 4 I 8.12.7, version duale de l'annotation marginale). La validité de la condition $\mathrm{e}_{\pi_0}$ résulte du fait qu'elle est valable dans $C^{\wedge}$ (où $C$ est la catégorie mentionnée dans 1°, $\mathrm{c}_{\pi_0}$), et du fait que les $\varprojlim$ dans $\Mc$, et les $\varinjlim$ grandes devant $\pi_0$ dans $\Mc$, se calculent dans $C^{\wedge}$. On trouve d'ailleurs de même : % [sic] assertions -> assertion, singular subject "La première"

\textbf{Corollaire 4.6.4.} \textit{Si $\Mc$ est une catégorie $\U$-$\pi_0$-paratopos, alors pour tout cardinal $\pi \geqslant \pi_0$, et tout catégorie $J$ telle que} $\mathrm{card}\, J \leqslant \pi$\textit{, le foncteur} % [sic] tout catégorie -> toute catégorie, gender agreement
\[
\underline{\mathrm{Hom}}(J,\Mc) \xrightarrow{\ \varprojlim\limits_{J}\ } \Mc
\]
\textit{est $\pi$-accessible, ou ce qui revient au même, pour toute petite catégorie $I$ grande devant $\pi$, le foncteur}
\[
\underline{\mathrm{Hom}}(I,M) \xrightarrow{\ \varinjlim\limits_{I}\ } \Mc % [sic] M -> \mathcal{M}, the first display has \mathcal{M}
\]
\textit{commute aux $\varprojlim$ de cardinal $\leqslant \pi$.}

\msp{page 152}

\dgrblock{dgrb1008}{Exactness comes from the set-valued model}
\begin{dcomment}{exposition}{p151}{Exactness comes from the set-valued model}
Grothendieck wants the cardinal used for generation also to control the passage between limits and colimits. This is available because the small-category presentation gives a faithful calculation of the particular colimits in question, not because an arbitrary colimit-preserving reflection is left exact.

Let $C$ admit colimits indexed by categories of cardinal at most $\pi_0$. The objects of $\Ind_{\pi_0}(C)$ may be identified with the presheaves which take these colimits to limits. One implication follows because representables have this property and $\pi_0^+$-filtered colimits of sets commute with such limits. Conversely, if $F$ has the property, then its category of elements $C/F$ is $\pi_0^+$-filtered: the colimit in $C$ of any bounded diagram of elements receives the unique compatible element supplied by preservation. The empty diagram gives nonemptiness. The canonical Yoneda presentation of $F$ is therefore a presentation of the required filtered kind.

Limits of these presheaves are pointwise, since limits commute with limits. Their $\pi_0^+$-filtered colimits are pointwise by the set-theoretic interchange theorem. Evaluating any proposed comparison map at $c\in C$ consequently reduces it to the corresponding comparison in sets. In particular, for $\pi\geq\pi_0$ and $|\Fl J|\leq\pi$, the map
\[
 \varinjlim_{i\in I}\varprojlim_{j\in J}F_{ij}
 \longrightarrow
 \varprojlim_{j\in J}\varinjlim_{i\in I}F_{ij}
\]
is an isomorphism whenever $I$ is $\pi^+$-filtered. This proves Corollary~4.6.4 and includes the equalizer and base-change assertions in Theorem~4.6.3(2).

The converse recognition argument in part~3 is stated by reference to a corrected use of \citet[Expos\'e~I, Proposition~9.11]{SGA4I1972}. Its conclusion allows an increase of cardinal. The set of subobjects and the set of generating maps into each object enter that argument, so exactness at one prescribed rank does not itself put the generators at that rank. The question in footnote~59 asks precisely whether this increase can be avoided.
\end{dcomment}


D'autre part, le fait que $\mathrm{e}_{\pi_0}$ implique $\alpha(\pi_0)$ et $\beta(\pi_0)$ est immédiat.

La partie 3°), en ce qui concerne l'équivalence des conditions a) à d), est contenue dans 1°, et l'équivalence avec d) a été vue également (la seule chose non tautologique étant que la condition c), savoir $\Mc \simeq \mathrm{Ind}_{\pi_0}(C)$ avec $C$ petite et stable par $\varinjlim$ de cardinal $\leqslant \pi_0$, implique l'existence d'une filtration cardinale, ce qui a été vu dans 4.5.9, et réaffirmé dans 4.6.2 2° ci-dessus). D'autre part, que ces conditions a) à d) impliquent \circled{e} , et que \circled{e} implique e$'$, a été vu dans 2°. Donc il reste à voir que les conditions de e$'$ impliquent a), et pour ceci, il reste à établir que les éléments de $\Mc$ sont accessibles. Mais c'est ce qui est établi dans SGA 4 I 9.11, à cela près que l'énoncé de la condition c), deuxième cas, est canulé, et devrait être remplacé par la condition $\beta(\pi_0)$ (la condition $\alpha(\pi_0)$ figurant dans la condition a) de loc. cit.).

\msp{page 153}

\dgrblock{dgrb1009}{Exactness does not determine the generating cardinal}
\begin{dcomment}{development}{p152}{Exactness does not determine the generating cardinal}
Grothendieck suspects that the fixed-cardinal converse in footnote~59 is false. A topological example gives a negative answer for every infinite cardinal in his convention. It makes the obstruction visible in open covers: colimits may be exact while no nonzero object can be detected at the proposed size.

\textbf{Proposition.}\ours\quad For every infinite $\pi_0$ there is a Grothendieck topos with a small dense subcategory and exact filtered colimits which is not a $\pi_0$-paratopos. In fact its only $\pi_0^+$-presentable object is the initial object.

Put $\lambda=\pi_0^+$, $\kappa=\lambda^+$, and choose a set $X$ of cardinal $\kappa$. Give it the topology whose nonempty opens are the complements of subsets of cardinal less than $\kappa$. This is a topology: finite unions of the excluded small subsets remain small, and an arbitrary union of nonempty opens contains one such open. Every nonempty open has cardinal $\kappa$. Let $\mathcal E=\operatorname{Sh}(X)$.

We first describe an obstruction valid for any noninitial sheaf $F$. Its support $U$, the image of $F\to1$ regarded as an open subset of $X$, is nonempty. Equivalently, $U$ is the set of points where $F$ has a nonempty stalk; it is open because a germ comes from a local section. Choose distinct points $x_\alpha\in U$ for $\alpha<\lambda$, and set
\[
 U_\beta=U\setminus\{x_\alpha:\beta\leq\alpha<\lambda\}
 \qquad(\beta<\lambda).
\]
Each $U_\beta$ is open, since the removed set has cardinal at most $\lambda<\kappa$. The opens increase, their union is $U$, and each is a proper subset of $U$. The ordinal $\lambda$ is a $\lambda$-filtered indexing poset: every set of fewer than $\lambda$ ordinals below it is bounded below it.

Write $F_\beta=F\times U_\beta$, where an open is identified with its subterminal sheaf. Then
\[
 F=\varinjlim_{\beta<\lambda} F_\beta .
\]
This equality can be checked on stalks. At a point in $U$ the point eventually belongs to $U_\beta$, and the stalk then equals $F_x$; outside $U$ every stalk is empty. Each map $F_\beta\to F$ is a monomorphism. If $F$ were $\lambda$-presentable, its identity would factor through one of these monomorphisms. A monomorphism with a right inverse is an isomorphism, forcing the supports $U_\beta$ and $U$ to be equal. This contradicts their construction. Thus no noninitial $F$ is $\lambda$-presentable.

For completeness, $\mathcal E$ has the other properties in the question; the general sheaf-theoretic facts used here are in \citet[Expos\'e~II, Theorem~4.1 and Corollary~4.1.1]{SGA4I1972}. The category of opens is small, and its representable sheaves give the usual dense generating family. All colimits exist by sheafifying presheaf colimits. Filtered colimits are left exact: they are computed by first taking a pointwise filtered colimit of presheaves and then applying the left exact associated-sheaf functor; finite limits of sheaves agree with finite limits of presheaves. Hence $e_{\pi_0}$ holds, as do $\alpha(\pi_0)$ and $\beta(\pi_0)$. Since $X$ is nonempty, the topos is nondegenerate and cannot be generated by its initial object alone. It is therefore not a $\pi_0$-paratopos.

The construction uses the standard sheaf-theoretic relation between support and open covers; it does not contradict the variable-cardinal theorem, since a sheaf topos is locally presentable at a sufficiently large rank. We checked the primary SGA~4 material on sheaves and accessibility and the representation descriptions in \citet{AdamekRosicky1994}; searches concerning compact topologies and presentable sheaves did not locate this exact co-small formulation. The proof is supplied as a development of the manuscript's question, without a claim of priority.
\end{dcomment}


\textbf{Remarque 4.6.5.} Les développements des sections 4.4, 4.5, 4.6 précisent et affirment considérablement, à bien des égards, la théorie des filtrations cardinales de SGA 4 I 9. Il en est ainsi, tout au moins, dans le cas des $\U$-catégories $\Mc$ stables par petites $\varinjlim$ . Dans loc. cit. on travaille avec des $\U$-catégories légèrement plus générales, celles satisfaisant
\printedpage{50}
% head: A. Grothendieck : Dérivateurs, Ch. XVIII, Édition des Universités de Montpellier II et Paris VII
% foot: Version 17 décembre 2025
la \textit{condition} $L$ : il existe un cardinal infini $\pi_0$ tel que $\Mc$ soit stable par petites $\varinjlim$ \textit{grandes devant} $\pi_0$. Je présume cependant qu'on doit pouvoir améliorer de façon similaire la théorie des filtrations cardinales \textit{accessibles}, en s'inspirant du corollaire 9.18 de loc. cit., établissant une équivalence de catégories
\[
\Mc \;\simeq\; \mathrm{Ind}_{\pi_0}(C)\;,
\]
pour un cardinal $\pi_0$ et une petite catégorie $C$ convenable. Ici, on ne peut plus supposer $C$ stable par tel ou tel type de $\varinjlim$ . Néanmoins, 4.4.15 nous dit que sans autre hypothèse sur $C$, $\Mc \overset{\text{déf}}{=} \mathrm{Ind}_{\pi_0}(C)$ satisfait bien à $\mathrm{L}_{\pi_0}$, i.e. est stable (et même stable dans $C^{\wedge}$)

\msp{page 154}

par petites $\varinjlim$ grandes devant $\pi_0$.

Ceci vu, j'ai envie d'appeler \textit{accessible} une catégorie $\Mc$ qui admet une filtration cardinal, et qui est de plus faiblement accessible \texttt{[i.e.]} dont tout objet est accessible. Et de hasarder une % [sic] filtration cardinal -> filtration cardinale, gender agreement

\textbf{Conjecture 4.6.6} (\footnote[61]{Tout ceci sera considérablement précisé dans la théorie qui sera développé dans les sections 4.7, 4.16.}). \textit{Soit $\Mc$ une $\U$-catégorie. Conditions équivalentes :} % [sic] (in footnote 61) développé -> développée, agreement with "la théorie"
\begin{itemize}[leftmargin=4em]
\item[\textit{a)}] \textit{$\Mc$ satisfait la condition $L$, et admet une petite catégorie strictement génératrice formée d'objets accessibles.}
% printed: a vertical double arrow $\Uparrow$ at the left, between a) and b)
\item[\llap{$\Uparrow$\ }\textit{b)}] \textit{Il existe une petite catégorie $C$, et un cardinal infini $\pi_0$, tel que $\Mc$ soit équivalente à la catégorie $\mathrm{Ind}'_{\pi_0}(C)$.}
% printed: a vertical double arrow $\Uparrow$ at the left, between b) and c)
\item[\llap{$\Uparrow$\ }\textit{c)}] \textit{$\Mc$ est “accessible”, i.e. est faiblement accessible (tout objet accessible) et admet une filtration cardinale.}
\end{itemize}
On a b $\Longrightarrow$ a $\Longleftarrow$ c de façon tautologique, et c $\Longrightarrow$ b a été vu dans loc. cit. 9.18, donc on a la suite
\[
\mathrm{c} \;\Longrightarrow\; \mathrm{b} \;\Longrightarrow\; \mathrm{a}\;,
\]
donc il resterait à établir a $\Longrightarrow$ c.

\msp{page 155}

Le travail accompli (et la sueur versée) dans 4.5 nous préservera d'ailleurs, le moment venu, d'erreurs fort naturelles dans la description d'une filtration cardinale, à partir d'une sous-catégorie $S$ strictement génératrice : ce n'est ni $\mathrm{Filt}^{\pi}$, ni $\mathrm{Filt}'^{\pi}$, mais quelque chose de plus subtil -- peut-être faut-il prendre encore la catégorie $\Mc_{\pi}$, ou alors saturer $S$ par les limites inductives de cardinal $\leqslant \pi$ qui veulent bien être représentables dans $\Mc$ : on trouve alors une $\mathrm{Filt}''^{\pi} \subset \Mc_{\pi}$, et si ces $\mathrm{Filt}''^{\pi}$ donnent bien une filtration cardinale, on doit avoir forcément $\Mc_{\pi} \subset \mathrm{Filt}''^{\pi}$, donc finalement $\mathrm{Filt}''^{\pi} = \Mc_{\pi}$. Donc je m'attends à ce qu'en tout état de cause, la filtration cardinale (si elle existe) se fasse par les $\Mc_{\pi}$. La condition a) de loc. cit. 9.12 ($\Mc_{\pi}$ équivalente à une petite catégorie), et la majoration de cardinaux dans loc. cit. c), ne doit pas faire de difficultés, et la condition b) est OK. La difficulté

\msp{page 156}

\dgrblock{dgrb1010}{Removing cocompleteness changes the problem}
\begin{dcomment}{context}{p155}{Removing cocompleteness changes the problem}
Grothendieck is now testing which part of the reconstruction survives when many colimits do not exist. The small objects can still be intrinsically accessible, but the category $\mathcal M_\pi/X$ need not be filtered. A collection of small objects is therefore not yet a collection of small stages of every object.

The modern definition of an accessible category makes the missing presentation part explicit. For a regular $\lambda$, one asks for $\lambda$-filtered colimits and a small full category of $\lambda$-presentable objects from which every object is obtained by such a colimit. Equivalently, the category is a suitable filtered completion of a small category, after allowing retracts among its presentables. This is the structural alternative in Conjecture~4.6.6(b). Its relation to the source's proposed ``cardinal filtration'' requires the generalized filtration developed next; one should not substitute a modern definition into condition~(a) and thereby turn the unproved implication into a tautology.

There is a useful chronological distinction. \citet{GabrielUlmer1971} had treated the cocomplete theory, while \citet{MakkaiPare1989} had already developed accessible categories and their categorical constructions before these manuscript pages. \citet{AdamekRosicky1994} subsequently provided a systematic account of both theories. These references place the present calculations in an established mathematical development; they do not establish that Grothendieck was using those treatments here. The source's immediate interlocutor is its own SGA~4 terminology, whose meaning of ``accessible category'' is reconsidered explicitly in \S4.8.

What has been proved on these pages is enough for the cocomplete case. What remains to be proved for a general completion is control of the indices and their retracts. That is why the next section studies three related filtrations instead of assuming that a single bound on a presentation is already the correct invariant.
\end{dcomment}


provient du fait que $\Mc_{\pi}$ n'étant pas même stable par sommes et conoyaux de double-flèches (quand ceux-ci ne sont pas représentés dans $\Mc$), la catégorie $\Mc_{\pi}/x$ (pour $x$ dans $\Mc$) n'a aucune raison d'être filtrante, pour autant que je vois, ni a fortiori
\printedpage{51}
% head: A. Grothendieck : Dérivateurs, Ch. XVIII, Édition des Universités de Montpellier II et Paris VII
% foot: Version 17 décembre 2025
d'être grande devant $\pi$. Donc il n'est pas trop clair comment dès lors représenter $x$ comme $\varinjlim\limits_{I} x_i$, les $x_i$ dans $\Mc_{\pi}$ et $I$ grande devant $\pi$.

Il faut bien dire que dans SGA 4 I 9, je ne donne aucun \textit{exemple} d'une filtration cardinale sur une catégorie où les petites $\varinjlim$ quelconques ne seraient représentables (cf. loc. cit. 9.13). Peut-être après tout est-ce contre-nature de vouloir elles au delà ? ? Il faut peut-être attendre d'y être forcé, pour essayer de développer une théorie en dehors du cadre des paratopos.

\msp{page 157}

\textbf{4.6.7.} Nous sommes à présent en mesure de reformuler les kyrielles de conditions un peu techniques a) à f) du théorème 2.6 (p. 25 dans XVIII) concernant le caractère “spécial à gauche” de l'ensemble $\mathrm{Mon}(\Mc) \subset \mathrm{Fl}(\Mc)$ :

\textbf{Proposition 4.6.7} (\footnote[62]{Peut-être on pourrait appeler ces catégories des catégories “\textit{paratopos parabéliens}” ? ?}). \textit{Les conditions préliminaires sur la $\U$-catégorie $\Mc$, assurant la validité du théorème XVIII 2.6, sous une forme très légèrement modifié (savoir en modifiant dans 2.6 d) la condition de stabilité de} $\mathrm{Filt}^{\pi}$ \textit{par quotients, en une propriété d'exactitude d) ci-dessous), équivaut à l'ensemble de six conditions suivantes a) à f).} % [sic] modifié -> modifiée, agreement with "forme"; équivaut -> équivalent, plural subject "Les conditions"
\[
\left\{
\begin{array}{ll}
\textit{a)} & \Mc\ \textit{est stable par petites}\ \varinjlim\ . \\[1.5ex]
\textit{b)} & \Mc\ \textit{admet une petite sous-catégorie strictement génératrice.} \\[1.5ex]
\textit{c)} & \textit{Les limites inductives filtrantes dans}\ \Mc\ \textit{sont exactes.}
\end{array}
\right.
\]
\textit{Ces conditions impliquent que $\Mc$ est une} catégorie $\U$-paratopos \textit{(cf. 4.6.3, 3°), et en particulier que les objets de $\Mc$ sont} accessibles, \textit{et que $\Mc$ admet une filtration cardinale $(\Mc_{\pi})_{\pi \geqslant \pi_0}$ telle que pour tout $\pi \geqslant \pi_0$, $\Mc_{\pi}$ soit stable par $\varinjlim$ de cardinal $\leqslant \pi$ (en fait, on aura $\Mc_{\pi} =$ sous-catégorie pleine des objets $\pi$-accessibles de $\Mc$). Aussi, elles impliquent la stabilité de $\Mc$ par petites $\varprojlim$ . D'autre part, les trois conditions suivantes sont des conditions de stabilité pour les monomorphismes.}
\begin{itemize}[leftmargin=4em]
\item[\textit{d)}] \textit{Soit $f : X \longrightarrow Y$ une flèche dans $\Mc$ et $R = X \times_Y X \rightrightarrows X$, alors $X/R \longrightarrow Y$ est un monomorphisme -- en d'autres termes, $f$ se factorise en un composé $ip$, avec $i$ monomorphe et $p$ un épimorphisme effectif.}
\end{itemize}

\msp{page 158}

\begin{itemize}[leftmargin=4em]
\item[\textit{e)}] \textit{Si $A_1$, $A_2$ sont deux sous-objets d'un objet $B$ de $\Mc$, et $A_0 = A_1 \cap A_2 \overset{\text{déf}}{=} A_1 \times_B A_2$, alors $A_1 \amalg_{A_0} A_2 \longrightarrow B$ est un monomorphisme. (Noter la proche parenté des conditions d) et e).)}
\item[\textit{f)}] \textit{L'ensemble} $\mathrm{Mon}(\Mc)$ \textit{des monomorphismes dans $\Mc$ est stable par cochangement de base.}
\end{itemize}
En fait, il n'est pas clair que ces conditions, et notamment d), suffisent à entraîner la stabilité des $\Mc_{\pi}$ par quotients (seulement par quotients effectifs), comme postulé dans 2.6 d). Mais dans l'énoncé et la démonstration donnée dans XV cor. 7.4 (p. 98-99), on se contentait d'une forme un peu plus faible de cette propriété de stabilité, en se bornant aux quotients $Z'$ de $Z \in \mathrm{Ob}\,\Mc_{\pi}$ \textit{tels} qu'il n'existe pas de sous-objet $Z'_0$ de $Z'$, distinct de $Z'$, et tel que $Z \longrightarrow Z'$ se factorise par $Z'_0$ ; on exigeait donc que $Z' \in \mathrm{Ob}\,\Mc_{\pi}$. Mais si la condition d) est satisfaite, il s'ensuit aussitôt qu'un quotient $Z'$ de $Z$ ayant cette propriété supplémentaire est nécessairement effectif, d'où résulte bien que $Z' \in \mathrm{Ob}\,\Mc_{\pi}$.

\printedpage{52}
% head: A. Grothendieck : Dérivateurs, Ch. XVIII, Édition des Universités de Montpellier II et Paris VII
% foot: Version 17 décembre 2025
\msp{page 159}

\dgrblock{dgrb1011}{The three conditions on monomorphisms}
\begin{dcomment}{examples}{p158}{The three conditions on monomorphisms}
Grothendieck returns to the purpose of the size theory: constructing and controlling monomorphisms without losing the boundedness needed by a transfinite argument. The three categorical conditions separate an image factorization, the formation of a union, and a pushout. They agree in familiar exact settings but express different demands.

In an abelian category, condition~(d) is the equality of image and coimage. For subobjects $A_1,A_2\subset B$, the kernel of $A_1\oplus A_2\to B$, $(a_1,a_2)\mapsto a_1+a_2$, is the image of
\[
 A_1\cap A_2\longrightarrow A_1\oplus A_2,
 \qquad t\longmapsto(t,-t).
\]
The pushout over the intersection is the quotient by this kernel, so it embeds in $B$. This proves~(e). If $A\hookrightarrow B$ is monic and $f:A\to C$, its pushout is
$(B\oplus C)/\{(a,-f(a)):a\in A\}$, with the evident categorical cokernel interpretation. The map from $C$ is monic: an element $(0,c)$ in the displayed image forces $a=0$ because $A\to B$ is monic, and then $c=0$. The same statement follows from the corresponding kernel diagram in any abelian category. This proves~(f).

The distinction is visible in groups. For the two coordinate subgroups of $C_2\times C_2$, the pushout over their trivial intersection is $C_2*C_2$; its map to the direct product kills the nontrivial commutator, so~(e) fails. For~(f), take the normal Klein four subgroup $V\subset A_4$ and a surjection $q:V\to C_2$ with kernel generated by one of its nonidentity elements. In the pushout of $V\hookrightarrow A_4$ along $q$, that element is killed. Its conjugates in $A_4$ are the other two nonidentity elements of $V$, so all of $V$ is killed. An element of $V$ mapping to the generator of $C_2$ is therefore killed too. The induced map from $C_2$ is not monic. For commutative unital rings, the pushout of $\mathbb Z\hookrightarrow\mathbb Q$ along $\mathbb Z\to\mathbb F_p$ is the zero ring, so $\mathbb F_p\to0$ gives another failure of~(f).

The last paragraph of the source uses an extremal quotient: a map which cannot factor through a proper subobject. Under~(d), its effective-epi--mono factorization has an invertible mono, hence the quotient is effective. This reduction is correct. One must still keep track of the size of the kernel-pair presentation when asserting preservation of $\pi$-accessibility; finite-colimit closure alone does not bound a kernel pair of an arbitrary quotient.
\end{dcomment}


\textbf{Corollaire 4.6.8.} \textit{Les conditions d), e), f) de 4.6.7 sont satisfaites si $\Mc$ est une catégorie abélienne. Donc l'ensemble des conditions a) à f) est satisfait dans chacun des deux cas suivants :}
\begin{itemize}[leftmargin=4em]
\item[\textit{1)}] \textit{$\Mc$ est une catégorie abélienne où les limites inductives filtrantes sont représentables et exactes, et admettant une petite catégorie strictement génératrice} (\footnote[63]{C'est le cas étudié dans “Tohoku”.}).
\item[\textit{2)}] \textit{$\Mc$ est un topos.}
\end{itemize}
\textbf{Remarque 4.6.9.} Je ne suis par sûr de connaître d'autres exemples où les conditions de 4.6.7 sont vérifiées, que ceux signalés dans 4.6.8 ci-dessus (\footnote[64]{ou qui s'en déduisent trivialement (produits etc.)}). On notera que le cas 1°) qui y est mentionné est celui de Tohoku, où j'établis l'existence de “suffisamment d'objets injectifs” dans une telle catégorie abélienne. Dans Tohoku, je me borne à supposer l'existence d'une petite sous-catégorie, \textit{génératrice pour les monomorphismes}, condition en apparence plus faible que celle énoncée dans 4.6.8 1°). Mais dans une catégorie abélienne, les deux notions (et leurs diverses variantes dans SGA 4 I 7.1) sont équivalentes en vertu de SGA 4 I 7.3 (i). % [sic] par sûr -> pas sûr, negation "ne ... pas"

\msp{page 160}

Par contre, il semble que les conditions e) f) ne sont pratiquement jamais vérifiées pour les catégories définies par les espèces de structure algébrique “définissables par produits finis” sans plus, tel les groupes, les anneaux (de toutes sortes), etc. Mais la catégorie des groupes \textit{commutatifs}, étant abélienne, est une exception. (On peut d'ailleurs l'interpréter comme $\Mc \boxtimes \Mc$, où $\Mc = (\mathrm{Gr})$, et on peut se demander dès lors s'il y a quelque généralisations aux catégories paratopos de cette forme, $\Mc$ étant une catégorie paratopos \ldots) % [sic] tel les groupes -> tels les groupes, agreement; quelque généralisations -> quelques généralisations, plural

\msp{page 161}

\dgrblock{dgrb1012}{Pointed sets and the return to homological algebra}
\begin{dcomment}{examples}{p160}{Pointed sets and the return to homological algebra}
Grothendieck asks for examples beyond abelian categories and topoi. Pointed sets give a small but useful additional example of all six conditions. It shows that the monomorphism properties do not by themselves force additivity or the topos axioms.

The category $\Ens_*$ has all colimits, obtained from disjoint unions and quotient identifications while identifying the specified basepoints. Finite pointed sets form an essentially small dense category, and every pointed set is the filtered union of its finite pointed subsets. Filtered colimits and finite limits are computed on underlying sets with their induced basepoints, so filtered colimits are exact. These verify~(a), (b), and~(c).

For~(d), a pointed map factors through its image, with a surjection onto that image; the surjection is the coequalizer of its kernel pair. For~(e), two pointed subsets of a pointed set have intersection containing the basepoint, and their pushout over that intersection is precisely their union. For~(f), the pushout of an inclusion $A\subset B$ along a pointed map $A\to C$ is the quotient of $B\amalg C$ identifying each $a$ with its image. Such identifications do not identify two distinct elements of $C$: any passage through $B$ can return to $C$ only through the same element of the embedded subset $A$. Thus $C$ still embeds in the pushout.

This category is not abelian: the coproduct of two pointed two-element sets has three elements, whereas their product has four. It is not a nondegenerate topos either: its initial and terminal objects coincide, while a topos with $0\simeq1$ is degenerate. Thus the example answers the concrete request for further settings without suggesting a classification.

The reference to T\^ohoku has a different force. \citet[Theorem~1.10.1]{Grothendieck1957Tohoku} proves the existence of enough injectives under AB5 and a generator. The present conditions seek a categorical environment in which analogous constructions with monomorphisms can be made beyond the additive case. Finally, the familiar identification of group objects in groups with abelian groups follows from interchange: the internal multiplication and the ambient multiplication have the same unit and satisfy
$(ab)*(cd)=(a*c)(b*d)$; substituting units shows that the two multiplications coincide and are commutative. This explains the concrete phenomenon behind the closing suggestion, but supplies no general theorem for the speculative construction $\mathcal M\boxtimes\mathcal M$.
\end{dcomment}


\section*{4.7\quad $\mathrm{Ind}_{\pi}C$ ($C$ quelconque) et filtrations cardinales généralisées.}

Je sens le besoin de me lancer dans des développements à ce sujet, pour parvenir à une meilleure compréhension de la notion de \textit{partie accessible} d'un paratopos, et du théorème de stabilité pour les paratopos par certaines opérations-clef sur ceux-ci.

Dans la suite, $C$ désigne une catégorie équivalente à une petite catégorie, $C_0$ une sous-catégorie réduite exhaustive, $\pi_0$ un cardinal infini. Nous posons
\[
\Mc \;=\; \mathrm{Ind}_{\pi_0}C \;\subset\; C^{\wedge} \qquad \text{(\footnotemark[65])}\;. \tag*{(4.7.1)}
\]
\footnotetext[65]{\textbf{N.B.} Je désigne maintenant par $\mathrm{Ind}C$, $\mathrm{Ind}_{\pi}C$ ce qui précédemment était noté $\mathrm{Ind}'C$, $\mathrm{Ind}'_{\pi}C$.}
On se propose d'introduire sur $\Mc$ diverses filtrations $\pm$ “cardinales”, par des propriétés correspondantes sur un objet $F$ de $\Mc \subset C^{\wedge}$, que je vais à présent introduire, en m'inspirant de 4.5.2 (p. 107) et 4.5.6 (p. 116).

\textbf{Proposition 4.7.1.} \textit{Posons} $\pi_1 = \sup(\pi_0, \mathrm{card}\, C_0)$. \textit{Soit $\pi \geqslant \pi_0$ un cardinal, et considérons les conditions suivantes (dépendant de $\pi$) sur un objet $F$ de $\Mc$ :}
\begin{itemize}[leftmargin=4em]
\item[\textit{a)}] $\mathrm{card}\, C_0/F \leqslant \pi$ \textit{(i.e.} $\mathrm{card}\,\mathrm{Fl}\, C_0/F \leqslant \pi$\textit{).}
\item[\textit{a$'$)}] $\mathrm{card}\,\mathrm{Ob}\, C_0/F \leqslant \pi$.
\item[\textit{a$''$)}] $\mathrm{card}\, F(x) \leqslant \pi$ \textit{pour tout $x \in$} $\mathrm{Ob}\, C$ \textit{(}\textbf{N.B.} \textit{il suffit de l'exiger pour $x$ dans $C_0$).}
\end{itemize}

\printedpage{53}
% head: A. Grothendieck : Dérivateurs, Ch. XVIII, Édition des Universités de Montpellier II et Paris VII
% foot: Version 17 décembre 2025
\msp{page 162}

\dgrblock{dgrb1013}{Presentation size and intrinsic smallness}
\begin{dcomment}{exposition}{p161}{Presentation size and intrinsic smallness}
Grothendieck now removes the existence of small colimits from the hypotheses on $C$. This changes the task. A presheaf may still be assembled from objects of $C$ by sufficiently filtered colimits, but a construction performed on such presentations need not be available inside $C$ itself. He therefore keeps several measures of size in play until he can prove that they agree.

For this section put $\lambda=\pi_0^+$. A category ``grande devant $\pi_0$'' is a $\lambda$-filtered category: every diagram having at most $\pi_0$ arrows admits a cocone. The category $\mathcal M=\Ind_{\pi_0}C$ is thus the ordinary $\lambda$-Ind completion, regarded as a full subcategory of $\widehat C=\Ens^{C^{\mathrm{op}}}$. Its $\lambda$-filtered colimits are computed pointwise. Grothendieck's $\mathcal M_\pi$ consists of the $\pi^+$-presentable objects, those $F$ for which $\Hom(F,-)$ preserves $\pi^+$-filtered colimits. The successor is essential: his diagrams have cardinality $\leq\pi$, whereas modern regular-cardinal notation uses cardinality strictly below the filtering cardinal.

The three tests in Proposition 4.7.1 now have distinct meanings. Conditions (a)--(a$''$) measure a particular presheaf through its elements. Condition (b) bounds a presentation by representables while retaining $\lambda$-filteredness. Condition (c) measures the object intrinsically in $\mathcal M$. The intermediate (c$'$) is intrinsic in the larger category $\widehat C$ and can initially be stronger. Finally $\Filt'^\pi$ is a closure under repeated bounded filtered colimits, whereas $\Filt^\pi$ asks for one presentation. The prime has changed meaning from the preceding cocomplete discussion; here its definition is the one in (b$'$), not a presentation by arbitrary colimits in an ambient cocomplete category.

Two easy implications deserve to be made explicit. Representables are $\pi^+$-presentable in $\widehat C$, and a colimit of at most $\pi$ such presentable objects is again $\pi^+$-presentable, because $\pi^+$-filtered colimits of sets commute with limits over categories of at most $\pi$ arrows. This gives $\Filt'^\pi\subseteq(\widehat C)_\pi$. Since $\mathcal M\hookrightarrow\widehat C$ preserves the filtered colimits in question, ambient presentability implies presentability in $\mathcal M$. Neither argument bounds the number of objects of $C$.

The cardinal $\pi_1$ enters when one replaces the total number of elements by a pointwise bound. If $|\Fl C_0|\leq\pi$ and every $|F(c)|\leq\pi$, then $C_0/F$ has at most $\pi$ arrows. Without the bound on $C_0$, this summation is unavailable. The printed unconditional arrow (a$''$)$\Rightarrow$(c$'$) must therefore be treated separately; the example below shows an actual failure, rather than merely an incomplete proof.
\end{dcomment}


\begin{itemize}[leftmargin=4em]
\item[\textit{b)}] \textit{On a $F = \varinjlim\limits_{i \in I} x_i$, avec les $x_i \in$} $\mathrm{Ob}\, C$\textit{, $I$ une catégorie (pas nécessairement ordonnée) grande devant $\pi_0$, et telle que} $\mathrm{card}\, I \leqslant \pi$. \textit{(}\textbf{N.B.} \textit{en vertu de 4.4.2 (p.79), pour $I$ donnée, ceci équivaut à l'existence d'un foncteur cofinal $I \longrightarrow C/F$, ou encore $I \longrightarrow C_0/F$.)} (\footnote[66]{\textbf{N.B.} On désignera par $\mathrm{Filt}^{\pi} \subset \Mc$ la sous-catégorie strictement pleine de $\Mc$ formée des $F$ satisfaisant cette condition.}).
\item[\textit{b$'$)}] \textit{On a $F \in$} $\mathrm{Ob}\,\mathrm{Filt}'^{\pi}$\textit{, où} $\mathrm{Filt}'^{\pi}$ \textit{est la plus petite sous-catégorie strictement pleine de $\Mc$ contenant $C$, et stable par $\varinjlim$ grandes devant $\pi_0$ et de cardinal $\leqslant \pi$.}
\item[\textit{c)}] \textit{On a $F \in$} $\mathrm{Ob}\,\Mc_{\pi}$\textit{, i.e. $F$ est un objet $\pi$-accessible de $\Mc$.}
\item[\textit{c$'$)}] \textit{On a $F \in$} $\mathrm{Ob}\, C^{\wedge}_{\pi}$\textit{, i.e. $F$ est un objet $\pi$-accessible de $C^{\wedge}$.}
\end{itemize}
\textit{Ceci posé, on a le diagramme d'implications suivant.}
\[
\begin{tikzcd}[column sep=large, row sep=large]
\mathrm{a} \arrow[r, Rightarrow] \arrow[dr, Rightarrow] & \mathrm{a}' \arrow[r, Rightarrow] & \mathrm{a}'' \arrow[dr, Rightarrow] & & \\
& \mathrm{b} \arrow[r, Rightarrow] & \mathrm{b}' \arrow[r, Rightarrow] & \mathrm{c}' \arrow[dr, Rightarrow] \arrow[ulll, dotted, bend right=50, "\text{si } \pi \geqslant \pi_1"'] & \\
& & & & \mathrm{c}\;, \arrow[ull, dotted, bend left=30, "\text{si } \pi = \pi^{\pi_0}"']
\end{tikzcd}
\]
\textit{où les deux implications en pointillés, c$'$ $\Longrightarrow$ a et c $\Longrightarrow$ b, sont conditionnelles soumises à la condition $\pi \geqslant \pi_1$, resp. $\pi = \pi^{\pi_0}$.}

\msp{page 163}

\textbf{Corollaire 4.7.2.}
\begin{itemize}[leftmargin=4em]
\item[\textit{1°)}] \textit{Si $\pi \geqslant \pi_1$, les six conditions a, a$'$, a$''$, b, b$'$, c$'$ sont équivalentes, et la condition en question (qu'on peut écrire $F \in$} $\mathrm{Ob}\,\mathrm{Filt}^{\pi}$ \textit{où $F \in$} $\mathrm{Ob}\,\mathrm{Filt}'^{\pi}$ \textit{indifféremment) implique c :} % [sic] où -> ou, alternative "ou ... indifféremment"
\[
\mathrm{Filt}^{\pi} \;=\; \mathrm{Filt}'^{\pi} \;\subset\; \Mc_{\pi} \qquad \textit{si } \pi \geqslant \pi_1 \tag*{(4.7.2.1)}
\]
\textit{(alors qu'en général, pour $\pi \geqslant \pi_0$, on a}
\[
\mathrm{Filt}^{\pi} \;\subset\; \mathrm{Filt}'^{\pi}\ (\,\subset C^{\wedge}_{\pi} \cap \Mc\,) \;\subset\; \Mc_{\pi} \qquad \textit{si } \pi \geqslant \pi_0\textit{).} \tag*{(4.7.2.2)}
\]
\item[\textit{2°)}] \textit{Quand $\pi = \pi^{\pi_0}$, p.ex. quand $\pi$ est de la forme $2^{c}$, avec $c \geqslant \pi_0$, alors les conditions b, b$'$, c, c$'$ sont équivalentes, i.e. on a}
\[
\mathrm{Filt}^{\pi} \;=\; \mathrm{Filt}'^{\pi} \;=\; (C^{\wedge})_{\pi} \cap \Mc \;=\; \Mc_{\pi}\;. \tag*{(4.7.3.2)} % [sic] (4.7.3.2) -> (4.7.2.3), the neighbouring labels are 4.7.2.1, 4.7.2.2, 4.7.2.4
\]
\item[\textit{3°)}] \textit{Quand $\pi \geqslant \pi_1$} et \textit{$\pi = \pi^{\pi_0}$, alors les sept conditions a), a$'$), a$''$), b) , b$'$), c) , c$'$) sont équivalentes, donc en particulier}
\[
\mathrm{Filt}^{\pi} \;=\; \mathrm{Filt}'^{\pi} \;=\; C^{\wedge}_{\pi} \cap \Mc \;=\; \Mc_{\pi} \qquad (\pi \geqslant \pi_1 \textit{ et } \pi = \pi^{\pi_0}\textit{).} \tag*{(4.7.2.4)}
\]
\end{itemize}

% ---- from der-ch18_p056-060.tex
\printedpage{54}
% head: A. Grothendieck : Dérivateurs, Ch. XVIII, Édition des Universités de Montpellier II et Paris VII

\textsc{Démonstration} de 4.7.1. Les implications a $\Longrightarrow$ a$'$ $\Longrightarrow$ a$''$ sont tautologiques, ainsi que b $\Longrightarrow$ b$'$. L’implication b$'$ $\Longrightarrow$ c$'$ provient de la stabilité de $C^{\wedge}_{\pi}$ par $\varinjlim$ de cardinal $\leqslant \pi$, et du fait que les $\varinjlim$ du type envisagé (en particulier, grandes devant $\pi_0$) \textit{dans} $\mathcal{M}$ sont les mêmes que dans $C^{\wedge}$, i.e. se calculent “argument par argument” (cf. 4.4.15, p. 99). Ceci, plus le fait que l’inclusion $\mathcal{M} \hookrightarrow C^{\wedge}$ est pleinement fidèle, implique aussi c$'$ $\Longrightarrow$ c.

\msp{page 164}
L’implication a$''$ $\Longrightarrow$ c$'$ se voit comme dans 4.5.6 (p. 119) en appliquant le lemme 4.5.8 1° et 2°, au cas de $C^{\wedge} = \underline{\mathrm{Hom}}(C^{\circ}, \mathcal{E})$, où $\mathcal{E} =$ (Ens). L’implication conditionnelle inverse c$'$ $\Longrightarrow$ a$''$, pour $\pi \geqslant \pi_1$, se voit comme dans loc. cit., en appliquant ce lemme 4.5.8. Il reste donc à établir encore les trois implications
\[
\left\{
\begin{array}{ccll}
\text{a} & \Longrightarrow & \text{b} & \\[4pt]
\text{a}'' & \Longrightarrow & \text{a} & \text{si } \pi \geqslant \pi_1 \\[4pt]
\text{c} & \Longrightarrow & \text{b} & \text{si } \pi = \pi^{\pi_0}.
\end{array}
\right. \tag*{$(*)$}
\]
Mais a $\Longrightarrow$ b est tautologique, en prenant $I = C_0/F$, puisqu’on sait que $C_0/F$ est grande devant $\pi_0$ et que $F = \varinjlim\limits_{i \in I} x_i$, pratiquement par définition de $\mathcal{M} = \underbrace{\Ind_{\pi}(C)}_{\simeq \Ind_{\pi} C_0} \subset C^{\wedge}$. Donc il reste à établir les deux implications conditionnelles de $(*)$. L’implication a$''$ $\Longrightarrow$ a pour $\pi \geqslant \pi_1$ s’établit comme dans 4.5.2 (cf. page 109), c’est essentiellement sorital. Reste à établir
\[
\text{c} \;\Longrightarrow\; \text{b} \qquad \text{si} \;\; \pi \;=\; \pi^{\pi_0} \;. \tag*{$(*)$}
\]
Nous aurons besoin du résultat auxiliaire :

\msp{page 165}
\textbf{Lemme 4.7.3.} \textit{Soit $\pi \geqslant \pi_0$ tel que $\pi^{\pi_0} = \pi$ (donc en fait $\pi > \pi_0$). Pour tout objet $F$ de $\mathcal{M}$, on a}
\[
F \;=\; \varinjlim_{i \in I} x_i \qquad \text{(\footnotemark[67])}\;. \tag*{(4.7.3.1)}
\]
\footnotetext[67]{les $x_i$ dans $\Filt^{\pi}$, $I$ ensemble ordonné grand devant $\pi$.}
\textit{Si de plus $\pi' > \pi$ et $F \in \Ob \Filt^{\pi'}$, alors on peut prendre $I$ telle que}
\[
\card I \;\leqslant\; \pi'^{\pi} \;. \tag*{(4.7.3.2)}
\]

\dgrblock{dgrb1014}{A pointwise singleton with a large presentation}
\begin{dcomment}{examples}{p164}{A pointwise singleton with a large presentation}
Grothendieck wants the presheaf description to give a manageable test for accessibility. A thin category separates the pointwise size from the amount of data needed to present the presheaf, and makes the missing hypothesis in (a$''$)$\Rightarrow$(c$'$) visible without any algebraic complications.

Fix an infinite $\pi\geq\pi_0$, and let $\kappa>\pi$ be a regular cardinal. Take $C$ to be the ordinal $\kappa$, with one arrow $\alpha\to\beta$ when $\alpha\leq\beta$. This is a small category in the working universe, and it is $\pi^+$-filtered: a family of at most $\pi$ ordinals below $\kappa$ has a supremum below $\kappa$. Write $h_\beta=C(-,\beta)$. Its value at $\alpha$ is a singleton if $\alpha\leq\beta$ and empty otherwise. Consequently
\[
 F=\varinjlim_{\beta<\kappa}h_\beta
\]
is the terminal presheaf, with $F(\alpha)=\{*\}$ at every object. In particular $F\in\Ind_{\pi_0}C$, and condition (a$''$) holds even with the bound $1$.

Yet there is no natural transformation $F\to h_\beta$: evaluation at $\beta+1$ would require a map from a singleton to the empty set. Thus
\[
 \varinjlim_{\beta<\kappa}\Hom(F,h_\beta)=\varnothing,
 \qquad
 \Hom\left(F,\varinjlim_{\beta<\kappa}h_\beta\right)=\{1_F\}.
\]
The same presentation witnesses failure of $\pi^+$-presentability both in $\widehat C$ and in $\Ind_{\pi_0}C$. Conditions (c$'$) and (c) fail. All arrows and colimits used here are ordinary categorical ones; there is no issue of a choice of model for the Ind completion.

The repair to Proposition 4.7.1 is to impose $\pi\geq\pi_1$ on the arrow (a$''$)$\Rightarrow$(c$'$), or to use condition (a), which bounds the whole category of elements. Under that size hypothesis, the six-way equivalence of Corollary 4.7.2(1) is unaffected. One should not extend that equivalence to small $\pi$ merely because all the values of a presheaf happen to be small.
\end{dcomment}


En effet, on a
\[
F \;=\; \varinjlim_{j \in J} x_j \qquad \text{avec } J = C_0/F \qquad \text{(\footnotemark[68])}\;,
\]
\footnotetext[68]{où $J$ peut aussi se remplacer par une catégorie cofinale dans [?] $C_0/F$ grande devant $\pi_0$.}
avec $J$ grande devant $\pi_0$. Soit alors
\[
I \;\;=\;\; \parbox[t]{0.62\textwidth}{ensemble des sous-catégories $J'$ de $J$ \textit{grandes devant} $\pi_0$, et telles que $\card J' \leqslant \pi$.} \tag*{(4.7.3.3)}
\]

\printedpage{55}
% head: A. Grothendieck : Dérivateurs, Ch. XVIII, Édition des Universités de Montpellier II et Paris VII

Je dis que $I$ (ordonné par inclusion) \textit{est grand devant} $\pi$ (\footnote[69]{cf. démonstration 4.7.5, page 168-171.}). Cela résulte aussitôt du lemme 4.5.5, p. 113, (énoncé d’abord, de façon erronée, en supposant seulement $\pi > \pi_0$ et non $\pi^{\pi_0} = \pi$). Posons alors, pour tout $J' \in I$,
\[
x_{J'} \;=\; \varinjlim_{j \in J'} x_j \;,
\]
on a donc
\[
x_{J'} \;\in\; \Ob \Filt^{\pi}
\]
par définition de $I$ et de $\Filt^{\pi}$. D’autre part,
\[
J' \;\longmapsto\; x_{J'} \;:\; I \;\longrightarrow\; \Filt^{\pi}
\]
\msp{page 166}
est un foncteur, donc
\[
\varinjlim_{J' \in I} x_{J'} \;\simeq\; \varinjlim_{j \in J} x_j \;\simeq\; x \;, % [sic] x -> F, the colimit is the object F of (4.7.3.1)
\]
d’où la première assertion de 4.7.3. Reste la majoration de cardinaux pour $F \in \Ob \Filt^{\pi}$. Dans ce cas, quitte à remplacer $C_0/F$ par une catégorie cofinale grande devant $\pi_0$, $\card J \leqslant \pi'$. On voit alors tout de suite que $\card I \leqslant \pi'^{\pi_0}$, ce qui est la majoration annoncée.

Nous pouvons maintenant prouver l’implication $(*)$, page 164, i.e.
\[
\mathcal{M}_{\pi} \;\subset\; \Filt^{\pi} \qquad (\text{pour } \pi = \pi^{\pi_0}).
\]
Soit donc $F$ dans $\mathcal{M}_{\pi}$, i.e. $\pi$-accessible, écrivons $F$ sous la forme (4.7.3.2), avec $I$ grande devant $\pi$. Par définition de la $\pi$-accessibilité, appliqué à $F \xrightarrow{\mathrm{id}} F \simeq \varinjlim F_i$, on voit que $F$ est un facteur direct d’un $F_i$, donc d’un objet de $\Filt^{\pi}$. Faute de savoir que $\Filt^{\pi}$ est stable par facteurs directs, on ne peut en conclure que $F$ est dans $\Filt^{\pi}$ (aussi je ne peux prouver c $\Longrightarrow$ b), mais seulement que $F$ est dans $\Filt'^{\pi}$, qui, lui, est stable par $\varinjlim$ grandes devant $\pi_0$, de cardinal $\pi$ donné $> \pi_0$, donc aussi \msp{page 167} stable par facteurs directs. Cela achève la démonstration de 4.7.1, et en même temps on obtient :

\dgrblock{dgrb1015}{Removing the retract without increasing the bound}
\begin{dcomment}{exposition}{p166}{Removing the retract without increasing the bound}
Grothendieck's factorization of $1_F$ is the decisive accessibility argument, but it initially produces a retract of one stage. He explicitly stops short of the claimed implication (c)$\Rightarrow$(b). The missing step has a positive answer. It belongs to the retract-removal result of Makkai--Par\'e, Proposition 2.3.11; an accessible verification and a stronger closure statement are given by \citet[Corollaries 4.13 and 4.15, including footnote 2]{Chen2021Sifted}. We give the elementary construction needed here, keeping track of the number of arrows.

\textbf{Retract lemma.} Let $\pi\geq\pi_0$ be infinite. If a presheaf $H$ is a retract of
\[
 F=\varinjlim_{i\in I}h_{c_i},
 \qquad I\text{ is }\pi_0^+\text{-filtered},\qquad |\Fl I|\leq\pi,
\]
then $H$ has a presentation of the same form with an indexing category having at most $\pi$ arrows. No closure property of $C$ is needed.

\textit{Proof.} Choose $s:H\to F$ and $r:F\to H$ with $rs=1_H$, and put $e=sr$. Let $q_i:h_{c_i}\to F$ be the presentation maps. By Yoneda and pointwise computation of the colimit, for each $i$ there are an object $t(i)$ of $I$ and a morphism $v_i:c_i\to c_{t(i)}$ satisfying
\[
 q_{t(i)}h_{v_i}=e q_i.
\]
Put $\alpha_i=rq_i$. Define $J$ to have the same \emph{labelled} objects as $I$. A morphism from label $i$ to label $j$ is an actual morphism $u:c_i\to c_j$ satisfying $\alpha_jh_u=\alpha_i$. Inside this category retain only the subcategory generated by the images of all arrows of $I$ and by the arrows $v_i:i\to t(i)$. The latter are morphisms over $H$, since $re=r$. Finite compositions of an infinite set of at most $\pi$ generators give at most $\pi$ arrows, so $|\Fl J|\leq\pi$.

We check filteredness with all compatibility conditions included. Let $D:K\to J$ have at most $\pi_0$ arrows. For each vertex $a$ of $K$, use the chosen map $v_{D(a)}$; choose an $I$-cocone on their targets and call its vertex $k$. This produces maps
\[
 w_a:c_{D(a)}\longrightarrow c_k,
 \qquad q_kh_{w_a}=s\alpha_{D(a)},
\]
each a composite of a chosen $v_i$ with an $I$-arrow. For an arrow $f:a\to b$ of $K$, the two maps $w_bD(f)$ and $w_a$ have equal composites into $F$, because the $\alpha_i$ are a cocone over $H$. Equality in the filtered colimit
\[
 F(c_{D(a)})=\varinjlim_{i\in I}C(c_{D(a)},c_i)
\]
is witnessed after an $I$-arrow out of $k$. There are at most $\pi_0$ such equalities. The $\pi_0^+$-filteredness of $I$ gives a single further arrow $k\to l$ which makes all of them hold simultaneously: take a cocone on the diagram consisting of $k$, the chosen witnesses, and the arrows from $k$ to them. The resulting maps $c_{D(a)}\to c_l$ are arrows of $J$ and form the required cocone. This also handles the empty diagram, since $I$ is nonempty. Thus $J$ is $\pi_0^+$-filtered.

Finally the maps $\alpha_i$ induce $\varinjlim_Jh_{c_i}\to H$. At an object $c$ of $C$, every element of $H(c)$ lifts along $s$ to $F(c)$ and therefore comes from some $C(c,c_i)$; the map is surjective. Suppose two such representatives $u:c\to c_i$ and $u':c\to c_j$ have the same image in $H(c)$. After composition with $v_i$ and $v_j$, their images in $F(c)$ are equal, since both are the image under $s$ of that common element. Filteredness supplies a common $I$-stage where the two representatives agree. The arrows to that stage, and $v_i,v_j$, all belong to $J$, so the original representatives agree in the colimit over $J$. The map is injective at every $c$, hence an isomorphism.\hfill$\square$

Apply the lemma to the retract produced on manuscript p.~166. At $\pi^{\pi_0}=\pi$ it proves the originally intended equality
\[
 \Filt^\pi=\Filt'^\pi=(\widehat C)_\pi\cap\mathcal M
 =\mathcal M_\pi.
\]
Thus the statement of Corollary 4.7.2(2) survives, while the proof actually written on these pages requires this additional argument. The construction uses the author's explicit permission that an indexing category need not be ordered; replacing it by a cofinal poset without a cardinal estimate would lose the point of the lemma.
\end{dcomment}


% marginal vertical bar beside Corollaire 4.7.4 and its statement
\textbf{Corollaire 4.7.4.} \textit{Si $\pi \geqslant \pi_0$ est tel que $\pi^{\pi_0} = \pi$, alors on a}
\[
\Filt'^{\pi} \;=\; (\Filt^{\pi})^{\natural} \;, \tag*{(4.7.4.1)}
\]
\textit{i.e.} $\Filt'^{\pi}$ \textit{est formée des objets $F$ de $\mathcal{M}$ qui sont facteurs directs d’un objet de} $\Filt^{\pi}$.

En effet, comme $\Filt^{\pi} \subset \Filt'^{\pi}$, et que $(\Filt^{\pi})^{\natural} \subset \Filt'^{\pi}$, et que $\Filt'^{\pi}$ est stable par facteurs directs, on a $(\Filt^{\pi})^{\natural} \subset \Filt'^{\pi}$ pour tout $\pi \geqslant \pi_0$. Mais inversement, comme $\Filt'^{\pi} \subset \mathcal{M}_{\pi}$, l’argument précédent montre que pour $\pi = \pi^{\pi_0}$, on a $\Filt'^{\pi} \subset \mathcal{M}_{\pi} \subset (\Filt^{\pi})^{\natural}$, ce qui prouve 4.7.3.

\textbf{Remarque 4.7.4.1} \texttt{[deux fois 4.7.4.1]}. Bien sûr, par 4.7.2, 1°, dès que $\pi \geqslant \pi_1$, on a même $\Filt^{\pi} = \Filt'^{\pi}$ (sans avoir à supposer $\pi = \pi^{\pi_0}$), donc \textit{si} $\pi > \pi_0$ (ce qui sera la cas si $\pi_1 > \pi_0$, i.e. $\pi_1 \neq \pi_0$). $\Filt'^{\pi}$ étant stable par facteurs directs, $\Filt^{\pi}$ l’est aussi, (4.7.4.1) est satisfait trivialement. Mais si $\pi = \pi_1 = \pi_0$, on a $\Filt^{\pi} = \Filt^{\pi_0} = \Filt'^{\pi_0} =$ sous-catégorie des foncteurs \textit{représentables} dans $C$ (équivalent à $C$), et celle-ci n’est égale à $C^{\natural}$ dans $\mathcal{M}$ que si $C$ est karoubienne, cf. \ldots{} ci-dessous. % [sic] la cas -> le cas, gender of "cas"

\printedpage{56}
% head: A. Grothendieck : Dérivateurs, Ch. XVIII, Édition des Universités de Montpellier II et Paris VII

\msp{page 168}
\textbf{4.7.5.} Pour que la démonstration de 4.7.1 soit complète, il me faudrait encore donner une démonstration de 4.5.5, p. 113, qui est le lemme technique clef pour établir le

\dgrblock{dgrb1016}{The initial level already contains idempotent splittings}
\begin{dcomment}{examples}{p167}{The initial level already contains idempotent splittings}
Grothendieck's concern about retracts also reaches the endpoint $\pi=\pi_0$. Here the literal assertion that $\Filt^{\pi_0}$ consists only of representables is false under the definition just given. The obstruction is present in the smallest possible nontrivial idempotent category.

Let $E$ have one object $*$ and endomorphism monoid $\{1,e\}$ with $e^2=e$ and $e\neq1$. Every diagram $K\to E$, of any size, admits a cocone whose components are $e$: for either possible arrow $a$ of $E$, one has $ea=e$. Hence $E$ is $\lambda$-filtered for every regular $\lambda$. It has no terminal object, since its unique object has two endomorphisms.

Take $C=E$. The representable $h_*$ has value $\{1,e\}$. The natural endomorphism induced by $e$ sends both elements to $e$, and its image is the singleton presheaf $P$. The colimit of the diagram $E\to\widehat E$ that sends the object to $h_*$ and $e$ to this endomorphism is $P$. It is a $\pi_0^+$-filtered presentation with just two indexing arrows, so $P\in\Filt^{\pi_0}$. It is not representable, since its value has one element rather than two.

This also clarifies the phrase ``cofinale \`a $e$'' in footnote 76. The unique functor $E\to e$ is final, as it is for any nonempty filtered category. This neither makes $E$ equivalent to $e$ nor supplies a final functor $e\to E$. In particular it does not say that the colimit of every diagram on $E$ is one of its values. The relevant conclusion of the smallness assumption is a retract, as the next argument shows.

The corrected general statement is
\[
 \Filt^{\pi_0}=\Filt'^{\pi_0}
 =\mathcal M_{\pi_0}\simeq\Kar(C).
\]
Indeed, if $I$ is $\pi_0^+$-filtered and $|\Fl I|\leq\pi_0$, the identity diagram of $I$ has a cocone $u_i:i\to i_*$. For any diagram $D:I\to\widehat C$, these arrows induce a map from its colimit $L$ to $D(i_*)$. Its composite with $D(i_*)\to L$ is $1_L$, by the colimit universal property. Thus a bounded presentation at the initial level gives a retract of a representable. Conversely every retract of a representable is the colimit of the two-arrow idempotent category just used. The same argument shows that this class is closed under further colimits of the prescribed size, proving the equality with $\Filt'^{\pi_0}$.

For the equality with $\mathcal M_{\pi_0}$, write a $\pi_0^+$-presentable object as its defining $\pi_0^+$-filtered colimit of representables and factor its identity through one stage. Conversely representables are presentable and retracts preserve presentability. This is the ordinary categorical form of the idempotent-completion statement in \citet[Lemma 5.4.2.4]{Lurie2009HTT}. In particular the conclusion does not require $C$ itself to be Karoubian. What Karoubian completeness provides is the further identification $\Kar(C)\simeq C$.
\end{dcomment}


\textbf{Lemme 4.7.5.1} (utilisé dans la démonstration de 4.7.3). \textit{Soient $\pi_0$, $\pi$ deux cardinaux infinis, avec $\pi \geqslant \pi_0$ et $\pi^{\pi_0} = \pi$ (donc $\pi > \pi_0$). Soit $J$ une catégorie grande devant $\pi_0$. Alors l’ensemble $I$, ordonné par inclusion, des sous-catégories $J'$ de $J$ qui sont grandes devant $\pi_0$, et telles que} $\card J' \leqslant \pi$, \textit{est grand devant $\pi$. (De plus, si} $\card J \leqslant \pi'$, \textit{on a} $\card I \leqslant \pi'^{\pi}$\textit{.)}

Il faut donc montrer le

\textbf{Corollaire 4.7.5.2} (\footnote[70]{Comparer avec la démonstration de 4.6.1.2.}) (c’est le lemme 4.5.5 non encore démontré). \textit{Sous les conditions de 4.7.5, soit $K \subset \mathrm{Fl} J$ tel que} $\card K \leqslant \pi$. \textit{Alors il existe $J' \in I$ tel que $K \subset J'$.}

On sait, $\pi$ étant infini, que la catégorie $K'$ engendrée par $K$ est aussi de cardinal $\leqslant \pi$, donc on peut supposer que $K$ “est” une sous-catégorie de $J$. On va construire une $J' \in I$ contenant $K$, en construisant par récurrence transfinie une suite croissante $K_{\alpha}$ de sous-catégories de $J$, avec $\card K_{\alpha} \leqslant \pi$, indexée par les ordinaux $\alpha$ tels que \fbox{$\card \alpha \leqslant \pi$} :

\msp{page 169}
\begin{enumerate}
\item[1°)] $K_0 = K$.
\item[2°)] $K_{\alpha} = \bigcup_{\beta < \alpha} K_{\beta}$ ($= \varinjlim K_{\beta}$) si $\alpha$ est un ordinal limite (donc $\card K_{\alpha} \leqslant \pi$, puisque les $\card K_{\beta} \leqslant \pi$).
\item[3°)] $K_{\alpha+1}$ déduit de $K_{\alpha}$ de la façon suivante. Si $K_{\alpha}$ est grande devant $\pi_0$, on prend $K_{\alpha+1} = K_{\alpha}$ (donc, compte tenu de 2°, on aura $K_{\alpha'} = K_{\alpha}$ pour $\alpha' \geqslant \alpha$ -- la construction transfinie s’arrête \ldots). Sinon, $K_{\alpha+1}$ s’obtient ainsi :
\begin{enumerate}
\item[a)] pour toute partie $A$ de cardinal $\leqslant \pi_0$ de $\Ob K_{\alpha}$, on choisit un majorant $x_A$ dans $J$ (on a envie de le prendre dans $K_{\alpha}$, s’il y a un majorant dans $K_{\alpha}$)\,;
\item[b)] pour tout couple $x$, $y$ d’objets de $K_{\alpha}$, et toute partie de cardinal $\leqslant \pi_0$ de $\mathrm{Hom}_{K_{\alpha}}(x,y)$, on choisit une flèche $u_B : y \longrightarrow z_B$ dans $J$, qui égalise les $v_{\beta} \in B$ (et on prend cette flèche dans $K_{\alpha}$, quand c’est possible)\,;
\item[c)] on prend pour $K_{\alpha+1}$ la plus petite sous-catégorie contenant les $\mathrm{id}_{x_A}$ et les $u_B$.
\end{enumerate}
\end{enumerate}

Il faut vérifier que
\[
\card K_{\alpha+1} \;\leqslant\; \pi \;.
\]
Le nombre des choix possibles pour $A \in \mathfrak{P}_{\leqslant \pi_0}(\Ob J)$ \msp{page 170} est $\leqslant \pi^{\pi_0} = \pi$, le nombre des choix possibles pour $B \in \mathfrak{P}_{\pi_0}(\mathrm{Fl}\, J)$ est aussi $\leqslant \pi^{\pi_0} = \pi$, donc le nombre des flèches de la forme $\mathrm{id}_{x_A}$ ou $u_B$ est $\leqslant \pi + \pi = \pi$. Mais alors la sous-catégorie $K_{\alpha+1}$ de $J$ engendrée par cet ensemble de flèches est de cardinal $\leqslant 1$. Soit alors % [sic] \leqslant 1 -> \leqslant \pi, the statement to be verified is card K_{\alpha+1} \leqslant \pi
\printedpage{57}
% head: A. Grothendieck : Dérivateurs, Ch. XVIII, Édition des Universités de Montpellier II et Paris VII
$\alpha_0$ le plus petit ordinal tel que
\[
\card \alpha_0 \;>\; \pi_0 \;,
\]
on a donc
\[
\card \alpha_0 \;=\; \text{cardinal successeur de } \pi_0 \leqslant \pi \;,
\]
puisque $\pi^{\pi_0} = \pi$, d’où $\pi > \pi_0$ (car $\pi^{\pi_0} = \pi_0^{\pi_0} \geqslant 2^{\pi_0} > \pi_0$). Ainsi, $\alpha_0$ fait partie des ordinaux $\alpha$ admissibles (i.e. tels que $\card \alpha \leqslant \pi$). Je dis que $K_{\alpha_0}$ \textit{est grande devant} $\pi_0$ (ce qui prouvera 4.7.5.2, en prenant $J' = K_{\alpha_0}$). Cela résulte immédiatement du fait que
\[
K_{\alpha_0} \;=\; \bigcup_{\alpha < \alpha_0} K_{\alpha} \qquad (\text{réunion \textit{croissante}}),
\]
% double marginal vertical bar beside the following two lines
et que l’ensemble des $\alpha < \alpha_0$ est \textit{grande devant} $\pi_0$ (\footnote[71]{expliciter ceci en un

\textbf{Lemme.} \textit{Soit $\pi_0$ un cardinal, $\alpha_0 = \alpha_{\pi'_0}$ ($\pi'_0$ le successeur de $\pi_0$) le plus petit ordinal tel que} $\card \alpha_0 > \pi_0$, \textit{alors $I_{\alpha_0} = \{\alpha \mid \alpha < \alpha_0\}$ est grand devant $\pi_0$.}}). Donc tout ensemble d’objets ou de flèches de $K_{\alpha_0}$, de cardinal $\leqslant \pi_0$, est contenu dans un $K_{\alpha}$ avec $\alpha < \alpha_0$.

\msp{page 171}
Cela implique, s’il s’agit d’un ensemble d’objets, qu’il y a un majorant commun dans $K_{\alpha+1}$, a fortiori dans $K_{\alpha_0}$. S’il s’agit d’une partie d’un $\mathrm{Hom}_{K_{\alpha}}(x,y)$, qu’il y a un égalisateur $y \longrightarrow z$ dans $K_{\alpha+1}$, a fortiori dans $K_{\alpha_0}$. Donc le critère $\mathrm{PF}_{\pi}$ 1,2 de 4.4.4 (page 89) est satisfait par $K_{\alpha_0}$, ce qui achève la démonstration de 4.7.5.2.

\dgrblock{dgrb1017}{Closing a small diagram under filteredness witnesses}
\begin{dcomment}{exposition}{p170}{Closing a small diagram under filteredness witnesses}
Grothendieck is building a categorical version of a downward size argument: a prescribed small part of a filtered indexing category should lie inside a small filtered subcategory. The fixed-point assumption $\pi^{\pi_0}=\pi$ is exactly what pays for adjoining all witnesses of size at most $\pi_0$ at each step. The proof is sound after two pieces of construction data, implicit in the prose but absent from clause 3(c), are restored.

Let $J$ be $\pi_0^+$-filtered and let $K$ be a set of at most $\pi$ arrows. First include their endpoints and identities and close under composition. This gives a subcategory $K_0$ with at most $\pi$ arrows. At a successor stage retain \emph{all} arrows of $K_\alpha$. For each set $A\subseteq\Ob K_\alpha$ of cardinality at most $\pi_0$, choose an object $z_A$ and arrows $a\to z_A$ for all $a\in A$, and adjoin the arrows as well as their target. Include the empty set here, so that nonemptiness is ensured even when $K_0$ is empty. For each pair of objects $x,y$ and each family $B\subseteq K_\alpha(x,y)$ of cardinality at most $\pi_0$, adjoin an arrow $y\to z_B$ which coequalizes $B$. Then close under identities and finite composition. At a limit stage take the union.

There are at most $\pi^{\pi_0}=\pi$ choices of a small family of existing arrows or objects. Each contributes at most $\pi_0$ new arrows. Closing under composition still leaves at most $\pi$ arrows. Thus all $K_\alpha$ with $|\alpha|\leq\pi$ have size at most $\pi$. The printed reference to subsets of $\Ob J$ and $\Fl J$ in the counting argument must refer to the current $K_\alpha$; there is no bound on $J$ itself. Likewise the printed final bound $1$ is $\pi$.

Set $\delta=\pi_0^+$, regarded as its initial ordinal. It is regular, and $\delta\leq\pi$ because $\pi^{\pi_0}=\pi$ forces $\pi>\pi_0$. A family of at most $\pi_0$ arrows of $K_\delta$ is contained in some earlier $K_\alpha$: take the supremum of their appearance stages, which is below $\delta$. The next stage supplies the required common target or common coequalizer. These two properties imply the full cocone criterion. For a diagram of at most $\pi_0$ arrows, first send all vertices to one target; for each diagram arrow compare the two induced maps to that target. Coequalize each pair, send the resulting targets to one common target, and then coequalize the resulting family of parallel arrows out of the first common target. This produces one map which coequalizes every pair. Thus $K_\delta$ is $\pi_0^+$-filtered.

This proves the extension lemma. It also proves that the poset $\mathscr I$ of all such subcategories is $\pi^+$-directed: a family of at most $\pi$ of them has union of arrow sets of cardinality at most $\pi$, and the extension construction supplies an upper bound. Every arrow of $J$ lies in an object of $\mathscr I$. A compatible family of cocones on the restrictions to those subcategories is therefore exactly a cocone on all of $J$. Applying the universal property twice gives
\[
 \varinjlim_{J'\in\mathscr I}\ \varinjlim_{j\in J'}D(j)
 \;\simeq\;\varinjlim_{j\in J}D(j).
\]
If $|\Fl J|\leq\pi'$, then $|\mathscr I|\leq(\pi')^\pi$, since every subcategory is determined by at most $\pi$ arrows. This is the announced bound in Lemma 4.7.3; the exponent $\pi_0$ printed in its proof on p.~166 is not the bound established by this counting.

The same construction is the familiar mechanism for raising the index of accessibility; compare \citet[Theorem 3.4]{Vasey2020Accessible}. In the higher categorical setting its poset form is used in \citet[Lemma 5.4.2.8 and Proposition 5.4.2.9]{Lurie2009HTT}. In each form the essential work is to retain enough witnesses while keeping their total size under control.
\end{dcomment}


\textbf{Corollaire 4.7.6} (de 4.7.1). \textit{Pour tout cardinal $\pi \geqslant \pi_0$, on a}
\[
\underbrace{\Filt'^{\pi}}_{\overset{\textit{si } \pi \geqslant \pi_1}{=}\ \Filt^{\pi}} \;\subset\; \mathcal{M}_{\pi} \;\subset\; \underbrace{\Filt^{\pi^{\pi_0}}}_{\overset{\textit{si } \pi^{\pi_0} \geqslant \pi_1}{=}\ \Filt^{\pi^{\pi_0}}} \qquad (= (\Filt^{\pi^{\pi_0}})^{\natural} \textit{ par cor. 4.7.4})\,, \tag*{(4.7.6.1)} % ?? uncertain: tiny condition under the second brace -- "si \pi^{\pi_0} \geqslant \pi_1" or "si \pi^{\pi}_0 \geqslant \pi_1"
\]
\textit{donc si $\pi \geqslant \pi_1$, on a aussi}
\[
\Filt^{\pi} \;\subset\; \mathcal{M}_{\pi} \;\subset\; \Filt^{\pi^{\pi_0}} \qquad \textit{(si $\pi^{\pi_0} \geqslant \pi_1$).} \tag*{(4.7.6.2)}
\]

En effet, $\Filt'^{\pi} \subset \mathcal{M}_{\pi}$ n’est autre que l’implication b$'$ $\Longrightarrow$ c de 4.7.1, d’autre part, on a $\pi \leqslant \pi^{\pi_0}$, d’où
\[
\mathcal{M}_{\pi} \;\subset\; \mathcal{M}_{\pi^{\pi_0}} \;=\; \Filt'^{\pi^{\pi_0}} \;,
\]
où la dernière égalité provient de $\mathcal{M}_{\pi'} = \Filt'^{\pi'}$ si $\pi'^{\pi_0} = \pi'$, i.e. de (4.7.2.3) (page 163), en faisant $\pi' = \pi^{\pi_0}$, compte tenu que $\pi'^{\pi_0} = (\pi^{\pi_0})^{\pi_0} = \pi^{\pi_0 \cdot \pi_0} = \pi^{\pi_0} = \pi$.

\msp{page 172}
\textbf{Proposition 4.7.7.} \textit{Considérons $C$, $C_0$, $\mathcal{M} = \Ind_{\pi_0}(C)$, $\pi_1 = \sup(\pi_0, \card \mathrm{Fl}\, C_0)$ comme précédemment. Pour tout cardinal $\pi \geqslant \pi_0$, soit $\mathcal{F}^{\pi}$ la sous-catégorie strictement pleine} $\Filt'^{\pi}$ \textit{(resp. $(C^{\wedge})_{\pi} \cap \mathcal{M}$, resp. $\mathcal{M}_{\pi}$) dans $\mathcal{M}$. Alors la filtration croissante $(\mathcal{F}^{\pi})_{\pi \geqslant \pi_0}$ de $\mathcal{M}$ par les sous-catégories $\mathcal{F}^{\pi}$ satisfait les conditions suivantes.}
\begin{enumerate}
\item[\textit{a)}] \textit{Pour tout $\pi \geqslant \pi_0$, $\mathcal{F}^{\pi}$ est équivalente à une petite catégorie.}
\item[\textit{b)}] \textit{$\mathcal{M}$ satisfait à $\mathrm{L}_{\pi_0}$ (i.e. est stable par $\varinjlim$ filtrantes grandes devant $\pi_0$), et pour tout $\pi \geqslant \pi_0$, $\mathcal{F}^{\pi}$ est stable par les limites filtrantes grandes devant $\pi_0$ et de cardinal $\leqslant \pi_0$ [?].}
\printedpage{58}
% head: A. Grothendieck : Dérivateurs, Ch. XVIII, Édition des Universités de Montpellier II et Paris VII
\item[\textit{c)}] \textit{Si $\pi \geqslant \pi_0$ est} tel que \fbox{$\pi^{\pi_0} = \pi$,} \textit{alors pour tout $F$ dans $\mathcal{M}$, on a un isomorphisme}
\[
F \;=\; \varinjlim_{i \in I} x_i \;,
\]
\textit{les $x_i$ dans $\mathcal{F}^{\pi}$, et $I$ catégorie ordonnée grande devant $\pi_0$. Si de plus $\pi' > \pi$ et $F \in \Ob F^{\pi'}$, on peut prendre $I$ telle que} % [sic] F^{\pi'} (italic F) -> \mathcal{F}^{\pi'}, the filtration is denoted \mathcal{F}^{\pi}
\[
\card I \;\leqslant\; \pi'^{\pi} \qquad \underline{\textit{si}}\ \pi'^{\pi} \geqslant \pi_1. \tag*{(4.7.7.1)}
\]
\item[\textit{d)}] \textit{Si $\pi \geqslant \pi_0$, on a $\mathcal{F}^{\pi} \subset \mathcal{M}_{\pi}$, i.e. les objets de $\mathcal{F}^{\pi}$ sont $\pi$-accessibles.}
\end{enumerate}
\textbf{N.B.} Ce sont là exactement, à part d), les conditions a), b), c) des filtrations cardinales de SGA 4 I 9.7.12, \msp{page 173} à cela près seulement que dans c), on se restreint ici aux ordinaux $\pi$ tels que $\pi^{\pi_0} = \pi$ (\footnote[72]{et \textit{de plus}, dans la majoration (4.7.7.1), aux cardinaux $\pi' > \pi$ tels que l’on ait de plus $\pi'^{\pi} \geqslant \pi_1$ (ce qui est automatique si on a $\pi \geqslant \pi_1$, ce qui est le cas, vu que $\pi = \pi^{\pi_0} \geqslant 2^{\pi_0}$, si on a $2^{\pi_0} \geqslant \pi_1$.)}). La condition supplémentaire d) est une “condition d’accessibilité”, cf. plus bas.

\textsc{Démonstration de} 4.7.7. a) Comme on a
\[
\Filt'^{\pi} \;\subset\; (C^{\wedge})_{\pi} \cap \mathcal{M} \;\subset\; \mathcal{M}_{\pi} \;,
\]
on est réduit à prouver que $\mathcal{M}_{\pi}$ est équivalente à une petite catégorie. On est ramené au cas $\pi \geqslant \pi_1$, $\pi = \pi^{\pi_0}$, auquel cas on a $\mathcal{M}_{\pi} = \Filt^{\pi}$. En fait, on trouve :

\textbf{Corollaire 4.7.7.1} \texttt{[deux fois 4.7.7.1]} (\footnote[73]{Comparer avec 4.7.19.}). \textit{Si $\pi \geqslant \pi_1$, on a}
\[
\left\{
\begin{array}{rcll}
\cardred \Filt^{\pi} & \leqslant & 2^{\pi} & \\[4pt]
\cardred \Filt'^{\pi} & \leqslant & \cardred \mathcal{M}_{\pi} \;\leqslant\; 2^{(\pi^{\pi_0})} & \qquad \textit{($= 2^{\pi}$ si $\pi^{\pi_0} = \pi$).}
\end{array}
\right.
\]

La deuxième majoration résulte de la première et de
\[
\mathcal{F}^{\pi} \;\subset\; \mathcal{M}_{\pi} \;\subset\; \Filt^{\pi^{\pi_0}} \qquad (\text{pour } \pi^{\pi_0} \geqslant \pi_1) \qquad \text{(\footnotemark[74]).}
\]
\footnotetext[74]{Cf. (4.7.6.2).}
Pour la première, on utilise la condition a$''$) de 4.7.1 (qui caractérise $\Filt^{\pi}$ pour $\pi \geqslant \pi_1$). Je passe sur les détails de la vérification.

\msp{page 174}
Prouvons 4.7.2 b). La condition $\mathrm{L}_{\pi_0}$ est connue, reste la stabilité de $\mathcal{F}^{\pi}$. Dans le cas de $\Filt'^{\pi}$, elle est acquise par définition. Dans les cas $C^{\wedge}_{\pi} \cap \mathcal{M}$ et $\mathcal{M}_{\pi}$, elle résulte de la stabilité de la notion de $\pi$-accessibilité par $\varinjlim$ de cardinal $\leqslant \pi$, soit dans $C^{\wedge}$, soit dans $\mathcal{M}$ (compte tenu que $C^{\wedge}$ et $\mathcal{M}$ ont une famille génératrice $C$ formée d’objets $\pi_0$-accessibles, donc $\pi$-accessibles).

Prouvons c) : c’est essentiellement le lemme 4.7.3, compte tenu que l’on a
\[
\Filt^{\pi} \;\subset\; \mathcal{F}^{\pi} \;,
\]

% ---- from der-ch18_p061-065.tex
\noindent\printedpage{59}% head: A. Grothendieck : Dérivateurs, Ch. XVIII, Édition des Universités de Montpellier II et Paris VII
% foot: Version 17 décembre 2025
sauf qu'il faut faire attention à la majoration de cardinaux, $\card I \leqslant \pi'^{\pi}$. Mais si $F \in \Ob\mathcal{F}^{\pi'}$, comme par (4.7.6.1)
\[
\mathcal{F}^{\pi'} \;\subset\; \underbrace{\mathcal{M}_{\pi'}}_{\subset\mathcal{M}_{\pi'^{\pi}}\subset\Filt'^{\pi'^{\pi}}} \;\subset\; \Filt'^{\pi'^{\pi_0}} ,
\]
donc a fortiori
\[
\mathcal{F}^{\pi'} \;\subset\; \Filt^{\pi'^{\pi}} \underbrace{\;=\;}_{\text{si } \pi'^{\pi}\geqslant\pi_1} \Filt^{\pi'^{\pi}} , % [sic] first \Filt^{\pi'^{\pi}} -> \Filt'^{\pi'^{\pi}}, otherwise the equality under the brace would be trivial (cf. the preceding display)
\]
d'où, par 4.7.3, avec $\pi'$ remplacé par $\pi'^{\pi}=\pi''$, la possibilité de prendre $I$ avec
\[
\card I \;\leqslant\; \pi''^{\pi} \;=\; (\pi'^{\pi})^{\pi} \;=\; \pi'^{\pi\cdot\pi} \;=\; \pi'^{\pi} ,
\]
ce qui achève la démonstration.

\noindent\msp{page 175}\textbf{Remarque 4.7.7.2.} J'ai été obligé, dans l'énoncé de la condition c), de faire intervenir un deuxième cardinal $\pi_1$, et ceci faut de savoir si, pour $\pi''$ tel que $\pi''=\pi''^{\pi_0}$, $\Filt^{\pi''}$ est stable par facteurs directs. (Il faudrait, pour bien faire, soit prouver qu'il est stable, donc qu'on a dans 4.7.1 c $\Longrightarrow$ b, i.e. $\mathcal{M}_\pi\subset\Filt^{\pi}$ (et non seulement c $\Longrightarrow$ b$'$, i.e. $\mathcal{M}_\pi\subset\Filt'^{\pi}$) si $\pi=\pi^{\pi_0}$, ou bien faire un contre-exemple!) Cet ingrédient un peu étrange, on sera hélas obligé de le traîner dans la définition des filtrations cardinales généralisées, si on ne veut s'interdire de regarder dans le cas présent $(\mathcal{F}^{\pi})_{\pi\geqslant\pi_0}$ comme une telle filtration, et ne s'autoriser que de regarder $(\mathcal{F}^{\pi})_{\pi\geqslant\pi_1}$ comme telle, ce qui me paraîtrait ridicule. Pratiquement (même si théoriquement l'introduction de $\pi_1$ dans la description d'une filtration cardinale généralisée était bel et bien nécessaire) on n'aura guère des cas où il faille faire mention de $\pi_1$ dans la condition c. En effet, $\pi\geqslant\pi_0$ et $\pi^{\pi_0}=\pi$ impliquent $\pi\geqslant 2^{\pi_0}$, et dès lors $\pi'^{\pi}\geqslant 2^{(2^{\pi_0})}$. Si donc $\pi_1\leqslant 2^{(2^{\pi_0})}$, i.e. $\card C_0\leqslant 2^{(2^{\pi_0})}$, alors la condition $\pi^{\pi'}\geqslant\pi_1$ dans c) est automatiquement satisfaite, et il n'y a donc pas à la mentionner. Or il est bien rare qu'on ait à faire à des petites catégories $C_0$, dont le cardinal soit $>$ puissance % [sic] faut -> faute, "faute de savoir" intended; [sic] \pi^{\pi'}\geqslant\pi_1 -> \pi'^{\pi}\geqslant\pi_1, the condition is written \pi'^{\pi}\geqslant\pi_1 elsewhere (cf. 4.7.8)
\msp{page 176}
du continu $c$, et a fortiori, telle que $\card C_0 > 2^{c}$. Si par contre $\card C_0\leqslant 2^{c}$, comme $2^{c}\leqslant 2^{(2^{\pi_0})}$ (puisque $\pi_0\geqslant\aleph_0$), on aura bien $\pi_1\leqslant 2^{(2^{\pi_0})}$, et on est tranquille.

\dgrblock{dgrb1018}{The bounded closure question has a positive answer}
\begin{dcomment}{later}{p175}{The bounded closure question has a positive answer}
Grothendieck introduces the second threshold $\pi_1$ because he has not yet removed retracts or shown that a bounded iteration can be replaced by one bounded presentation. The general closure result now available answers precisely this difficulty. For infinite $\pi>\pi_0$, set $\lambda=\pi_0^+$ and $\mu=\pi^+$. Then \citet[Corollary 4.13]{Chen2021Sifted} identifies the closure of representables under $\mu$-small $\lambda$-filtered colimits with the presheaves admitting a single such presentation. In the notation of these pages this says
\[
 \Filt^\pi=\Filt'^\pi\qquad(\pi>\pi_0).
\]
The endpoint follows from the idempotent calculation above. The fixed-point presentability statement is the Makkai--Par\'e result recovered in Chen's Corollary 4.15. This answers the question on pp.~175--176 affirmatively; it is not an open problem of the current theory.

The distinction between the two arguments remains useful. Retract closure alone repairs the inference in the proof of 4.7.1 and suffices at $\pi^{\pi_0}=\pi$. The statement about arbitrary repeated colimits also resolves $\Filt^\pi=\Filt'^\pi$ when no fixed-point condition holds. One cannot obtain the latter merely by writing a formal double colimit: a diagram of already presented objects has transition maps which need not have been specified at chosen presentation stages. The bounded closure theorem supplies the necessary control of these maps.

These improvements do not make every cardinal useful for $\mathcal M$. Equality of the two \emph{presentation} filtrations is distinct from equality with the \emph{intrinsic} filtration $\mathcal M_\pi$. The latter compares the ability to lift morphisms out of an object with the size of a presentation inside the chosen Ind completion. The arithmetic condition in the small-subcategory construction still has work to do there.
\end{dcomment}


\textbf{Définition 4.7.8.} Soit $\mathcal{M}$ une $\mathfrak{U}$-catégorie. On appelle $\mathfrak{U}$-\textit{filtration cardinale généralisée} (ou simplement filtration cardinale généralisée, quand $\mathfrak{U}$ est sous-entendu) la donnée d'un cardinal infini $\pi_0$, et d'une filtration $(\mathcal{F}^{\pi})_{\pi\geqslant\pi_0}$ de $\mathcal{M}$ par des sous-catégories strictement pleines $\mathcal{F}^{\pi}$, indexées par les $\mathfrak{U}$-cardinaux $\geqslant\pi_0$, et satisfaisant les conditions a) b) c) énoncées dans 4.7.7, où dans c) on pose
\[
\pi_1 \;=\; \sup(\pi_0,\cardred\mathcal{F}^{\pi_0}) \, . \tag*{(4.7.8.1)}
\]
(Donc si $\cardred\mathcal{F}^{\pi_0}\leqslant 2^{(2^{\pi_0})}$, il n'y a pas à tenir compte de la condition $\pi'^{\pi}\geqslant\pi_1$, cf. 4.7.7.1.) On dit que la filtration cardinale généralisée (fcg) est \textit{accessible}, si elle satisfait de plus la condition d). On dit qu'elle est \textit{canonique} si on a même
\[
\mathcal{F}^{\pi} \;=\; \mathcal{M}_{\pi} \tag*{(4.7.8.2)}
\]
pour tout $\pi\geqslant\pi_0$.

\textbf{Proposition 4.7.9.} \textit{Soient $(\mathcal{F}^{\pi})_{\pi\geqslant\pi_0}$ une fcg sur la $\mathfrak{U}$-catégorie $\mathcal{M}$, et $\pi>\pi_0$ un cardinal tel que $\pi^{\pi_0}=\pi$. Alors on a}
\printedpage{60}% head: A. Grothendieck : Dérivateurs, Ch. XVIII, Édition des Universités de Montpellier II et Paris VII
% foot: Version 17 décembre 2025
\msp{page 177}
\[
\left\{
\begin{aligned}
&\mathcal{M}_\pi\subset\mathcal{F}^{\pi}\quad\text{\textit{(si $\pi=\pi^{\pi_0}$),}}\\
&\text{\textit{donc} }\mathcal{M}_\pi=\mathcal{F}^{\pi}\text{ \textit{si de plus la fcg est} accessible.}
\end{aligned}
\right.
\tag*{(4.7.9.1)}
\]

\dgrblock{dgrb1019}{Reading the generalized filtration axioms}
\begin{dcomment}{exposition}{p176}{Reading the generalized filtration axioms}
Grothendieck's purpose is to retain only the properties of the preceding constructions that are needed to recover the category from its small objects. Two subscripts in the printed statement of 4.7.7 obscure that purpose. In (b), the closure bound is printed $\pi_0$ followed by the editors' question mark; the proof on p.~174 uses $\pi$. In (c), the presenting poset is printed large relative to $\pi_0$; Lemma 4.7.3 gives a poset large relative to $\pi$, and the factorization of a $\pi$-accessible object's identity on p.~177 requires this stronger assertion. We retain the edition's text, but the following discussion of its consequences uses those two repaired conditions explicitly.

In modern notation an accessible generalized filtration thus consists of increasing, essentially small, replete full subcategories $\mathcal F^\pi$, each consisting of $\pi^+$-presentable objects, with the specified bounded closure and generation properties. At a fixed point $\pi^{\pi_0}=\pi$, every object has a $\pi^+$-filtered presentation by objects of $\mathcal F^\pi$. An identity map out of a $\pi^+$-presentable object then factors through one stage, producing a retract. Closure under the permitted colimits supplies its splitting, and hence
\[
 \mathcal F^\pi=\mathcal M_\pi.
\]
This is the content of 4.7.9. The hypothesis on the presenting index cannot be weakened to mere $\pi_0^+$-filteredness in this step: presentability at the larger cardinal says nothing about preservation of an arbitrary less-filtered colimit.

The cardinal estimates in 4.7.7.1 express essential smallness quantitatively. Write $\rho\geq\pi_1$. A presheaf with at most $\rho$ elements at each object of a skeleton $C_0$ can be described by sets chosen inside $\rho$ and functions assigned to the arrows of $C_0$. There are at most $2^\rho$ choices: $|\Fl C_0|\leq\rho$, there are at most $2^\rho$ functions between subsets of $\rho$, and $(2^\rho)^\rho=2^\rho$. Between two such presheaves there are at most
\[
 (\rho^\rho)^{|\Ob C_0|}\leq 2^\rho
\]
natural transformations. Consequently a skeleton of this full category has at most $2^\rho$ arrows. Taking $\rho=\pi$ gives the first estimate for $\Filt^\pi$; taking $\rho=\pi^{\pi_0}$ gives the estimate for $\mathcal M_\pi$ after the inclusion of 4.7.6. This proves the details omitted on p.~173.

Finally, for any $\pi\geq\pi_0$, let $\rho=\pi^{\pi_0}$. Since $\rho^{\pi_0}=\rho$, monotonicity of presentability yields
\[
 \mathcal F^\pi\subseteq\mathcal M_\pi
 \subseteq\mathcal M_\rho=\mathcal F^\rho.
\]
The equality $\mathcal M_\pi=\mathcal M_\rho$ printed in the proof of 4.7.10 must be an inclusion. Even in sets, increasing the cardinal bound admits genuinely larger objects. The conclusion (4.7.10.1) has the correct direction.
\end{dcomment}


\textsc{Démonstration} (\footnote[75]{\texttt{[Dans ce qui suit, $\mathcal{F}^{\pi}$ a parfois été écrit $\mathcal{F}_{\pi}$, souvent corrigé en $\mathcal{F}^{\pi}$ par l'auteur. Nous écrivons $\mathcal{F}^{\pi}$ partout.]}}). Appliquant la condition c) des fcg à un $F\in\Ob\mathcal{M}_\pi$, on trouve que $F$ est facteur direct d'un $x_i\in\Ob\mathcal{F}^{\pi}$. Comme $\mathcal{F}^{\pi}$ est stable par $\varinjlim$ grandes devant $\pi_0$ et de cardinal $\leqslant\pi$, et que \textit{de plus} $\pi>\pi_0$, on conclut que :

\textbf{Lemme 4.7.9.1} \texttt{[deux fois 4.7.9.1]}. \textit{Si $\pi>\pi_0$ (\footnote[76]{Mais $\mathcal{F}^{\pi_0}$ n'est pas forcément stable par facteurs directs, comme on voit en prenant $\mathcal{M}=\Ind_{\pi_0}(C)$, $C$ une petite catégorie non karoubienne, et la filtration cardinale pour les $\Filt'^{\pi}$ -- on note que $\Filt'^{\pi_0}=\Filt^{\pi_0}\simeq e$, vu qu'une catégorie grande devant $\pi_0$ et de cardinal $\leqslant\pi_0$ est cofinale à $e$.}), $\mathcal{F}^{\pi}$ est stable par facteurs directs.}

D'où résulte que $F\in\mathcal{F}^{\pi}$, d'où (4.7.9.1).

\textbf{Corollaire 4.7.10.} \textit{Supposons la fcg accessible. Pour tout cardinal $\pi\geqslant\pi_0$, on a}
\[
\mathcal{F}^{\pi} \;\subset\; \mathcal{M}_\pi \;\subset\; \mathcal{F}^{\pi^{\pi_0}} \; ( \;=\; \mathcal{M}_{\pi^{\pi_0}}) \, . \tag*{(4.7.10.1)}
\]

En effet, $\mathcal{F}^{\pi}\subset\mathcal{M}_\pi$ par définition de fcg accessibles, et $\mathcal{M}_\pi=\mathcal{M}_{\pi^{\pi_0}}$ (puisque $\pi\leqslant\pi^{\pi_0}$), or $\mathcal{M}_{\pi^{\pi_0}}=\mathcal{F}^{(\pi^{\pi_0})}$ en vertu de (4.7.9.1) appliqué à $\pi'=\pi^{\pi_0}$ au lieu de $\pi$ (vu que $\pi'^{\pi_0}=\pi'$), d'où (4.7.10.1). On notera que cette inclusion généralise (4.7.9.1), elle donne $\mathcal{F}^{\pi}=\mathcal{M}_\pi$ si $\pi=\pi^{\pi_0}$.

\textbf{4.7.11.} Il me semble que les seules filtrations cardinales généralisées utiles sont celles qui sont \textit{accessibles.} Par la suite, je vais me borner à ce cas. Comme alors $\mathcal{F}^{\pi}\subset\mathcal{M}_\pi$, on en conclut
\msp{page 178}
que le foncteur canonique
\[
\Ind_\pi(\mathcal{F}^{\pi}) \;\hookrightarrow\; \mathcal{M} \tag*{(4.7.11.1)}
\]
prolongeant l'inclusion $\mathcal{F}^{\pi}\subset\mathcal{M}$, de façon à commuter aux $\varinjlim$ grandes devant $\pi$, est \textit{pleinement fidèle.} De plus, lorsque $\pi^{\pi_0}=\pi$, la condition c) de fcg nous dit entre autres que ce foncteur est essentiellement surjectif, donc que c'est une \textit{équivalence de catégories} :
\[
\Ind_\pi(\mathcal{F}^{\pi}) \xrightarrow{\ \sim\ } \mathcal{M} \qquad (\text{si } \pi=\pi^{\pi_0}). \tag*{(4.7.11.2)}
\]

\textbf{Proposition 4.7.12.} \textit{Soit $(\mathcal{F}^{\pi})_{\pi\geqslant\pi_0}$ une fcg accessible (fcga) sur la $\mathfrak{U}$-catégorie $\mathcal{M}$, et soit $\pi$ un cardinal $\geqslant\pi_0$. Les conditions suivantes sont équivalentes.}
\begin{enumerate}[label=\textit{\alph*)}]
\item \textit{Le foncteur pleinement fidèle (4.7.11.1) est essentiellement surjectif, i.e. est une équivalence de catégories.}
\item \textit{Pour tout $F\in\Ob\mathcal{M}$, $\mathcal{F}^{\pi}/F$ est une catégorie filtrante et grande devant $\pi$.}
\end{enumerate}

\noindent\printedpage{61}% head: A. Grothendieck : Dérivateurs, Ch. XVIII, Édition des Universités de Montpellier II et Paris VII
% foot: Version 17 décembre 2025
\textit{Quand ces conditions sont satisfaites, et si $\pi>\pi_0$, on a}
\[
\mathcal{F}^{\pi} \;=\; \mathcal{M}_\pi \, . \tag*{(4.7.12.1)}
\]

\textsc{Démonstration}. L'équivalence de a) et b) résulte aussitôt de 4.4.2 et de la définition des catégories grandes devant $\pi$ (p.~81). Que ces
\msp{page 179}
conditions impliquent $\mathcal{F}^{\pi}=\mathcal{M}_\pi$, i.e. $\mathcal{M}_\pi\subset\mathcal{F}^{\pi}$ dans le cas où $\pi>\pi_0$, se voit comme (4.7.9.1) en utilisant le lemme 4.7.9.1.

\textbf{Définition 4.7.13.}
\begin{enumerate}[label={$\arabic*^{\circ}$)}]
\item Un cardinal $\pi\geqslant\pi_0$ est dit \textit{cardinal utile} pour la fcga $(\mathcal{F}^{\pi})_{\pi\geqslant\pi_0}$, s'il satisfait les conditions équivalentes a) et b) de 4.7.12, et si de plus dans le cas où $\pi=\pi_0$, on a $\mathcal{F}^{\pi}=\mathcal{M}_\pi$. (Donc si $\pi$ est un cardinal utile, on a toujours $\mathcal{F}^{\pi}=\mathcal{M}_\pi$.) Dans le cas contraire, on dit que $\pi$ est un \textit{cardinal superflu} (pour le fcga donnée). % [sic] le fcga donnée -> la fcga donnée, gender agreement (feminine)
\item Un cardinal $\pi$ est dit \textit{régulier} pour la fcg donnée, s'il est utile, et si de plus la majoration de cardinaux énoncée dans (4.7.7.1) est valable.
\end{enumerate}

\textbf{4.7.14} Ainsi, tous les cardinaux $\pi\geqslant\pi_0$ tels que $\pi^{\pi_0}=\pi$ sont utiles, et même réguliers pour la filtration cardinale \texttt{[généralisée]} donnée. Notons que si tous les cardinaux $\pi\geqslant\pi_0$ sont utiles, alors on a $\mathcal{F}^{\pi}=\mathcal{M}_\pi$ pour tout $\pi\geqslant\pi_0$, donc la fcga est canonique. Notons aussi que dans tous les cas, si $(\mathcal{F}^{\pi})_{\pi\geqslant\pi_0}$ est une fcga, alors $(\mathcal{M}_\pi)_{\pi\geqslant\pi_0}$ en est une aussi.

\noindent\msp{page 180}En effet, la condition a) des fcg (cf. 4.7.7) résulte de la même condition pour $(\mathcal{F}^{\pi})$, et de (4.7.10.1), $\mathcal{M}_\pi\subset\mathcal{F}^{\pi^{\pi_0}}$. La condition b) est une propriété générale de la $\pi$-accessibilité, compte tenu que $(\mathcal{F}^{\pi_0})_0$ est une petite sous-catégorie génératrice formée d'objets $\pi_0$-accessibles de $\mathcal{M}$. La condition d) est satisfaite tautologiquement, et il reste à vérifier la condition c). Mais on a vu que pour $\pi=\pi^{\pi_0}$, on a $\mathcal{M}_\pi=\mathcal{F}^{\pi}$ (4.7.9.1), donc la condition c) pour $(\mathcal{M}_\pi)$ résulte de la même condition pour $(\mathcal{F}^{\pi})$. Ainsi, on trouve :

\dgrblock{dgrb1020}{Recovering the category at a useful cardinal}
\begin{dcomment}{exposition}{p179}{Recovering the category at a useful cardinal}
Grothendieck's useful cardinal is a cardinal at which the chosen level of the filtration actually presents the whole category. This is stronger than the assertion that its objects are small. The distinction explains both the comparison functor in 4.7.11 and the later preference for the canonical filtration.

Let $A\subseteq\mathcal M_\pi$ be an essentially small full subcategory, and assume $\mathcal M$ admits $\pi^+$-filtered colimits. Its extension
\[
 T:\Ind_\pi A\longrightarrow\mathcal M
\]
is fully faithful. Indeed, for presentations $X=\varinjlim_i a_i$ and $Y=\varinjlim_j b_j$ over $\pi^+$-filtered categories,
\[
 \Hom_{\mathcal M}(TX,TY)
 \simeq \varprojlim_i\varinjlim_j\Hom_A(a_i,b_j),
\]
using the universal property in the first variable and presentability in the second. This is also the Hom formula in the Ind category. The extension has all of $\mathcal M$ as its essential image precisely when every object has such a presentation.

The equivalent condition in 4.7.12 uses the canonical presentation. If $A/F$ is $\pi^+$-filtered and the restricted Yoneda functor determined by $A$ is fully faithful, then the canonical map
\[
 \varinjlim_{(a\to F)\in A/F}a\longrightarrow F
\]
is an isomorphism. To see this, apply $\Hom(a_0,-)$ for $a_0\in A$: presentability moves it inside the colimit, and the category-of-elements calculation identifies the resulting set with $\Hom(a_0,F)$. Full faithfulness of restricted Yoneda reflects the isomorphism. If $T$ is already an equivalence, the same category-of-elements description proves the converse immediately. However, density of a fixed-point level does not by itself imply density of every smaller level. The stated generalized-filtration axioms do not supply that missing hypothesis. The next block exhibits a counterexample to (b)$\Rightarrow$(a) in the proposition as printed; the repair is to require the level $A$ to be dense, or to require its canonical cocones to be colimit cocones.

At a useful level every $\pi^+$-presentable object is a retract of a member of $A$. Taking $A=\mathcal F^\pi$ and using its retract closure gives $\mathcal F^\pi=\mathcal M_\pi$. This explains the equality in the endpoint clause of Definition 4.7.13, with usefulness understood through the realization equivalence rather than the faulty bare slice criterion; it is not an additional form of cardinal arithmetic. Nor is ``regular for the filtration'' the usual set-theoretic regularity of the cardinal $\pi$. It means useful together with a specified quantitative presentation bound. The ordinary regular filtering cardinal in all these statements is $\pi^+$, regardless of the cofinality of $\pi$ itself.

Replacing $\mathcal F^\pi$ by $\mathcal M_\pi$ preserves the generalized filtration. Essential smallness follows from $\mathcal M_\pi\subseteq\mathcal F^{\pi^{\pi_0}}$. Bounded closure is the general fact that $\pi^+$-presentables are closed under colimits of diagrams having at most $\pi$ arrows whenever those colimits exist. For the quantitative part, if $F\in\mathcal M_{\pi'}$, first use
\[
 F\in\mathcal F^{(\pi')^{\pi_0}}.
\]
Then apply the old generation bound at the fixed-point level $\pi$. The new index has cardinal at most
\[
 \bigl((\pi')^{\pi_0}\bigr)^\pi=(\pi')^\pi.
\]
This supplies the small substitution omitted from the proof of 4.7.15: equality of the two filtrations at the generating level $\pi$ alone does not identify their arbitrary target levels $\pi'$.
\end{dcomment}


\dgrblock{dgrb1021}{A filtered slice need not be a presentation}
\begin{dcomment}{development}{p179}{A filtered slice need not be a presentation}
Grothendieck wants to recognize a useful level by its slices, so that the Ind realization can be checked without choosing presentations. This recognition needs a density hypothesis. The ordinary category-of-elements theorem applies to a presheaf on the chosen level; identifying its realization with the original object is a separate assertion. The following example shows that the axioms stated in 4.7.7--4.7.8 do not imply it. We checked the general recognition criteria for Ind completions cited above; none dispenses with the generation requirement. The example is supplied as a local audit result, without a priority claim.

\textbf{Counterexample to the slice criterion in 4.7.12}\ours.
For every infinite $\pi_0$, there is an accessible generalized filtration satisfying even the repaired strong forms of (b) and (c), such that $\mathcal F^{\pi_0}/X$ is $\pi_0^+$-filtered for every $X$, but $\Ind_{\pi_0}(\mathcal F^{\pi_0})\to\mathcal M$ is not essentially surjective.

\textit{Proof.} Let $\mathcal M=\Ens$, let $\theta=2^{\pi_0}$, and define
\[
 \mathcal F^\pi=
 \begin{cases}
 \{\varnothing\},&\pi_0\leq\pi<\theta,\\
 \Ens_{\leq\pi},&\pi\geq\theta.
 \end{cases}
\]
Take replete full subcategories; the first is equivalent to the one-object category. This is an increasing exhaustive filtration. Each level is essentially small. It consists of $\pi^+$-presentable sets. Sets have all the required filtered colimits. A filtered colimit of empty sets is empty, while a colimit indexed by a category with at most $\pi$ arrows of sets of size at most $\pi$ again has size at most $\pi$. Thus condition (b), even with the intended bound $\pi$, holds.

If $\pi^{\pi_0}=\pi$, then necessarily $\pi\geq2^{\pi_0}=\theta$. Every set $X$ is the union of its subsets of cardinality at most $\pi$, indexed by the $\pi^+$-directed poset $[X]^{\leq\pi}$. If $X\in\mathcal F^{\pi'}$ with $\pi'>\pi$, the indexing poset has cardinal at most $(\pi')^\pi$. The threshold determined by $\mathcal F^{\pi_0}$ is $\pi_1=\pi_0$, so its proviso is automatic. This verifies condition (c), with the stronger, intended $\pi^+$-filtered index. All accessible generalized-filtration axioms are satisfied.

For every set $X$ there is exactly one map $\varnothing\to X$. Thus $\mathcal F^{\pi_0}/X$ is a one-object one-arrow category and is $\pi_0^+$-filtered. But a filtered colimit of copies of $\varnothing$ remains $\varnothing$. The realization $\Ind_{\pi_0}(\mathcal F^{\pi_0})\to\Ens$ has only the empty set in its essential image, and is not an equivalence.\hfill$\square$

The correction is substantive. A slice may be filtered because the chosen level sees almost nothing. Requiring the level to be dense restores the proof in the preceding block. Alternatively replace condition (b) by the conjunction that the slice is $\pi^+$-filtered and that its canonical diagram has colimit $F$. The later conclusions about useful cardinals should use the equivalence in condition (a); the counterexample does not affect the fixed-point levels, where the generation axiom supplies actual presentations.
\end{dcomment}


\textbf{Proposition 4.7.15.} \textit{Si $(\mathcal{F}^{\pi})_{\pi\geqslant\pi_0}$ est une fcga sur la $\mathfrak{U}$-catégorie $\mathcal{M}$, il en est de même de $(\mathcal{M}_\pi)_{\pi\geqslant\pi_0}$. De plus, tout cardinal utile (resp. régulier) pour la première filtration, l'est pour la seconde (appelée la fcg} canonique associée \textit{à la filtration cardinale} \texttt{[généralisée]} \textit{donnée).}

La dernière assertion résulte du fait que si $\pi$ est utile pour $(\mathcal{F}^{\pi})$, alors on a $\mathcal{F}^{\pi}=\mathcal{M}_\pi$.

\textbf{4.7.16.} Il me semble qu'on peut dire que la fcg canonique associée à une fcga donnée
\msp{page 181}
est toujours ``meilleure'' que celle-ci (si celle-ci n'est canonique, i.e. si les deux ne sont égales), à cause des deux faits suivants :
\begin{enumerate}[label=\alph*)]
\item justement parce qu'elle est \textit{canonique}, i.e. ne dépend d'aucun autre choix que celui du seul cardinal $\pi_0$, et
\item parce qu'il y a a priori plus des cardinaux utiles et de cardinaux réguliers pour la fcg canonique.
\end{enumerate}
Donc au total, il me semble que les filtrations cardinales \texttt{[généralisées]} les seules utiles sont celles qui sont canoniques.

\noindent\printedpage{62}% head: A. Grothendieck : Dérivateurs, Ch. XVIII, Édition des Universités de Montpellier II et Paris VII
% foot: Version 17 décembre 2025
Notons qu'une filtration cardinale (au sens de SGA 4 I 9, que je conserve ici) qui est de plus \textit{accessible} (i.e. $\mathcal{F}^{\pi}\subset\mathcal{M}_\pi$ pour tout $\pi\geqslant\pi_0$), n'est autre qu'une fcga canonique, telle que \textit{tous} les cardinaux sont réguliers (a fortiori utiles) pour ladite filtration. Je présume (sans avoir cependant construit d'exemple) que les filtrations cardinales \texttt{[généralisées]} canoniques des catégories de la forme $\mathcal{M}=\Ind_{\pi_0}(C)$, $C$ une petite catégorie \textit{karoubienne} (pour que $\pi_0$ soit un cardinal utile) ne sont pas forcément des fc (non généralisées, i.e. au sens SGA), i.e. je présume qu'il arrive qu'il y ait des cardinaux $\pi\geqslant\pi_0$ non réguliers, et même superflus.

\noindent\msp{page 182}Il me faut enfin étudier ce qui se passe quand on change le cardinal initial $\pi_0$ d'une fcga.

\dgrblock{dgrb1022}{Useful cardinals and the accessibility spectrum}
\begin{dcomment}{status}{p181}{Useful cardinals and the accessibility spectrum}
Grothendieck suspects that an accessible category may have larger nonuseful cardinals. The relevant modern invariant is its \emph{accessibility spectrum}. The sharp relation $\kappa\mathrel{\triangleleft}\lambda$ between regular cardinals requires every set of size below $\lambda$ to have a cofinal family of its subsets of size below $\kappa$, of cardinality below $\lambda$. It controls the general raising theorem. The equation $\pi^{\pi_0}=\pi$ implies $\pi_0^+\mathrel{\triangleleft}\pi^+$, since a set of size at most $\pi$ has at most $\pi$ subsets of size at most $\pi_0$. A mere inequality between cardinals does not give this control.

The example in \citet[Remark 3.7 and Example 3.8]{Vasey2020Accessible}, attributed there to Ad\'amek--Rosick\'y, settles the existence question in 4.7.16. We give the details, including the explicit successor ranks and the unbounded choice needed for 4.8.6(4). The latter is an elementary verification of the cited criterion, not a priority claim.

Fix $\kappa=\aleph_1$, and let $\mathcal D$ be the category of nonempty $\kappa$-directed posets and monotone maps. For every regular $\lambda\geq\kappa$, its $\lambda$-filtered colimits are computed as filtered colimits of posets: representatives and order relations can be compared at a common stage, and fewer than $\kappa$ elements acquire an upper bound at one stage. We first verify that
\[
 P\text{ is }\lambda\text{-presentable in }\mathcal D
 \quad\Longleftrightarrow\quad |P|<\lambda.
\]
For the forward implication, adjoin a new greatest element $\top$ to $P$. The resulting poset $P^\top$ is the $\lambda$-directed union of the subposets $S\cup\{\top\}$ for $S\subseteq P$, $|S|<\lambda$. Every stage belongs to $\mathcal D$, because it has a greatest element. If $P$ is presentable, the inclusion $P\hookrightarrow P^\top$ factors through one stage. Its image then forces $P\subseteq S$, proving $|P|<\lambda$. Conversely, when $|P|<\lambda$, lift a map from $P$ into a $\lambda$-filtered colimit by choosing representatives for its elements and witnesses for all its order relations. There are fewer than $\lambda$ choices. A common stage makes them into a monotone map. Fewer than $\lambda$ equality witnesses likewise prove essential uniqueness. This is exactly preservation of the colimit by $\Hom(P,-)$.

Every object $P$ of $\mathcal D$ is a $\kappa$-directed union of its subposets having fewer than $\kappa$ elements and a greatest element. Indeed, the union of fewer than $\kappa$ such subposets has fewer than $\kappa$ elements; adjoin an upper bound in $P$ to obtain a larger subposet of the same kind. These subposets cover $P$. Their isomorphism classes form a set and they are $\kappa$-presentable by the criterion just proved. Thus $\mathcal D$ is $\aleph_1$-accessible. An idempotent on one of these presentables splits through its image, which is again $\kappa$-directed and has fewer than $\kappa$ elements. Taking a skeleton $C$ of the presentables therefore gives a small idempotent-complete category with
\[
 \mathcal D\simeq\Ind_{\aleph_0}C.
\]
This meets the Karoubian-base requirement in 4.7.16.

Now let $\theta>\kappa$ be any singular cardinal of countable cofinality, and put $\lambda=\theta^+$. We claim that the inclusion poset $[\theta]^{<\kappa}$ has no cofinal subset of cardinality at most $\theta$. Suppose a candidate family is $\mathscr A=\bigcup_{n<\omega}\mathscr A_n$, where $|\mathscr A_n|\leq\theta_n<\theta$, $\theta_n\geq\kappa$, and $\sup_n\theta_n=\theta$. Since every member has cardinality below $\kappa$, the union of $\mathscr A_n$ has cardinality at most $\theta_n$. Choose $a_n\in\theta$ outside that union. The countable set $\{a_n:n<\omega\}$ lies in $[\theta]^{<\kappa}$ but is contained in no member of $\mathscr A$, a contradiction.

The poset $Q=[\theta]^{<\kappa}$ itself belongs to $\mathcal D$. If it were a $\lambda$-filtered colimit of $\lambda$-presentables, each stage would have at most $\theta$ elements. For each $a<\theta$, choose a stage representing $\{a\}\in Q$. There are fewer than $\lambda$ such choices, so one later stage has image $Y\subseteq Q$ containing all these singletons. A monotone image of a $\kappa$-directed poset is $\kappa$-directed. Therefore $Y$ is cofinal in $Q$: any subset of $\theta$ of size below $\kappa$ is contained in an upper bound of its singletons in $Y$. But $|Y|\leq\theta$, contradicting the preceding paragraph. Thus $\mathcal D$ is not $\theta^+$-accessible.

For a first explicit pair one may take $\theta=\aleph_\omega$. Singular cardinals of countable cofinality are unbounded, so this \emph{one} category has arbitrarily large nonuseful manuscript cardinals $\pi=\theta$, although its initial cardinal $\pi_0=\aleph_0$ is useful. Moreover its small dense subcategory $C$ remains formed of $\theta^+$-presentables and it has $\theta^+$-filtered colimits. Hence it also gives the dense-presentable-generator separation requested in 4.8.6(1). Questions about the presentability ranks or internal sizes of individual objects are different; compare \citet[Definitions 2.5 and 2.10]{LiebermanRosickyVasey2019Sizes}.
\end{dcomment}


\dgrblock{dgrb1023}{Usefulness already gives the quantitative bound}
\begin{dcomment}{exposition}{p181}{Usefulness already gives the quantitative bound}
Grothendieck's distinction between useful and regular cardinals reflects the unresolved presentation problem at this point of the manuscript. Once the retract-free presentation theorem is available, a useful cardinal is regular, and the proviso involving $\pi_1$ can be omitted. We include the argument because the requirement that the index be ordered deserves its own cardinal check.

Suppose $\Ind_\pi(A)\simeq\mathcal M$, where $A=\mathcal F^\pi=\mathcal M_\pi$, and let $F\in\mathcal F^{\pi'}$ with $\pi'>\pi$. Set $\mu=(\pi')^\pi$. Then $\mu^\pi=\mu$, and $F$ is $\mu^+$-presentable. Apply the fixed-point presentation theorem, now with initial cardinal $\pi$ and generating category $A$: it gives a $\pi^+$-filtered category $J$ of at most $\mu$ arrows and a presentation of $F$ by objects of $A$. This estimate is independent of $|\Fl A|$.

To make the index ordered, put $\kappa=\pi^+$. We show that a $\kappa$-filtered category $J$ with $|\Fl J|\leq\mu=\mu^\pi$ has a cofinal functor from a $\kappa$-directed poset $P$ with $|P|\leq\mu$. Begin with a discrete poset labelled by the objects of $J$ and its evident functor $u$ to $J$. Build larger posets through $\kappa$ stages. Every existing element will have fewer than $\kappa$ predecessors. At a successor stage, for each subset $S$ of the current poset of size below $\kappa$, take its downward closure and adjoin a new upper element. This downward closure has size below $\kappa$; choose a cocone on its image in $J$ to define the value of the new element and all new transition maps.

We also impose cofinality. For each pair $p,q$ of existing elements and each fork $j\to u(p)$, $j\to u(q)$ in $J$, adjoin another upper element over the downward closure of $\{p,q\}$, choosing its image cocone to include this fork. The category of this augmented diagram still has fewer than $\kappa$ arrows, so such a cone exists. In both constructions the predecessors of a new element are exactly the selected downward-closed set. Distinct new elements are initially incomparable; composition is defined by the chosen cocone, so the functor remains well defined. At limit stages take unions.

There are at most $\mu^{<\kappa}=\mu$ choices at each stage and $\kappa\leq\mu$ stages. Thus $|P|\leq\mu$. Regularity of $\kappa$ ensures that any subset of $P$ of size below $\kappa$ is present at one earlier stage, and then receives an upper bound, so $P$ is $\kappa$-directed. For each $j\in J$, the comma category $j\backslash P$ is nonempty because the initial stage includes $j$. Any two of its objects are joined to a common object by one of the added fork cones. The functor $P\to J$ is therefore cofinal. Reindexing the presentation of $F$ along it retains the bound
\[
 |P|\leq\mu=(\pi')^\pi.
\]
This is the full ordered-index bound in Definition 4.7.13, with no threshold proviso. Thus useful-but-nonregular cardinals do not occur under the repaired realization definition. Superfluous cardinals still occur, as the accessibility-spectrum discussion explains. The source's distinction is a faithful record of the stage reached in its argument, but should not be carried into a statement of the current theory.
\end{dcomment}


\textbf{Proposition 4.7.17.} \textit{Soit $(\mathcal{F}^{\pi})_{\pi\geqslant\pi_0}$ une fcga sur la $\mathfrak{U}$-catégorie $\mathcal{M}$, et soit $\pi'_0$ un cardinal $\geqslant\pi_0$. Alors $(\mathcal{F}^{\pi})_{\pi\geqslant\pi'_0}$ est une fcga sur $\mathcal{M}$, qui est canonique si celle de départ est canonique. Pour un cardinal $\pi\geqslant\pi'_0$, celui-ci est utile pour $(\mathcal{F}^{\pi})_{\pi\geqslant\pi'_0}$ si et seulement s'il l'est pour $(\mathcal{F}^{\pi})_{\pi\geqslant\pi_0}$. Pour qu'il soit régulier pour $(\mathcal{F}^{\pi})_{\pi\geqslant\pi'_0}$, il} suffit \textit{qu'il le soit pour $(\mathcal{F}^{\pi})_{\pi\geqslant\pi_0}$, et c'est aussi nécessaire quand on suppose de plus que $2^{\pi}\geqslant\pi'_1\overset{\text{déf}}{=}\sup(\pi'_0,\cardred\mathcal{F}^{\pi'_0})$, a fortiori quand $\pi\geqslant\pi'_1$.}

\textsc{Démonstration}. Les conditions a) b) sont trivialement satisfaites par la deuxième filtration, étant satisfaites par la première, ainsi que d) si elle l'est par la première. Il en est de même de la condition c) si $\pi=\pi^{\pi'_0}$, qui implique en effet $\pi^{\pi_0}=\pi$, puisqu'en général
\[
\pi \;\leqslant\; \pi^{\pi_0} \;\leqslant\; \pi^{\pi'_0} \, .
\]
Reste à examiner la majoration de cardinaux
\[
\card I \;\longrightarrow\; \pi'^{\pi} \qquad \text{si } F\in\Ob\mathcal{F}^{\pi'} \ \underline{\text{et}}\ \pi'^{\pi}\geqslant\pi'_1 \, , % [sic] \longrightarrow -> \leqslant, a majoration (inequality) of cardinals is meant
\]
où cette fois $\pi'_1=\sup(\pi'_0,\cardred\mathcal{F}^{\pi'_0})$. En tout état de cause, on a $\pi'_1\geqslant\pi_1$, donc
\msp{page 183}
$\pi'^{\pi}\geqslant\pi'_1$ implique $\pi'^{\pi}\geqslant\pi_1$, donc on peut appliquer la condition c) des fcg pour $(\mathcal{F}^{\pi})_{\pi\geqslant\pi_0}$, pour conclure qu'elle reste encore vérifiée. Donc en somme, \textit{toutes} les conditions des fcg (resp. fcga) sont trivialement vérifiées pour la fcg \textit{restreinte} de la première. De même, le fait qu'un cardinal $\pi'\geqslant\pi'_0$ est utile est trivialement indépendant de quelle des deux filtrations on considère, elle ne dépend que de la donnée du cardinal $\pi'$ et de la sous-catégorie $\mathcal{F}^{\pi'}$ de $\mathcal{M}$. Reste la majoration de cardinaux, dans le cas où on s'intéresse à la \textit{régularité} de $\pi'$ pour les fcga envisagées. La majoration exigée est la même, mais la condition préalable de validité de cette majoration, savoir
\[
\pi'^{\pi} \;\geqslant\; \pi'_1 \qquad \text{resp.} \qquad \pi'^{\pi} \;\geqslant\; \pi_1 \, ,
\]
n'est pas la même. Comme $\pi'_1\geqslant\pi_1$, on voit que la condition préalable est moins restrictive dans le cas de la première filtration, que de la seconde, d'où le ``il suffit'' dans la dernière assertion. Mais lorsque $2^{\pi}\geqslant\pi'_1$, même la plus restrictive est automatique\msp{page 184}ment vérifiée, puisque $\pi'^{\pi}\geqslant 2^{\pi}\geqslant\pi'_1$, et alors la régularité de $\pi$ pour la deuxième filtration implique la régularité pour le première, q.e.d. % [sic] le première -> la première, gender agreement (feminine)

\noindent\printedpage{63}% head: A. Grothendieck : Dérivateurs, Ch. XVIII, Édition des Universités de Montpellier II et Paris VII
% foot: Version 17 décembre 2025
\textbf{Corollaire 4.7.18.} \textit{Supposons $\pi_0$ régulier pour la fcga (\footnote[77]{Il suffit pour cela que $\pi_0$ soit utile, i.e. $\mathcal{M}\xleftarrow{\ \sim\ }\Ind_{\pi_0}(\mathcal{F}^{\pi_0}=\mathcal{M}_{\pi_0})$, cela implique régulier, cf. 4.8.5.}). Si on restreint la fcga $(\mathcal{F}^{\pi})_{\pi\geqslant\pi_0}$ aux cardinaux $\pi\geqslant\pi'_0$, où $\pi'_0\geqslant\pi_1$, alors on trouve une fcga où la majoration de cardinaux (pour $\pi=\pi^{\pi_0}$) est valable sans condition préalable $\pi'^{\pi}\geqslant\pi'_1$ ; plus précisément, on a alors}
\[
\pi'_1 \;\leqslant\; 2^{({\pi'_0}^{\pi_0})} \qquad \text{\textit{(a fortiori $\pi'_1\leqslant 2^{(2^{\pi'_0})}$)},} \tag*{(4.7.18.1)}
\]
\textit{donc la condition préalable $\pi'^{\pi}\geqslant\pi'_1$ est automatiquement vérifiée si $\pi\geqslant\pi_0$, $\pi=\pi^{\pi'_0}$, $\pi'>\pi$ (cf. remarque 4.7.7.2).}

Ceci équivaut donc à la majoration de cardinaux suivante (appliquée au cas de $\pi=\pi'_0$).

\textbf{Lemme 4.7.19.} \textit{Supposons $\pi_0$ régulier pour la fcga $(\mathcal{F}^{\pi})_{\pi\geqslant\pi_0}$. Alors pour tout cardinal $\pi\geqslant\pi_1$ ($\overset{\text{déf}}{=}\sup(\pi_0,\cardred\mathcal{F}^{\pi_0})$), on a}
\[
\cardred\mathcal{F}^{\pi} \;\leqslant\; 2^{(\pi^{\pi_0})} \, , \tag*{(4.7.19.1)}
\]
\textit{et si $\pi'\geqslant\pi_1$ et $F\in\Ob\mathcal{F}^{\pi'}$, on a}
\[
\cardred\mathcal{F}^{\pi}/F \;\leqslant\; \pi'^{(\pi^{\pi_0})} \qquad \text{(\footnotemark[78])} \, . \tag*{(4.7.19.2)}
\]
\footnotetext[78]{\textbf{N.B.} J'ignore si on peut remplacer dans ces majorations $\pi^{\pi_0}$ par $\pi_0$.}

\noindent\msp{page 185}\textsc{Démonstration du lemme}. Comme $\pi_0$ est régulier, posant $C=\mathcal{F}^{\pi_0}=\mathcal{M}_{\pi_0}$, on trouve $\mathcal{M}\simeq\Ind_{\pi_0}(C)$, donc on peut supposer $\mathcal{M}=\Ind_{\pi_0}(C)$, et appliquer les résultats prouvés pour les catégories de cette forme. Comme $\mathcal{F}^{\pi}\subset\mathcal{M}_\pi$, $\mathcal{F}^{\pi'}\subset\mathcal{M}_{\pi'}$, les majorations (4.7.9.1) et (4.7.9.2) sont contenues dans % [sic] (4.7.9.1) et (4.7.9.2) -> (4.7.19.1) et (4.7.19.2), the majorations of Lemme 4.7.19 are meant
\[
\left\{
\begin{array}{rcll}
\cardred\mathcal{M}_\pi & \leqslant & 2^{(\pi^{\pi_0})} & \\
\cardred\mathcal{M}_\pi/F & \leqslant & \pi'^{(\pi^{\pi_0})} & \text{si } F\in\Ob\mathcal{M}_{\pi'},\ \text{avec}\\
 & & & \pi'\geqslant\pi\geqslant\pi_1.
\end{array}
\right\}
\ \text{pour } \pi\geqslant\pi_1 \qquad \text{(\footnotemark[79])}. \tag*{$(*)$}
\]
\footnotetext[79]{Comparer avec majorations $\cardred\Filt^{\pi}\leqslant 2^{\pi}$, $\cardred\Filt^{\pi}/F\leqslant\pi'^{\pi}$, si $F\in\Ob\Filt^{\pi'}$, avec $\pi'\geqslant\pi\geqslant\pi_1$, majorations plus agréables que celles de (4.7.19). Il faudrait quand même s'assurer sur un exemple si les majorations plus sympathiques ne seraient pas valables dans le cas général.}
La première de ces majorations n'est autre que la deuxième majoration dans 4.7.7.1. Notons qu'on a bien
\[
2^{(\pi^{\pi_0})} \;\leqslant\; 2^{(2^{\pi})} \, ,
\]
ce qui résulte en effet de
\[
\pi^{\pi_0} \;\leqslant\; 2^{\pi} \, ,
\]
qu'on prouve en notant que $\pi<2^{\pi}$, d'où $\pi^{\pi_0}\leqslant(2^{\pi})^{\pi_0}=2^{\pi\cdot\pi_0}=2^{\pi}$ puisque $\pi\geqslant\pi_0$. Il reste à établir la deuxième majoration
\[
\cardred\mathcal{M}_\pi/F \;\leqslant\; \pi'^{(\pi^{\pi_0})} \qquad \text{si } F\in\Ob\mathcal{M}_{\pi'},\ \text{avec } \pi'\geqslant\pi\geqslant\pi_1 \, . \tag*{$(*)$}
\]
(On n'en a pas besoin pour 4.7.18, mais ce sera sans doute utile par ailleurs.) On utilise $\mathcal{M}_{\pi'}\subset\Filt^{\pi'^{\pi_0}}$ si $\pi'\geqslant\pi_1$ (4.7.6.2), donc si $F\in\Ob\mathcal{M}_{\pi'}$, on a

\dgrblock{dgrb1024}{Changing the starting point and bounding a slice}
\begin{dcomment}{exposition}{p184}{Changing the starting point and bounding a slice}
Grothendieck wants the filtration to be intrinsic enough that deleting its first levels does not change the useful cardinals which remain. The only part of the definition that remembers the starting point quantitatively is the threshold $\pi_1$. Proposition 4.7.17 records this bookkeeping correctly once the repaired filtration axioms are used.

If $\pi_0'\geq\pi_0$ and $\pi^{\pi_0'}=\pi$, then
\[
 \pi\leq\pi^{\pi_0}\leq\pi^{\pi_0'}=\pi.
\]
The old generation property therefore applies. The new threshold $\pi_1'$ is at least $\pi_1$, so its size proviso is more restrictive and the old bound implies the new one. Usefulness concerns the same full subcategory $\mathcal F^\pi$ and the same realization functor, hence is unchanged. For regularity the converse needs the proviso to be automatic: $2^\pi\geq\pi_1'$ ensures $(\pi')^\pi\geq\pi_1'$ whenever $\pi'>\pi$. No additional category-theoretic argument is hidden here.

The slice estimate in 4.7.19 can be proved without the extra ordering $\pi'\geq\pi$ imposed midway through its proof. Assume only the stated $\pi,\pi'\geq\pi_1$ and set
\[
 a=\pi^{\pi_0},\qquad b=(\pi')^{\pi_0},\qquad
 \Gamma=(\pi')^a.
\]
Choose a skeleton of $C$ and of $\mathcal M_\pi$. Its object set and arrow set have cardinal at most $2^a$. For $G\in\mathcal M_\pi$ and $F\in\mathcal M_{\pi'}$, the preceding bounds give $|G(c)|\leq a$ and $|F(c)|\leq b$. Hence
\[
 \begin{split}
 |\Hom(G,F)|
 &\leq\prod_{c\in\Ob C_0}|F(c)|^{|G(c)|}\\
 &\leq(b^a)^{|\Ob C_0|}
 \leq\bigl((\pi')^{\pi_0}\bigr)^{a\pi_1}
 =(\pi')^a=\Gamma.
 \end{split}
\]
The product, rather than a bare sum, is the direct count of all component maps of a natural transformation; naturality can only reduce the number. The infinite-cardinal identities use $\pi_0,\pi_1\leq a$. Since $\pi'\geq2$, also $2^a\leq\Gamma$. There are therefore at most $\Gamma$ objects in the chosen representative slice, and at most
\[
 \Gamma^2\,2^a=\Gamma
\]
arrows. This proves $\cardred(\mathcal M_\pi/F)\leq(\pi')^{\pi^{\pi_0}}$ in exactly the range of the statement. All counts are counts on skeletons; the literal set of objects of a replete full subcategory need not itself be small.

For 4.7.18 one needs only the first count, with $\pi=\pi_0'$. It gives $\pi_1'\leq2^{(\pi_0')^{\pi_0}}\leq2^{2^{\pi_0'}}$. If $\pi^{\pi_0'}=\pi$, then $\pi\geq2^{\pi_0'}$, and $(\pi')^\pi\geq2^\pi\geq\pi_1'$. This is the precise reason the threshold disappears after increasing the initial level.
\end{dcomment}


% ---- from der-ch18_p066-070.tex
\printedpage{64}
% head: A. Grothendieck : Dérivateurs, Ch. XVIII, Édition des Universités de Montpellier II et Paris VII

\msp{page 186}

\dgrblock{dgrb1025}{Two different proposed improvements of the counting bounds}
\begin{dcomment}{development}{p185}{Two different proposed improvements of the counting bounds}
Grothendieck asks two different questions in the footnotes to 4.7.19. Footnote 78 literally asks whether the exponent $\pi^{\pi_0}$ can be replaced by $\pi_0$. Footnote 79 asks about the more plausible bounds $2^\pi$ and $(\pi')^\pi$. These must be separated: the first has an elementary negative answer already for sets. We record the calculation because it prevents a small printed exponent from being mistaken for a still-open refinement. This is a local verification given here, not a claim of priority.

\textbf{Failure of a uniform initial exponent}\ours.
Fix an infinite $\pi_0$. In $\Ens$, let $\mathcal F^\rho$ be the sets of cardinality at most $\rho$, and let $\pi=\pi'=2^{\pi_0}$. Then
\[
 \cardred\mathcal F^\pi>2^{\pi_0},
 \qquad
 \cardred(\mathcal F^\pi/1)>(\pi')^{\pi_0}.
\]
Thus neither occurrence of $\pi^{\pi_0}$ in 4.7.19 can uniformly be replaced by $\pi_0$.

\textit{Proof.} We first verify the categorical setting. A set $X$ with $|X|\leq\rho$ is $\rho^+$-presentable: a map from $X$ to a $\rho^+$-filtered colimit has at most $\rho$ elementwise witnesses, and all witnesses and equalities can be moved to a common stage. If $|X|>\rho$, its presentation as the union of subsets of cardinality at most $\rho$, indexed by the $\rho^+$-directed poset of those subsets, prevents the identity of $X$ from factoring through a stage. Thus $\mathcal F^\rho=\mathcal M_\rho$. Every set is a filtered union of its subsets of cardinality at most $\pi_0$, so $\pi_0$ is useful for this canonical filtration. A skeleton $C$ of $\mathcal F^{\pi_0}$ has at most $2^{\pi_0}$ arrows and at least that many endomorphisms of a $\pi_0$-element set; hence its reduced cardinal is exactly $2^{\pi_0}$ and $\pi_1=\pi$.

Now choose $X$ with $|X|=\pi$. It has exactly $\pi^\pi=2^\pi$ endomorphisms. Any skeleton containing a representative of its isomorphism class has the same number of endomorphisms of that representative. Therefore $\cardred\mathcal F^\pi\geq2^\pi>\pi=2^{\pi_0}$. In the slice over the singleton, the unique map $X\to1$ has precisely the same endomorphism set. Hence
\[
 \cardred(\mathcal F^\pi/1)\geq2^\pi
 >\pi=(2^{\pi_0})^{\pi_0}=(\pi')^{\pi_0}.
\]
The strict inequalities are Cantor's theorem, and all hypotheses of 4.7.19 are met.\hfill$\square$

In fact the elementary upper count also gives $\cardred\mathcal F^\pi=2^\pi$ in this example, so the first bound proposed in footnote 79 is sharp for sets. That observation does not prove its validity for arbitrary accessible categories. Our checks of the sources cited here, including the bounded-colimit theorem of Chen and the discussion of sizes by Lieberman--Rosick\'y--Vasey, did not locate a proof or counterexample for the full pair of numerical inequalities in footnote 79 with these exact hypotheses. We leave that specific comparison unsettled; we do not label it a recognized open problem of the literature.
\end{dcomment}


\[ \card F(x) \longrightarrow \pi'^{\pi_0} \qquad \forall\, x \in \Ob C , \] % [sic] \longrightarrow -> \leqslant, a cardinal inequality is expected (cf. the next display)

donc si $G \in \Ob \mathcal{M}_\pi \subset \Ob \Filt^{\pi^{\pi_0}}$, d'où
\[ \card G(x) \leqslant \pi^{\pi_0} \qquad \forall\, x \in \Ob C , \]
on trouve
\[ \card \Hom(G,F) \leqslant \underbrace{\underbrace{\card(\Ob C)}_{\leqslant \pi_1}\cdot(\pi'^{\pi_0})^{(\pi^{\pi_0})}}_{\pi'^{\pi_0\cdot\pi^{\pi_0}} = \pi'^{(\pi^{\pi_0})}} , \]
d'où
\[ \card \Ob \mathcal{M}_\pi/F \leqslant \underbrace{\underbrace{\card(\Ob \mathcal{M}_\pi)}_{\leqslant 2^{(\pi^{\pi_0})}}\cdot\pi'^{(\pi^{\pi_0})}}_{\gamma = \pi'^{\pi^{\pi_0}}} , \]
et enfin
\[ \card \Fl \mathcal{M}_\pi/F \leqslant \underbrace{\underbrace{\gamma\cdot\gamma}_{\gamma = \pi'^{(\pi^{\pi_0})}}\cdot\underbrace{\card(\Fl \mathcal{M}_\pi)}_{\leqslant 2^{(\pi^{\pi_0})}}}_{\pi'^{(\pi^{\pi_0})}} , \]
ce qui est la majoration annoncée (la même que celle prouvée dans 4.5.12, la démonstration étant essentiellement la même).

\msp{page 187}

\textbf{4.7.20.} Une fcg canonique est dite \textit{maximale} (où \textit{strictement canonique}) si elle n'est par déduite d'une fcg canonique strictement plus étendue $(\mathcal{F}^\pi)_{\pi\geqslant\pi_{00}}$ (donc avec $\pi_{00} < \pi_0$). Donc si $\mathcal{M}$ admet une fcga, alors elle admet même une fcg strictement canonique, et cette dernière est de plus unique. On va revenir là-dessus dans la section suivante. % [sic] où strictement -> ou strictement, alternative is meant; [sic] n'est par -> n'est pas, negation

\section*{4.8\quad Catégories $\mathfrak{U}$-accessibles}

\textbf{Proposition 4.8.1.} \textit{Soit $\mathcal{M}$ une $\mathfrak{U}$-catégorie, $\pi_0$ un cardinal infini. Considérons les conditions suivantes :}
\begin{itemize}[label={}]
\item[\textit{a)}] \textit{$\mathcal{M}$ satisfait à $L_{\pi_0}$ (i.e. est stable par petites $\varinjlim$ grandes devant $\pi_0$), et contient une petite sous-catégorie pleine $S$} strictement génératrice, \textit{formée d'objets $\pi_0$-accessibles.}
\item[\textit{b)}] \textit{Il existe une petite catégorie $C$, et une équivalence de catégories}
\[ \mathcal{M} \simeq \Ind_{\pi_0}(C) . \]
\item[\textit{c)}] \textit{$\mathcal{M}$ admet une filtration cardinale généralisée accessible $(\mathcal{F}^\pi)_{\pi\geqslant\pi_0}$, ayant $\pi_0$ comme cardinal initial, et telle que $\pi_0$ soit un cardinal utile pour ladite filtration.}
\end{itemize}
% foot: Version 17 décembre 2025

\printedpage{65}
% head: A. Grothendieck : Dérivateurs, Ch. XVIII, Édition des Universités de Montpellier II et Paris VII

\msp{page 188}
\begin{itemize}[label={}]
\item[\textit{c$'$)}] \textit{Comme c), avec $\pi_0$ cardinal régulier (et non seulement utile) pour la filtration cardinale généralisée accessible en question, et cette filtration étant supposée de plus} canonique.
\end{itemize}
\textit{On a alors le diagramme d'implications}
\[ c' \Longleftrightarrow c \Longleftrightarrow b \Longrightarrow a . \]

\textbf{4.8.1.1.} Il me faut ici préciser la terminologie, sur ce que j'entends par sous-catégorie (pleine) $C$ \textit{strictement génératrice} dans $\mathcal{M}$. Pour ceci, je me réfère à SGA 4 I 7.2, (i), aux conditions équivalentes b) et c) de loc. cit. :
\begin{itemize}[label={}]
\item[\circled{b}] $\forall\, F \in \Ob \mathcal{M}$, la flèche canonique
\[ \varinjlim_{i\in C/F} x_i \longrightarrow F \]
est un isomorphisme.
\item[\circled{c}] Le foncteur canonique
\[ \begin{array}{ccc}
\mathcal{M} & \longrightarrow & C^{\wedge} \\[1ex]
F & \longmapsto & (\underbrace{x}_{\in \Ob C} \longmapsto \underbrace{\Hom_{\mathcal{M}}(x,F)}_{\in \Ob (\mathrm{Ens})})
\end{array} \]
est pleinement fidèle.
\end{itemize}
Ces conditions impliquent :
\begin{itemize}[label={}]
\item[a)] $C$ est génératrice dans $\mathcal{M}$ par épimorphismes stricts, i.e. pour tout $F$ dans $\mathcal{M}$, la famille des flèches $x \longrightarrow F$ de source dans $C$, de but $F$, est épimorphique stricte.
\end{itemize}
Et lorsque $\mathcal{M}$ est stable par produits fibrés, cette condition a) est équivalente aux précédentes. Dans loc. cit. j'affirme

\msp{page 189}

\dgrblock{dgrb1026}{Density records all relations}
\begin{dcomment}{exposition}{p188}{Density records all relations}
Grothendieck now separates the existence of sufficiently filtered colimits from the assertion that a small family actually generates every object by such colimits. The distinction is indispensable once arbitrary colimits are unavailable. We continue to write $\lambda=\pi_0^+$: his $\pi_0$-accessible objects are the objects whose covariant Hom functors preserve $\lambda$-filtered colimits. His strictly generating functors are usually called \emph{dense} functors. Thus the principal assertion of 4.8.1 is that a category of the form $\Ind_{\pi_0}(C)$ has a dense small subcategory of $\lambda$-presentable objects. The converse with that very same cardinal requires a separate argument; it is not included in the proposition's displayed implication chain. For an arbitrary generalized filtration, the term useful must retain the comparison-equivalence condition of 4.7.12: filteredness of its initial overcategories alone does not imply density, as the earlier audit explains. The canonical ind-completion examples satisfy the density hypothesis.

The density criterion can be proved without assuming that the ambient category has arbitrary colimits. Let $i:C\to\mathcal M$, and set $N_iX(c)=\Hom(i(c),X)$. A natural transformation $N_iX\to N_iY$ is exactly a compatible family of maps $i(c)\to Y$ indexed by objects $(c,p:i(c)\to X)$ of $C/X$. Consequently
\[
 \operatorname{Nat}(N_iX,N_iY)
 \cong \lim_{(c,p)\in(C/X)^{\mathrm{op}}}\Hom(i(c),Y).
\]
Full faithfulness of $N_i$ says precisely that the evident cocone with vertex $X$ has the colimit universal property. The assertion supplies this particular colimit; it does not presuppose all colimits in $\mathcal M$. This is the content of the equivalence between b) and c) in 4.8.2, and explains the usefulness of the formulation inherited from SGA~4, I, 7.2 \citep{SGA4I1972}.

The passage through strict epimorphic families checks a different-looking compatibility condition. A family $(p_a:X_a\to X)$ is strict epimorphic when maps $u_a:X_a\to Y$ descend uniquely whenever $u_av=u_bw$ for every pair $v:Z\to X_a$, $w:Z\to X_b$ with $p_av=p_bw$. Density implies this property for the family of all maps from $i(C)$: its compatibility condition includes the arrows of $C/X$. Conversely, when $i$ is fully faithful, the compatibility along $C/X$ implies the required compatibility for arbitrary $Z$. Indeed, apply the joint epimorphicity of the family of all $i(c)\to Z$, which follows from the strict epimorphic hypothesis at $Z$, and use fullness to regard their composites with $v,w$ as arrows of $C$. This explains exactly where the full faithfulness hypothesis enters. It also prevents us from replacing a dense subcategory by a mere collection of generators without recording its morphisms.
\end{dcomment}


cette équivalence sans hypothèse restrictive sur $\mathcal{M}$ – mais en annotation marginale j'affirme que cette assertion est \textit{fausse}, en dehors de l'hypothèse d'existence des produits fibrés dans $\mathcal{M}$. Pourtant, je viens de vérifier la démonstration donnée dans loc. cit. de l'implication a $\Longrightarrow$ b, et elle est correcte. Mais peut-être l'annotation résulte-t-elle d'une confusion entre la condition a), et la condition suivante, plus faible a priori, et qui lui ressemble comme un frère :
\begin{itemize}[label={}]
\item[a$'$)] Il existe une famille de flèches $u_i : x_i \longrightarrow F$, de source dans $C$, et de but $F$, qui est épimorphique stricte (cf. SGA 4 I 10.3).
\end{itemize}
On a envie de dire que quand on \textit{agrandisse} une famille de flèches épimorphique stricte, en rajoutant d'autres flèches de même but, elle doit rester épimorphique stricte (“a fortiori”) – ce qui donne l'impression que a$'$) doit impliquer a), donc lui être équivalent. Mais si on essaye de prouver cette propriété hypothétique % [sic] agrandisse -> agrandit, verb form
% foot: Version 17 décembre 2025
\printedpage{66}
% head: A. Grothendieck : Dérivateurs, Ch. XVIII, Édition des Universités de Montpellier II et Paris VII
\msp{page 190}
de stabilité pour les familles épimorphiques strictes, on voit qu'on n'y arrive que si on suppose que la famille initiale est de plus \textit{épimorphique universelle} (si les produits fibrés existent dans $\mathcal{M}$), ou (sous la réserve précédente) si pour tout $x' \longrightarrow F$, le crible de $\mathcal{M}/F'$ image inverse du crible de $\mathcal{M}/F$ engendré par les $x_i \longrightarrow F$ est un “crible épimorphique”. Cela fait entrevoir en effet qu'il doit y avoir des cas où la condition a$'$) est satisfaite, mais non la condition a) (équivalente à b et à c).

Mais nous sommes alors confrontés à la situation très déplaisante : si $C$ est une sous-catégorie strictement génératrice de $\mathcal{M}$, il ne s'ensuit pas pour autant qu'une sous-catégorie pleine $C'$ de $\mathcal{M}$ telle que $C' \supset C$ soit également strictement génératrice\,! Or j'ai toujours utilisé précédemment, sans y réfléchir à deux fois, qui si $C$ est strictement génératrice, il en est de même de $C' \supset C$. Donc il faudra que je sois à présent très attentif, % [sic] qui si -> que si, conjunction
\msp{page 191}
et que je vérifie si cette erreur $\pm$ tacite rend inopérantes certaines démonstrations, voire même, rend faux certains énoncés. Aussi, il me faudrait donner un exemple où $C$ est strictement génératrice, mais $C' \supset C$ ne l'est pas.

Finalement, réflexion faite, ces ennuis craints disparaissent. Je vais énoncer un résultat en forme :

\textbf{Proposition 4.8.2.} \textit{Soient $C$, $\mathcal{M}$ deux $\mathfrak{U}$-catégories, $i : C \longrightarrow \mathcal{M}$ un foncteur, d'où un foncteur}
\[ \left\{\begin{aligned}
&\varphi : \mathcal{M} \longrightarrow C^{\wedge}\\
&\varphi(X) = (x \longmapsto \Hom_{\mathcal{M}}(i(x),X))
\end{aligned}\right. \tag*{(4.8.2.1)} \]
\textit{(foncteur composé $\mathcal{M} \xrightarrow{\textit{can.}} \mathcal{M}^{\wedge} \underbrace{\xrightarrow{i^*}}_{\substack{\textit{i.e. composition}\\ \textit{avec } i}} C^{\wedge}$). Pour tout $X$ dans $\mathcal{M}$, on considère la catégorie $C/X \simeq C/\varphi(X)$, et l'injection} [\,?] \textit{canonique (dans $\mathcal{M}^{\vee}$)}
\[ \varinjlim_{x\in C/X} i(x) \longrightarrow X . \tag*{(4.8.2.2)} \]
\begin{itemize}[label={}]
\item[\textit{1°)}] \textit{Considérons les conditions suivantes.}
\begin{itemize}[label={}]
\item[\circled{a}] \textit{Pour tout $X \in \Ob \mathcal{M}$, la famille des flèches $i(x) \longrightarrow X$, indexée par $\Ob C/X$, est épimorphique stricte.}

\msp{page 192}
\item[\circled{b}] \textit{Pour tout $X \in \Ob \mathcal{M}$, la flèche canonique 4.8.2.2. est un isomorphisme, i.e. fait de $X$ une $\varinjlim$ des foncteurs canoniques $C/X \longrightarrow \mathcal{M}$.}
\item[\circled{c}] \textit{Le foncteur $\varphi$ (4.8.2.1) est pleinement fidèle.}
\end{itemize}
\textit{On a alors les implications}
\[ \mathrm{b} \Longleftrightarrow \mathrm{c} \Longrightarrow \mathrm{a} , \]
\textit{et si on suppose $i$ pleinement fidèle, alors les trois conditions} a, b, c \textit{sont équivalentes.} (\textbf{N.B.} On dit que le foncteur $i : C \longrightarrow \mathcal{M}$ est \textit{strictement générateur}, si les conditions équivalentes b, c sont satisfaites. Si $C$ est une sous-catégorie de $\mathcal{M}$ (pas
% foot: Version 17 décembre 2025
\printedpage{67}
% head: A. Grothendieck : Dérivateurs, Ch. XVIII, Édition des Universités de Montpellier II et Paris VII
nécessairement pleine), on dit que $C$ est \textit{strictement génératrice} si le foncteur d'inclusion est strictement générateur, i.e. satisfait aux conditions équivalentes b), c). Pour ceci, il faut donc que la condition a) soit satisfaite, i.e. que pour tout $X$ dans $\mathcal{M}$, la famille des flèches $x \longrightarrow X$ de source dans $C$, de but $X$, soit épimorphique stricte, et cette condition est aussi suffisante quand $C$ est sous-catégorie pleine de $\mathcal{M}$.)
\item[\textit{2°)}] (\footnote[80]{\texttt{[2° se trouve sur la page 193.]}}). \textit{Supposons que $i$ soit strictement générateur} (\footnote[81]{Il suffit même qu'il satisfasse la condition plus faible a) de 1°).}). \textit{Alors pour tout objet $X$ de $\mathcal{M}$, la famille des flèches $i(x) \longrightarrow X$, indexée par $\Ob C/X$, est non seulement épimorphique stricte (condition a) dans 1°), mais de plus, toute famille de flèches de but $X$, incluant la famille précédente, est également épimorphique stricte.}
\item[\textit{3°)}] \textit{Considérons un diagramme commutatif de foncteurs de $\mathfrak{U}$-catégories :}

\msp{page 193}
\[ \begin{tikzcd}[column sep=small]
C \arrow[rr, "\alpha"] \arrow[dr, "i"'] & & C' \arrow[dl, "i'"] \\
& E &
\end{tikzcd} \]
\textit{Si le foncteur $i$ est strictement générateur} (\footnote[82]{En fait, il suffit qu'il satisfasse la condition plus faible c) de 1°).}), \textit{et si de plus $i'$ est pleinement fidèle, alors $i'$ est strictement générateur.}
\end{itemize}

\textbf{Corollaire 4.8.3.} \textit{Soit $\mathcal{M}$ une $\mathfrak{U}$-catégorie, $C$ une sous-catégorie strictement génératrice, $C'$ une sous-catégorie} pleine \textit{de $\mathcal{M}$ qui contient $C$, alors $C'$ est strictement génératrice.}

C'est en effet un cas particulier de 4.8.2 3°.

\msp{page 194}

\textsc{Démonstration de 4.8.2}. La démonstration de 1°) est celle de SGA 4 I 7.2 (i). (Où on suppose que $C$ est une sous-catégorie pleine de $\mathcal{M}$, mais la démonstration donnée prouve en fait l'énoncé donné ici.) Prouvons 2°). Rajoutons donc à la famille des
\[ i(x) \longrightarrow X \qquad (x \in \Ob C/X) \]
une famille
\[ X_\alpha \xrightarrow{p_\alpha} X \qquad (\alpha \in A) \]
de nouvelles flèches, et prouvons que la famille totale ainsi obtenue est encore épimorphique stricte. On doit donc prouver que si $Y$ est un objet de $\mathcal{M}$, et si on a une famille de flèches $u_x : i(x) \longrightarrow Y$, $u_\alpha : X_\alpha \longrightarrow Y$, satisfaisant la condition de compatibilité habituelle pour des couples de flèches :
\[ (v_x,v_y) \qquad\qquad \begin{tikzcd}
& i(x) \arrow[dr, dotted, "p_x"] & \\
Z \arrow[ur, "v_x"] \arrow[dr, "v_y"'] & & X\ , \\
& i(y) \arrow[ur, dotted, "p_y"'] &
\end{tikzcd} \]
ou
% foot: Version 17 décembre 2025
\printedpage{68}
% head: A. Grothendieck : Dérivateurs, Ch. XVIII, Édition des Universités de Montpellier II et Paris VII
\[ (v_x,v_\alpha) \qquad\qquad \begin{tikzcd}
& i(x) \arrow[dr, dotted, "p_x"] & \\
Z \arrow[ur, "v_x"] \arrow[dr, "v_\alpha"'] & & X\ , \\
& X_\alpha \arrow[ur, dotted, "p_\alpha"'] &
\end{tikzcd} \]
\msp{page 195}
ou
\[ (v_\alpha,v_\beta) \qquad\qquad \begin{tikzcd}
& X_\alpha \arrow[dr, dotted] & \\
Z \arrow[ur, "v_\alpha"] \arrow[dr, "v_\beta"'] & & X\ , \\
& X_\beta \arrow[ur, dotted] &
\end{tikzcd} \]
savoir les relations
\[ \left\{\begin{aligned}
u_xv_x &= u_yv_y\\
u_xv_x &= u_\alpha v_\alpha\\
u_\alpha v_\alpha &= u_\beta v_\beta ,
\end{aligned}\right. \tag*{($*$)} \]
alors il existe une flèche
\[ u : X \longrightarrow Y \]
(nécessairement unique, d'après la propriété c) d'épimorphicité stricte de la famille initiale indexée par $C/X$) telle qu'on ait
\[ \left\{\begin{aligned}
up_x &= v_x \quad \text{pour tout } x \in \Ob C/X\\ % [sic] v_x -> u_x, the arrow u p_x goes to Y (cf. up_\alpha = u_\alpha below)
up_\alpha &= v_\alpha \quad \text{pour tout } \alpha \in A. % [sic] v_\alpha -> u_\alpha, cf. the relation up_\alpha = u_\alpha below
\end{aligned}\right. \tag*{($*$)} \]
Or d'après la propriété a), il existe un unique $u : X \longrightarrow Y$, de telle façon que la première des relations ($*$) soit satisfaite, pour tout $x \in \Ob C/X$. Donc il s'agit seulement de prouver que pour cet $u$, on a pour tout $\alpha$ la relation

\msp{page 196}
\[ up_\alpha = u_\alpha . \]
Or, en appliquant la propriété a) à $X_\alpha$ au lieu de $X$, on veut que la famille des flèches % [sic] veut -> voit, sense of the sentence
\[ i(y) \xrightarrow{q_y} X_\alpha \]
(indexée par $C/X_\alpha$) est épimorphique stricte, a fortiori épimorphique. Donc pour prouver $up_\alpha = u_\alpha$, il suffit de prouver qu'on a
\[ (up_\alpha)q_y = u_\alpha q_y \tag*{($**$)} \]
pour toute $y$ dans $C/X_\alpha$. Mais c'est là une conséquence de la compatibilité ($*$) (deuxième formule), appliquée au cas suivant,
% foot: Version 17 décembre 2025

% ---- from der-ch18_p071-075.tex
\printedpage{69}
% head: A. Grothendieck : Dérivateurs, Ch. XVIII, Édition des Universités de Montpellier II et Paris VII

\[
\begin{tikzcd}[column sep=small]
 & i(y) \arrow[dr, dotted, "p_\alpha q_y=p_x"] & \\
Z=i(y) \arrow[ur, "\mathrm{id}"] \arrow[dr, "q_y"'] & & X \ , \\
 & X_\alpha \arrow[ur, dotted, "p_\alpha"'] &
\end{tikzcd}
\]
on a donc, en posant $x=(y,p_\alpha q_y)\in\Ob C/X$,
\[
\underbrace{u_x\circ\mathrm{id}_{i(y)}}_{u_x} \ = \ u_\alpha q_y \ ,
\]
donc pour prouver $(**)$, il reste à vérifier
\[
u_x \ = \ \underbrace{(up_\alpha)q_y}_{u(p_\alpha q_y)\,=\,up_x} \ ,
\]
\msp{page 197}
qui est en effet vérifiée par définition de $u_x$. Cela achève de prouver $2^\circ$.

Prouvons $3^\circ$. Comme $i'$ est pleinement fidèle, en vertu de $1^\circ$, pour prouver que $i'$ est strictement génératrice, il suffit d'établir que pour tout objet $X$ de $\mathcal{M}$, la famille des flèches $i'(x')\longrightarrow X$, indexée par $\Ob C'/X$, est épimorphique stricte. Or cette famille contient la famille des $\underbrace{i'(\alpha(x))}_{=\,i(x)}\longrightarrow X$, indexée par $\Ob C/X$. Donc il résulte de l'hypothèse sur $i$, et de $2^\circ$, que cette famille est épimorphique stricte. On a gagné ! Ouf \dots

\textsc{Démonstration de 4.8.1.} On a
\[
\mathrm{c}' \ \Longrightarrow \ \mathrm{c} \ \Longrightarrow \ \mathrm{b} \ \Longrightarrow \ \mathrm{a} \, ,
\]
la première implication tautologique, la seconde résultant du fait que $\pi_0$ cardinal régulier pour la fcga envisagée, implique que
\[
\mathcal{M} \ \xleftarrow{\ \sim\ } \ \Ind_{\pi_0}(\mathcal{F}^{\pi_0}) \ ,
\]
$\mathcal{F}^{\pi_0}$ équivalente à une petite catégorie, d'où b).

\msp{page 198}
Enfin b) $\Longrightarrow$ a) résulte du fait que $\Ind_{\pi_0}(C)$ est stable par $\varinjlim$ grandes devant $\pi_0$ (4.4.15), et que $C$ est une sous-catégorie strictement génératrice de $\Ind_{\pi_0}(C)$, formée d'objets $\pi_0$-accessibles. Il nous faut expliciter le

\dgrblock{dgrb1027}{The apparent difficulty with enlargement disappears}
\begin{dcomment}{exposition}{p197}{The apparent difficulty with enlargement disappears}
Grothendieck is checking that his earlier constructions survive enlargement of the generating category. The worry on manuscript pages 190--191 is withdrawn by Proposition 4.8.2; it is not an unresolved defect in the preceding theory. The useful fact is stronger than the permanence of density. Suppose that for every object $Z$, the family of all maps $i(c)\to Z$ is strict epimorphic. Adjoin arbitrary maps $q_a:Z_a\to X$ to the canonical family over $X$. Given compatible descent data on the enlarged family, its restriction supplies a unique $u:X\to Y$. To prove $uq_a=u_a$, precompose with every $i(c)\to Z_a$. Each resulting equality is already among the compatibility equations of the enlarged family. These maps are jointly epimorphic, so $uq_a=u_a$. Thus this particular strict epimorphic family remains strict after enlargement, with no pullback hypothesis.

There is also a short proof directly with natural transformations. If $C\subset C'\subset\mathcal M$, $C'$ is full, and $C$ is dense, a transformation $N_{C'}X\to N_{C'}Y$ restricts to a unique arrow $u:X\to Y$ on $C$. For $p:c'\to X$ and every $a:c\to c'$, naturality gives
\[
 \eta_{c'}(p)a=\eta_c(pa)=upa.
\]
The maps from $C$ to $c'$ are jointly epimorphic, hence $\eta_{c'}(p)=up$. Faithfulness follows by restriction. Therefore $C'$ is dense. This proof isolates the global hypothesis missing from the tempting assertion that \emph{every} strict epimorphic family is preserved by arbitrary enlargement: the testing families are available at every newly adjoined domain.

Two local notational slips should not obscure this argument. In the display on manuscript page~195 the required equations are $up_x=u_x$ and $up_\alpha=u_\alpha$, not equations with the auxiliary comparison maps $v_x,v_\alpha$. Footnote~82 should refer to condition a), as does footnote~81: condition c) is the density condition itself, so it is not a weaker assumption. Both readings are present in the edition and are retained in the transcription.
\end{dcomment}


\textbf{Lemme 4.8.4.} \textit{Soit $C$ une $\mathfrak{U}$-catégorie, $\pi_0$ un cardinal infini. Alors $C$ est strictement génératrice dans $\mathcal{M}=\Ind C$ et dans $\mathcal{M}=\Ind_{\pi_0}C$.}

\msp{page 199}
En effet, le foncteur canonique $\mathcal{M}\longrightarrow C^\wedge$ est pleinement fidèle, et on applique le critère c) de 4.8.2.

Pour achever de prouver 4.8.1, il reste donc à établir l'implication
\[
\mathrm{a} \ \Longrightarrow \ \mathrm{c}' \, .
\]

\printedpage{70}
% head: A. Grothendieck : Dérivateurs, Ch. XVIII, Édition des Universités de Montpellier II et Paris VII

Supposons donc que $\mathcal{M}$ satisfasse à $\mathrm{L}_{\pi_0}$, et admette une petite sous-catégorie strictement génératrice $S$, formée d'objets $\pi_0$-accessibles, donc un foncteur canonique
\[
\Ind_{\pi_0}S \ \longrightarrow \ \mathcal{M} \ ,
\]
prolongeant l'inclusion $S\hookrightarrow\mathcal{M}$, de façon à commuter aux $\varinjlim$ grandes devant $\pi_0$, et du fait que $S\subset\mathcal{M}_{\pi_0}$, ce foncteur est \textit{pleinement fidèle.} Mais comment prouver qu'il est essentiellement surjectif ? Ce ne sera d'ailleurs pas vrai sans autre hypothèse sur $S$. L'idée naturelle est de saturer $S$ pour les $\varinjlim$
\msp{page 200}
de cardinal $\leqslant\pi_0$ représentables dans $\mathcal{M}$, et notamment par facteurs directs, mais même alors (et même en admettant que $\mathcal{M}_{\pi_0}$ lui-même est équivalent à une petite catégorie, et prenant $C=\mathcal{M}_{\pi_0}$), il n'est pourtant pas clair que le foncteur pleinement fidèle $\Ind_{\pi_0}C\longrightarrow\mathcal{M}$ soit essentiellement surjectif, ce qui équivaudrait aussi à dire que pour tout $F$ dans $\mathcal{M}$, $C/F$ est grande devant $\pi_0$. Il n'est en effet pas même clair que $C/F$ soit filtrante.

Donc finalement, je doute qu'on ait a) $\Longrightarrow$ c$'$) ou seulement a) $\Longrightarrow$ b), la situation n'est pas bien comprise visiblement.

Mais prouvons quand même
\[
\mathrm{b} \ \Longrightarrow \ \mathrm{c}' \, .
\]
Pour ceci, on peut supposer $\mathcal{M}=\Ind_{\pi_0}C$, $C$ petite. On sait alors que la filtration de $\mathcal{M}$ par les $(\mathcal{M}_\pi)_{\pi\geqslant\pi_0}$ est une fcga (prop. 4.7.7 et définition 4.7.8), et même une fcg \textit{canonique.} De plus, je dis que le cardinal $\pi_0$ est \textit{utile} et même régulier pour cette filtration (ce qui achèvera de prouver c$'$). On a un résultat un peu plus précis :

\msp{page 201}
\textbf{Corollaire 4.8.5.} \textit{Soient $C$ une catégorie équivalente à une petite catégorie, $\pi_0$ un cardinal infini, $\mathcal{M}=\Ind_{\pi_0}(C)\subset C^\wedge$. Alors :}
\begin{enumerate}
\item[\textit{1$^\circ$)}] $\pi_0$ \textit{est un cardinal} utile \textit{et même régulier pour la fcg canonique} $(\mathcal{M}_\pi)_{\pi\geqslant\pi_0}$ \textit{de} $\mathcal{M}$.
\item[\textit{2$^\circ$)}] $\mathcal{M}_{\pi_0}=C^\natural$ \textit{(sous-catégorie pleine de $\mathcal{M}=\Ind_{\pi_0}C$ formée des facteurs directs dans $\mathcal{M}$ d'objets de $C$), et $\mathcal{M}_{\pi_0}$ est une enveloppe de Karoubi de $C$.}
\end{enumerate}
(\footnote{\textbf{Corollaire 4.8.5.1.} \textit{Si $C$ est karoubienne, alors $\mathcal{M}_{\pi_0}\subset C^\wedge$ est la catégorie formée des foncteurs} représentables \textit{dans} $C\overset{\sim}{\hookrightarrow}\mathcal{M}_{\pi_0}$.}, \footnote{\textbf{4.8.5.2.} Expliciter aussi ceci : Soit $u : C\longrightarrow C'$ un foncteur, avec $C$, $C'$ équivalentes à des petites catégories. Conditions équivalentes :
\begin{enumerate}
\item[a)] $u$ induit $\Kar C\xrightarrow{\ \sim\ }\Kar C'$.
\item[a$'$)] $u$ pleinement fidèle, et tout objet de $C'$ est isomorphe à un facteur direct d'un objet de l'image essentielle.
\item[b)] $\Ind(u) : \Ind(C)\longrightarrow\Ind(C')$ est une équivalence.
\item[c)] $\Ind_{\pi_0}(u) : \Ind_{\pi_0}(C)\longrightarrow\Ind_{\pi_0}(C')$ est une équivalence ($\pi_0$ un cardinal infini donné).
\item[d)] $u_! : C^\wedge\longrightarrow C'^\wedge$ est une équivalence.
\item[d$'$)] $u_* : C'^\wedge\longrightarrow C^\wedge$ est une équivalence \texttt{[plutôt} $u_* : C^\wedge\longrightarrow C'^\wedge$\texttt{]}.
\item[d$''$)] $u^* : C'^\wedge\longrightarrow C^\wedge$ est une équivalence.
\end{enumerate}}).

\printedpage{71}
% head: A. Grothendieck : Dérivateurs, Ch. XVIII, Édition des Universités de Montpellier II et Paris VII

\textsc{Démonstration.}

$1^\circ$) On sait déjà que
\[
\Ind_{\pi_0}\mathcal{M}_{\pi_0} \ \longrightarrow \ \mathcal{M} \tag*{$(*)$}
\]
est pleinement fidèle, et pour voir que c'est essentiellement surjectif, on note que l'on a
\[
\mathcal{M} \ \simeq \ \Ind_{\pi_0}C \ \hookrightarrow \ \Ind_{\pi_0}\mathcal{M}_{\pi_0} \ \longrightarrow \ \mathcal{M} \ ,
\]
le composé étant isomorphe au foncteur identique, est essentiellement surjectif, donc aussi $(*)$. Cela prouve donc que $\pi_0$ est un cardinal utile pour la filtration envisagée, prouvons que ce cardinal est même régulier, donc que pour
\[
F \ \in \ \Ob\mathcal{M}_\pi \ , \qquad \text{avec } \pi>\pi_0,\ \pi^{\pi_0}\geqslant\pi_1
\]
on a
\[
F \ = \ \varinjlim_I x_i \ , \qquad x_i\in\Ob\mathcal{M}_{\pi_0},\ I \text{ grande devant } \pi_0,\ \card I\leqslant\pi^{\pi_0}.
\]
\msp{page 202}
Pour ceci, on utilise (4.7.6.2), page 171,
\[
\mathcal{M}_\pi \ \subset \ \Filt^{\pi^{\pi_0}} \qquad (\text{si } \pi\geqslant\pi_0,\ \pi^{\pi_0}\geqslant\pi_1),
\]
et la définition de $\Filt^{\pi'}$ (pour $\pi'\geqslant\pi_0$) nous montre que c'est là la conclusion cherchée.

Prouvons donc $2^\circ$). Comme $\mathcal{M}_{\pi_0}$ est stable par les $\varinjlim$ inductives finies qui sont représentables dans $\mathcal{M}$ (SGA 4 I 9.9), il est stable par facteurs directs, et comme il contient $C$, on a donc $C^\natural\subset\mathcal{M}_{\pi_0}$. L'inclusion inverse $\mathcal{M}_{\pi_0}\subset C^\natural$ se voit de la façon standard, en écrivant un objet $F$ de $\mathcal{M}_{\pi_0}$ comme $\varinjlim\limits_I x_i$, les $x_i$ dans $C$, $I$ grande devant $\pi_0$, alors $F$ est facteur direct d'un des $x_i$.

Le fait que $C^\natural$ est une \textit{enveloppe de Karoubi} de $C$ dans $\mathcal{M}\subset C^\wedge$ provient du fait que $\mathcal{M}$ est stable par facteurs directs, et que $C$ est une sous-catégorie pleine de $\mathcal{M}$, cf. SGA 4 I (\,?\,\dots).

\textbf{4.8.6. Questions ouvertes.} On a ainsi prouvé 4.8.1 revu et corrigé. Il reste la question :
\begin{enumerate}
\item[1)] Donner un exemple où 4.8.1 a) est vérifié, mais non b) etc. (\footnote{Il suffit de trouver un cardinal non régulier, pour la fcg canonique d'une $\mathcal{M}=\Ind_{\pi_0}C$.}).
\item[2)] Dégager des conditions supplémentaires serviables qui, jointes à a), impliquent b), donc c, c$'$.
\end{enumerate}

\msp{page 203}
Signalons aussi pour mémoire les questions suivantes.
\begin{enumerate}
\item[3)] Soit $\mathcal{M}=\Ind_{\pi_0}(C)$, et $\pi>\pi_0$. Alors $\Filt^\pi$ est-il stable par facteurs directs ? Cela est-il vrai du moins si $\pi=\pi^{\pi_0}$, ou ce qui revient au même (cf. 4.7.3), a-t-on dans ce cas $\Filt^\pi=\Filt'^{\pi}$ ? Si oui, il serait inutile dans la description 4.7.7 des filtrations cardinales, dans la majoration de cardinaux (4.7.7.1), $\card I\leqslant\pi'^{\pi}$, de mettre l'hypothèse restrictive $\pi'^{\pi}\geqslant\pi_1$.
\end{enumerate}

\dgrblock{dgrb1028}{The compact part is an idempotent completion}
\begin{dcomment}{exposition}{p202}{The compact part is an idempotent completion}
Grothendieck identifies the part of an ind-completion that is already present at its initial cardinal. This is the first place where the completion by retracts must be treated as a structural operation. For $\mathcal M=\Ind_{\pi_0}(C)$, every representable is $\lambda$-presentable, and so is every retract of a representable. Conversely, write a $\lambda$-presentable $F$ as a $\lambda$-filtered colimit $\varinjlim_i y(c_i)$. Preservation of that colimit by $\Hom(F,-)$ makes $1_F$ factor through one of its terms:
\[
 F\xrightarrow{s}y(c_i)\xrightarrow{r}F,
 \qquad rs=1_F.
\]
Hence $\mathcal M_{\pi_0}$ is exactly the retract closure of $y(C)$. Idempotents split in $\mathcal M$: a filtered presentation of the splitting is obtained by iterating the idempotent along an ordinal of cofinality greater than $\pi_0$, or by the finite idempotent indexing category used earlier in this commentary. Since idempotents split in the presheaf category, their splitting objects lie in the ind-completion. We have therefore proved
\[
 \mathcal M_{\pi_0}\simeq\Kar(C),\qquad
 \Ind_{\pi_0}(\mathcal M_{\pi_0})\simeq\mathcal M.
\]
The second equivalence follows either from the same presentations or by factoring the identity through the fully faithful comparison functor. It is the useful-cardinal assertion of 4.8.5; the separate quantitative notion called regularity in the manuscript additionally uses the presentation-size estimate from 4.7.6.

The list in footnote~84 is the corresponding Morita invariance statement. A fully faithful $u:C\to C'$ whose image contains every object of $C'$ up to retract induces an equivalence on idempotent completions. Restriction then induces an equivalence of presheaf categories: any set-valued presheaf extends uniquely, up to unique isomorphism, across a split idempotent, because its value at that retract must be the image of the induced idempotent of a set. The left and right Kan extensions are the inverse equivalences. They also preserve the indicated ind-subcategories. Conversely, an equivalence induced on $\Ind_{\pi_0}$ identifies its $\lambda$-presentable objects, and therefore identifies $\Kar(C)$ with $\Kar(C')$. The same argument at $\lambda=\aleph_0$ gives ordinary $\Ind$. This proves the equivalences in the footnote; the variance of $u_*$ must be $C^\wedge\to C'^\wedge$, as the editors already indicate.

The modern higher-categorical formulation identifies the compact objects of an ind-completion with the idempotent completion of its input \citep[Lemma~5.4.2.4]{Lurie2009HTT}. The ordinary proof above shows which part of that statement is already fully present here: the identity of a compact object factors through a single stage, and the resulting projector is the only extra object that completion has introduced at this cardinal.
\end{dcomment}


Pour avoir une réponse affirmative à cette question, pour $\pi$ donné, il suffirait (et probablement faudrait) d'avoir une réponse affirmative à celle-ci : Soient $\mathcal{D}$, $\mathcal{D}_0$ deux catégories grandes devant $\pi_0$, $\mathcal{D}_0$ facteur direct de $\mathcal{D}$, $\mathcal{D}$ admettant une catégorie cofinale $I$ grande devant $\pi_0$, de cardinal $\leqslant\pi$. Alors en est-il de même de $\mathcal{D}_0$ ? Si $p$ est une rétraction de $\mathcal{D}$ sur $\mathcal{D}_0$, la difficulté est que je ne sais pas si $p$ est foncteur cofinal. En fait, je vois que
\printedpage{72}
% head: A. Grothendieck : Dérivateurs, Ch. XVIII, Édition des Universités de Montpellier II et Paris VII
\msp{page 204}
ce n'est pas toujours le cas, même si $\mathcal{D}$ est finie. Prenons p.ex. pour $\mathcal{D}_0$ une catégorie ayant objet final $e$, et une flèche $x\longrightarrow e$ qui n'est pas un isomorphisme, et admettant une section (p.ex. $\mathcal{D}_0$ la sous-catégorie pleine de $\boldsymbol{\Delta}$ formée par $\boldsymbol{\Delta}_0$, $\boldsymbol{\Delta}_1$ : \begin{tikzcd}[cramped] \boldsymbol{\Delta}_0 \arrow[r, shift left, "\partial_0"] \arrow[r, shift right, "\partial_1"'] & \boldsymbol{\Delta}_1 \arrow[l, bend left=60, "\kappa"] \end{tikzcd} ). Soit $\mathcal{D}$ déduit de $\mathcal{D}_0$ en ajoutant un objet final $\tilde{e}$. Alors $\mathcal{D}_0$, $\mathcal{D}$ sont filtrantes, et on obtient une rétraction de $\mathcal{D}$ sur $\mathcal{D}_0$, envoyant $\tilde{e}$ sur $x$. Ce foncteur n'est pas cofinal, car un foncteur cofinal doit envoyer $\tilde{e}$ en $e$.

Plus généralement, si $\mathcal{D}_0$ est une catégorie telle qu'il existe une flèche
\[
\alpha : e_{\mathcal{D}_0^\wedge} \ \longrightarrow \ x_0 \tag*{$(*)$}
\]
dans $\mathcal{D}_0^\wedge$, avec $x_0\in\Ob\mathcal{D}_0$, alors on voit que $\mathcal{D}_0$ est grande devant $\pi_0$ quel que soit le cardinal infini $\pi_0$, et que d'autre part, désignant encore par $\mathcal{D}$ la catégorie déduite de $\mathcal{D}_0$ en lui ajoutant un objet final, il y a une rétraction de $\mathcal{D}$ sur $\mathcal{D}_0$ qui envoie $e_{\mathcal{D}}$ dans $x_0$. Cette rétraction n'est cofinale que si $x_0$ est un objet final de $\mathcal{D}_0$, i.e. si et seulement si $(*)$ est un isomorphisme. Supposons qu'il n'en soit pas ainsi. Mais il y a un foncteur cofinal de la sous-catégorie $I_0$ de $\mathcal{D}_0$, telle que $\Ob I_0=\{x_0\}$, $\Fl I_0=\End_{I_0}(x_0)=\{p,\mathrm{id}_{x_0}\}$, où $p$ est le composé $x_0\xrightarrow{\text{can.}}e_{\mathcal{D}_0^\wedge}\xrightarrow{\alpha}x_0$. Donc pas de contre-exemple de cette façon !

\msp{page 205}
\begin{enumerate}
\item[4)] Considérons $\mathcal{M}=\Ind_{\pi_0}(C)$, $C$ petite, et considérons la filtration canonique $(\mathcal{M}_\pi)_{\pi\geqslant\pi_0}$. Donner des exemples où il y a des cardinaux non utiles, ou mieux, ou pour tout cardinal $\pi$, existe un cardinal $\pi'>\pi$ qui soit non utile. (Cela donnera l'exemple demandé dans $1^\circ$.) D'autre part, donner des exemples de cardinaux utiles qui ne soient réguliers -- si possible, comme tantôt, de tels cardinaux aussi grands qu'on veut.
\end{enumerate}

\dgrblock{dgrb1029}{A retraction can fail cofinality without losing a small presentation}
\begin{dcomment}{examples}{p204}{A retraction can fail cofinality without losing a small presentation}
Grothendieck's example separates two operations that are easy to confuse: retracting a category and choosing a final functor into it. Take the full subcategory $D_0\subset\boldsymbol\Delta$ on $[0]$ and $[1]$. Its terminal object is $[0]$, while $[1]\to[0]$ has two sections. Choose one, $s:[0]\to[1]$. Adjoin a new terminal object $t$ to $D_0$, obtaining $D$. The maps
\[
 a\longrightarrow[0]\xrightarrow{s}[1]
 \qquad(a\in D_0)
\]
form a natural cocone, so they define a retraction $p:D\to D_0$ with $p(t)=[1]$ and $p|_{D_0}=1$. This $p$ is not final. If it were, then for every $a\in D_0$ the colimit of $\Hom(a,p(-))$ would be a singleton. Since $D$ has terminal object $t$, that colimit is $\Hom(a,[1])$. Taking $a=[0]$ gives two elements.

Nevertheless a finite final indexing category remains. More generally a section $1\to y(x)$ of the terminal presheaf map $y(x)\to1$ determines an idempotent $e:x\to x$. For every arrow $v:a\to x$, naturality gives $ev=\alpha_a$, independent of $v$. Let $E$ have one object and endomorphisms $1,e$, with $e^2=e$. The functor $E\to D_0$ picking $x$ is final: in $(a\downarrow E)$ every object $v:a\to x$ maps, using $e$, to $\alpha_a$. This comma category is nonempty and connected. The category $E$ is itself large in front of every cardinal: the constant components $e$ give a cocone on every diagram in $E$. Thus a failed retraction argument supplies no counterexample to bounded ind-presentations.

Question 4.8.6(3) is answered affirmatively by the retract-presentation argument in our commentary on \S4.7. The conclusion holds for the manuscript's indexing convention, which allows non-ordered filtered categories. If one insists on a poset presentation, a finite indexing category with a nontrivial idempotent can no longer be used without changing its size; this is precisely why the indexing convention belongs in the statement. At fixed points $\pi^{\pi_0}=\pi$, the equality $\Filt^\pi=\Filt'^\pi$ then follows from 4.7.3. The example here explains why the proof should work on presentations themselves, rather than assuming that an arbitrary categorical retraction is final.
\end{dcomment}


\textbf{Définition 4.8.7.} Soit $\mathcal{M}$ une $\mathfrak{U}$-catégorie, $\pi_0$ un cardinal infini. On dit que $\mathcal{M}$ est $\pi_0$-\textit{accessible} si $\mathcal{M}$ est équivalente à une catégorie $\Ind_{\pi_0}(C)$, $C$ une petite catégorie, i.e. si elle satisfait aux conditions équivalentes b, c, c$'$ de 4.8.1 (impliquant la condition a), sans doute strictement plus faible). On dit que $\mathcal{M}$ est accessible s'il existe un cardinal tel que $\mathcal{M}$ soit $\pi_0$-accessible (\textbf{N.B.} si $\pi_0$ est un tel cardinal, alors tout cardinal $\pi>\pi_0$ tel que $\pi^{\pi_0}=\pi$, notamment les cardinaux $\pi'^{\pi_0}$, avec $\pi'\geqslant 2$, ont la même propriété), ou ce
\msp{page 206}
qui revient au même d'après 4.8.1, si $\mathcal{M}$ admet une fcga, ou encore une fcg canonique.

\textbf{Commentaire 4.8.8.} Cette définition est en contradiction avec la terminologie introduite dans SGA 4 I 9.4 (p. 114) ; on désignera dorénavant sous les noms de ``catégorie faiblement $\pi_0$-accessible'', ``catégorie faiblement accessible'' les notions introduites dans 9.4 -- i.e. les catégories admettant une sous-catégorie pleine génératrice (pas même strictement génératrice) dont les objets sont $\pi_0$-accessibles (resp. accessibles).

Notons que si $\mathcal{M}$ est $\pi_0$-accessible (resp. accessible), la condition de 4.8.1 nous dit que $\mathcal{M}$ est faiblement $\pi_0$-accessible (resp. faiblement accessible), et même qu'il existe une petite sous-catégorie pleine \textit{strictement} génératrice, qui soit formée d'objets $\pi_0$-accessibles (resp. accessibles). Je doute fort que cette dernière condition nécessaire soit aussi suffisante pour
\printedpage{73}
% head: A. Grothendieck : Dérivateurs, Ch. XVIII, Édition des Universités de Montpellier II et Paris VII
que $\mathcal{M}$ soit $\pi_0$-accessible (resp. accessible). Mais en vertu de 4.6.3.3 (p. 150), il en est ainsi si on suppose $C$ stable par petites $\varinjlim$ , ou ce qui revient au même, moyennant $\mathrm{L}_{\pi_0}$, stable par limites inductives de cardinal $\leqslant\pi_0$.

\msp{page 207}
\textbf{Proposition 4.8.9.} \textit{Soit $\mathcal{M}$ une $\mathfrak{U}$-catégorie accessible (4.8.7), $\pi_0$ un cardinal tel que $\mathcal{M}$ satisfait à $L_{\pi_0}$. Conditions équivalentes :}
\begin{enumerate}
\item[\textit{a)}] \textit{$\mathcal{M}$ est une catégorie paratopos (cf. 4.6.3, $3^\circ$, p. 150).}
\item[\textit{b)}] \textit{$\mathcal{M}$ est stable par petites limites inductives.}
\item[\textit{b$'$)}] \textit{$\mathcal{M}$ est stable par limites inductives de cardinal $\leqslant\pi_0$.}
\item[\textit{b$''$)}] \textit{(Si $\mathcal{M}$ est équivalente à une catégories} $\Ind_{\pi_0}C$\textit{, avec $C$ une petite catégorie} karoubienne.\textit{) La catégorie $C$ est stable (en elle-même) par limites inductives de cardinal $\leqslant\pi_0$.} % [sic] une catégories -> une catégorie, accord du nombre
\item[\textit{c)}] \textit{$\mathcal{M}$ est stable par petites limites projectives.}
\item[\textit{c$'$)}] \textit{(Sous l'hypothèse $\mathcal{M}\simeq\Ind_{\pi_0}(C)$, et quand on se donne un cardinal $\pi$ tel que $\pi=\pi^{\pi_0}\geqslant\pi_1=\sup(\pi_0,\card\textup{réd}\,C)$.) $\mathcal{M}_\pi$ est stable par $\varprojlim$ finies, et $\mathcal{M}$ est stable par petits produits infinis (ou encore, par $\varprojlim$ filtrantes).}
\item[\textit{c$''$)}] \textit{Il existe un cardinal infini $\pi$, et une petite catégorie $\mathcal{D}$ stable par $\varprojlim$ finies, telle que $\mathcal{M}\simeq\Ind_\pi(\mathcal{D})$, de plus $\mathcal{M}$ stable par produits infinis (ou encore par $\varprojlim$ filtrantes).}
\end{enumerate}
(\footnote{\textbf{N.B.} L'équivalence de c$'$, c$''$ avec les autres conditions n'est pas établie (ici, mais seulement dans 4.11 plus bas.)}, \footnote{\textbf{N.B.} Les critères c$'$, c$''$ me paraissent plus techniques et moins utiles \texttt{[c$'$, c$''$ ci-dessus sont mis en parenthèses]}.}).

\dgrblock{dgrb1030}{Accessibility and the chosen cardinal}
\begin{dcomment}{status}{p206}{Accessibility and the chosen cardinal}
Grothendieck's change of terminology brings the chapter to the modern structural notion: an accessible category is generated by a small subcategory under sufficiently filtered colimits, with the generators presentable at the same cardinal. His earlier weak notion only asked for a small generating family of accessible objects. The cardinal is part of the assertion. In particular, the existence of a dense family of $\mu$-presentable objects and of $\mu$-filtered colimits does not by itself provide the filtered presentation required in the definition.

Questions 4.8.6(1) and the nonuseful-cardinal part of (4) have a common concrete answer. Let $\mathcal K$ be the category of nonempty $\aleph_1$-directed posets and monotone maps. It is $\aleph_1$-accessible: each object is the $\aleph_1$-directed union of its countable subposets with a greatest element. Its $\lambda$-presentable objects, for regular $\lambda\geq\aleph_1$, are precisely the posets of cardinal less than $\lambda$. Our commentary on \S4.7 proves these facts and the resulting obstruction, following \citet[Example~3.8]{Vasey2020Accessible}. For every singular $\theta>\aleph_1$ of countable cofinality, the poset $[\theta]^{<\aleph_1}$ has no cofinal subset of cardinal at most $\theta$. If it were a $\theta^+$-filtered colimit of $\theta^+$-presentables, a single stage image would contain all $\theta$ singleton subsets. That image would be $\aleph_1$-directed and of cardinal at most $\theta$, hence would be a forbidden cofinal subset. Thus $\mathcal K$ is not $\theta^+$-accessible, although its original small dense family of countable presentables remains $\theta^+$-presentable. It satisfies 4.8.1(a), but not (b), at manuscript cardinal $\theta$. Taking arbitrarily large such $\theta$ answers the unbounded version of (4) for this one category.

The modern change-of-cardinal theorem explains the allowed upward steps through a sharp-inequality hypothesis; see \citet[Proposition~5.4.2.9 and Remark~5.4.2.10]{Lurie2009HTT}. Grothendieck's sufficient equation $\pi^{\pi_0}=\pi$ is a concrete way to ensure the requisite closure of small subsets under the previous stage's operations. The example shows why the cardinal condition has mathematical content.

For question~(2), cocompleteness is a sufficient condition already established in 4.6.3. More economically, adjoining the existence of all colimits with at most $\pi_0$ arrows to $L_{\pi_0}$ supplies cocompleteness by the earlier decomposition of a small diagram into bounded pieces. The closure of the chosen generators under those colimits then makes the overcategories filtered at the required cardinal. This is a repair by adding operations that create cones, rather than merely adding more generators.

The useful-but-nonregular alternative in question~(4) is excluded by the repaired presentation argument of \S4.7. If $\pi$ is useful in the sense of the comparison equivalence, write $\mathcal M=\Ind_\pi(\mathcal M_\pi)$. For $\pi'>\pi$ put $\mu=(\pi')^\pi$, so that $\mu^\pi=\mu$. Given $F\in\mathcal M_{\pi'}\subseteq\mathcal M_\mu$, the fixed-point presentation theorem, with base category $\mathcal M_\pi$, supplies a $\pi^+$-filtered presentation of $F$ of size at most $\mu$. The bounded cofinal-poset construction in our commentary on \S4.7 makes this index ordered without increasing its size. This is precisely the quantitative regularity required in 4.7.13, without an additional threshold $\pi_1$. The nonuseful cardinals supplied above therefore exhaust the possible kind of failure asked for in question~(4).
\end{dcomment}


\textsc{Démonstration.} Les implications suivantes sont tautologiques, ou déjà $\pm$ connues :
\[
\begin{tikzcd}[column sep=large, row sep=large]
\mathrm{a} \arrow[r, Rightarrow] & \mathrm{b} \arrow[r, Rightarrow, "\text{cf. 4.4.14}"] \arrow[dl, Rightarrow] & \mathrm{b}' \arrow[r, Rightarrow, "\bullet"] & \mathrm{b}'' \\
\mathrm{c}\ ( \arrow[r, Rightarrow, "\bullet"] & \mathrm{c}' \arrow[r, Rightarrow] & \mathrm{c}'')\ , &
\end{tikzcd}
\]
\msp{page 208}
où seules les implications pointées b$'\Longrightarrow$ b$''$ et c $\Longrightarrow$ c$'$ demandent une démonstration. Pour b$'\Longrightarrow$ b$''$, on note que l'hypothèse $C$ karoubienne implique que $C\simeq\mathcal{M}_{\pi_0}$ (4.8.5, $2^\circ$), or on sait que $\mathcal{M}_{\pi_0}$ est stable dans $\mathcal{M}$ par limites inductives de cardinal $\leqslant\pi_0$ (et celles-ci sont représentables dans $\mathcal{M}$ en vertu de b$'$), donc $\mathcal{M}_{\pi_0}$ admet des $\varinjlim$ du type envisagé, donc $C$ aussi.

c $\Longrightarrow$ c$'$. Notons que si $\mathcal{M}=\Ind_{\pi_0}(C)$, alors les $\varprojlim$ représentables dans $\mathcal{M}$ se calculent dans $C^\wedge$. Il s'ensuit que si $\pi\geqslant\pi_1$, alors $\Filt^\pi$ est stable par $\varprojlim$ finies ($\Filt^\pi$ étant donné alors par la condition a$''$ de 4.7.1), d'autre part, si de plus $\pi^{\pi_0}=\pi$, $\pi$ est un cardinal utile, donc $\mathcal{M}\simeq\Ind_\pi(\mathcal{D})$, où $\mathcal{D}=\Filt^\pi=\mathcal{M}_\pi$. Il suffit maintenant de voir que $\Ind_\pi(\mathcal{D})$ est alors stable par $\varprojlim$ finies, pour conclure qu'il en est de même de $\mathcal{M}$, donc que $\mathcal{M}$ est stable par limites projectives quelconques, s'il l'est par petites produites ou par petites $\varprojlim$ filtrantes, d'où c$'$. Il reste donc à prouver le lemme suivant. % [sic] petites produites -> petits produits, accord

% ---- from der-ch18_p076-080.tex
\printedpage{74}
% head: A. Grothendieck : Dérivateurs, Ch. XVIII, Édition des Universités de Montpellier II et Paris VII

\textbf{Lemme 4.8.10.} \textit{Soient $C$ une petite catégorie, $\pi_0$ un cardinal infini, $\pi_1 = \sup(\pi_0, \card\red C)$. Supposons $\Ind_{\pi_0}$ stable par $\varprojlim$ finies (resp. et d'un certain type). Alors pour tout $\pi \geqslant \pi_1$, $\Filt^{\pi}$ est stable par lesdites $\varprojlim$ . (Donc aussi $\cM_{\pi} = \Filt^{\pi}$, si $\pi = \pi^{\pi_0}$.) Plus} % [sic] \Ind_{\pi_0} stable -> \Ind_{\pi_0}C stable, the argument C appears to be omitted (cf. the following sentence)
\msp{page 209}
\textit{généralement, soit $J$ une petite catégorie telle que $\Ind_{\pi_0}C$ soit stable par $\varprojlim$ de type $J$, et soit $c = \card\red J$, supposons que $\pi^{c} = \pi$ (p.ex. $c$ fini). Alors $\Filt^{\pi}$ est stable dans $\cM$ par $\varprojlim$ de type $J$. (Donc il en est de même de $\cM_{\pi} = \Filt^{\pi}$, si de plus $\pi^{\pi_0} = \pi$.)}

On utilise le critère a$''$ de 4.7.1, savoir $\card F(x) \leqslant \pi$ pour tout $x \in C$, pour l'appartenance à $\Filt^{\pi}$. On peut supposer $J$ réduite, donc $\card\Ob J \leqslant c$. Comme le foncteur d'inclusion $\Ind_{\pi_0}(C) \longrightarrow C^{\wedge}$ commute aux $\varprojlim$ quelconques, le calcul de $\varprojlim\limits_{j\in J} x_j$ dans $C^{\wedge}$ fibre par fibre montre que dans $C^{\wedge}$, on a
\[
X \ \defeq\ \varinjlim_{j\in J} x_j \ = \ \underbrace{\prod_{j\in\Ob J} x_j}_{\text{produit dans } C^{\wedge}} \quad ,
\] % [sic] \varinjlim_{j\in J} x_j -> \varprojlim_{j\in J} x_j, the text concerns the projective limit (cf. the preceding sentence and (4.8.10.1))
d'où
\[
\card X(x) \ \leqslant\ \prod_{j\in\Ob J} \underbrace{\card x_j(x)}_{\leqslant\pi} \ \leqslant\ \pi^{c}\ ,
\]
d'où (sans autre hypothèse sur $\pi$ que $\pi \geqslant \pi_1$)
\[
\varprojlim_{j\in J} x_j \ \in\ \Ob\Filt^{\pi^{c}}\ , \tag*{(4.8.10.1)}
\]
d'où la conclusion voulue si $\pi^{c} = \pi$, q.e.d.

\msp{page 210}

\dgrblock{dgrb1031}{Limits are bounded by products}
\begin{dcomment}{exposition}{p209}{Limits are bounded by products}
Grothendieck's cardinal estimate makes finite limits manageable at sufficiently large stages of the filtration. It uses the restricted Yoneda embedding, not a special property of the object being limited. An existing limit in $\mathcal M$ is carried to the pointwise limit of its restricted representable functors, because each $\Hom(c,-)$ preserves limits. If $X=\lim_{j\in J}X_j$, then
\[
 X(c)\hookrightarrow\prod_{j\in\Ob J}X_j(c),
 \qquad |X(c)|\leq\pi^{|\Ob J|}
\]
whenever $|X_j(c)|\leq\pi$. The source display has two slips: an inductive limit in place of a projective one, and an equality with the product in place of an inclusion. A limit imposes the compatibility equations associated with the arrows of $J$; deleting those equations generally enlarges it. The next displayed inequality is correct and is all the proof needs.

For example, let $J$ be the parallel-pair category and take two functions $f,g:A\rightrightarrows B$. Its limit is the equalizer $\{a\in A:f(a)=g(a)\}$, with the redundant $B$-coordinate determined by $a$. For $A=B=\{0,1\}$, $f=1$, and $g$ the transposition, that limit is empty, while the product over the two objects of $J$ has four elements. Thus the equality cannot be read literally even in sets. For a finite nonempty shape and infinite $\pi$, the estimate becomes $\pi^{|\Ob J|}=\pi$; the empty shape is the terminal presheaf and has singleton values. The latter should be treated separately when the notation $\Filt^{\pi^0}=\Filt^1$ would fall outside the filtration's declared range. This proves the intended finite-limit bound without changing any substantive conclusion of 4.8.10.
\end{dcomment}


Nous allons prouver les implications supplémentaires
\[
\text{b}'' \ \Longrightarrow\ \text{b}\ , \qquad \text{c} \ \Longrightarrow\ \text{b}\ , \qquad \text{c}'' \ \overset{?}{\Longrightarrow}\ \text{c}\ ,
\]
ce qui achève la démonstration.

L'implication b$''$ $\Longrightarrow$ a est contenue dans 4.6.3.

Pour l'implication c $\Longrightarrow$ b, on note que la stabilité de $\cM$ par petites $\varinjlim$ équivaut à l'assertion que toute petite limite projective de foncteurs $\cM \xrightarrow{\ \xi_i\ } \Ens$ représentables, est représentable. Or les objets de $\cM$ étant accessibles ($\cM$ étant accessible), les foncteurs représentables $\xi_i$ sont accessibles, donc aussi leur $\varprojlim$ (SGA 4 I 9.7). D'autre part, le foncteur $\xi$ est exact à gauche, comme $\varprojlim$ de foncteurs représentables donc exactes à gauche. Donc le fait que $\xi$ est représentable résulte du % [sic] exactes -> exacts, agreement with « foncteurs »

\textbf{Lemme 4.8.11.} \textit{Soit $\cM$ une catégorie accessible, stable par $\varprojlim$ finies, $\cM \longrightarrow$} Ens \textit{un foncteur. Pour que $\xi$ soit représentable, il faut et il suffit qu'il soit exact à gauche et accessible. Donc (SGA 4 I 8.12.8.1) si $\cM$ est stable par petites $\varprojlim$ , $\xi$ est représentable si et seulement s'il commute auxdites $\varprojlim$ , et qu'il soit accessible.}

% foot: Version 17 décembre 2025

\printedpage{75}
% head: A. Grothendieck : Dérivateurs, Ch. XVIII, Édition des Universités de Montpellier II et Paris VII

\msp{page 211}

\dgrblock{dgrb1032}{Prorepresentability is the missing word}
\begin{dcomment}{examples}{p210}{Prorepresentability is the missing word}
Grothendieck is trying to recover colimits from completeness by representing a limit of covariant representable functors. That strategy is sound, but the first sentence of Lemma~4.8.11, as printed, is false: accessibility and preservation of finite limits imply \emph{prorepresentability}, not representability. This is exactly the conclusion of the cited SGA~4, I, 8.13.2(b). The following sentence, which adds all small limits, is the representability statement; SGA~4, I, 8.12.8.1 supplies the additional step \citep{SGA4I1972}.

An explicit example keeps the two conclusions apart. Put $A_n=\{0,1\}^n$, with the truncation surjections $A_{n+1}\to A_n$, and define a covariant functor on sets by
\[
 T(X)=\varinjlim_{n<\omega} X^{A_n}.
\]
The transition maps repeat a function on the two extensions of each binary word. Thus $T(X)$ is the set of functions from $\{0,1\}^{\mathbb N}$ to $X$ which depend on finitely many coordinates. It is equally the functor of continuous maps from the Cantor space to the discrete space $X$; compactness ensures that such a map factors through some $A_n$.

Each functor $X^{A_n}$ preserves all limits and filtered colimits, since $A_n$ is finite. Filtered colimits commute with finite limits in sets. It follows, by interchanging two colimits for the accessibility assertion, that $T$ is finitary and preserves finite limits. However, the comparison
\[
 T(\{0,1\}^{\mathbb N})\longrightarrow
 \prod_{k<\omega}T(\{0,1\})
\]
does not contain the family of coordinate functions in its image. A preimage would have to be the identity function on $\{0,1\}^{\mathbb N}$, which depends on every coordinate. Therefore $T$ fails to preserve this countable product and cannot be representable. Its defining inverse system $(A_n)$ is its prorepresentation. This is a fully explicit counterexample to the printed first sentence, with $\mathcal M=\mathrm{Set}$ satisfying even stronger hypotheses than those stated.

The application $c\Rightarrow b$ in 4.8.9 survives. If $D:K\to\mathcal M$ is small and $\mathcal M$ is accessible and complete, the functor
\[
 H(X)=\lim_{k\in K^{\mathrm{op}}}\Hom(Dk,X)
\]
preserves all small limits. Choose a sufficiently large useful cardinal so that all $Dk$ are presentable and $K$ has smaller cardinality. Interchange of its defining small limit with sufficiently filtered colimits makes $H$ accessible. The solution-set form of the adjoint functor theorem now represents $H$, and its representing object is $\varinjlim_KD$.

For completeness, the required solution set consists of all pairs $(A,a)$ with $A$ in a small presentable generating subcategory at a suitable cardinal and $a\in H(A)$. Every $(X,x)$ receives a map from such a pair because $H$ preserves the generating filtered colimits. The category of elements of $H$ is complete, since $H$ preserves all limits. A complete locally small category with a small weakly initial family has an initial object: take the product of the family to obtain a weakly initial object $P$, then equalize all endomorphisms of $P$. If $e:E\to P$ is that equalizer, choose $r:P\to E$ by weak initiality. Since $ere=e$, monicity gives $re=1_E$. For any $h:E\to E$, the endomorphism $ehr$ of $P$ satisfies $ehre=e$, hence $h=1_E$. Equalizers of two maps out of $E$, together with weak initiality of $E$, then show uniqueness of every such map. Thus $E$ is initial. This proves the representability step with the hypothesis actually needed.
\end{dcomment}


Le ``il faut'' est remis pour mémoire, on vient de le voir comme conséquence du fait que tout objet de $\cM$ est accessible. Le ``il suffit'' est essentiellement SGA 4 I 8.13.2.b), sauf que dans loc. cit. on suppose $\cM$ munie d'une fc \textit{stricte.} Mais la démonstration donnée dans loc. cit. montre aussi pour une fcga, en choisissant dans loc. cit. (p. 137) $\pi$ tel que $\xi$ soit $\pi$-accessible, et de plus $\pi$ soit utile pour la filtration cardinale généralisée envisagée.

Reste à prouver c$''$ $\overset{?}{\Longrightarrow}$ c, sous réserve que ce soit bien vrai. Ce serait une conséquence du lemme suivant (où $\cD$, $\pi$ sont remplacés par $C$, $\pi_0$).

\textbf{Lemme 4.8.12} (?) (\footnote[88]{Lemme non prouvé ici -- cf. démonstration 4.11.7 (non achevée).}). \textit{Soient $C_0$ une petite catégorie, $\pi_0$ un cardinal infini. Si $C$ est stable par $\varprojlim$ finies, alors $\Ind_{\pi_0}C$ également.} % [sic] C_0 -> C, the statement then speaks only of C

L'assertion similaire pour $\Ind(C)$, au lieu de $\Ind_{\pi}C$, est établie dans SGA 4 I 8.9.1 c). Mais il n'est pas clair sans plus que $\Ind_{\pi}C$ soit stable dans $\Ind(C)$ par $\varprojlim$ finies, dans le cas envisagé. De façon précise, si $F = \varprojlim\limits_{j\in J} F_j$, les $F_j$ dans $\Ind_{\pi}C$,
\msp{page 212}
$F$ dans $\Ind(C) \subset C^{\wedge}$, on sait que $C/F$ est filtrant et il faut voir qu'il est filtrante grande devant $\pi_0$. On est ramené aux trois cas où $J = \emptyset$, $J = e \amalg e$, et $J = \vcenter{\hbox{\begin{tikzcd}[ampersand replacement=\&, column sep=small, row sep=small]
\bullet \arrow[dr] \& \& \bullet \arrow[dl] \\
\& \bullet \&
\end{tikzcd}}}$. Les deux premiers cas sont essentiellement triviaux (car $e$, et un produit de deux catégories grandes devant $\pi_0$, sont grandes devant $\pi_0$), le troisième cas par contre n'a pas l'air trivial a priori. C'est lié visiblement à la question des représentations indicielles de foncteurs $J \longrightarrow \Ind_{\pi_0}(C)$, pour $\card\red J \leqslant \pi_0$ -- i.e. la validité de 4.4.13 b) en dehors de l'hypothèse que $C$ soit stable par $\varinjlim$ de cardinal $\leqslant \pi_0$. Il nous faudra y revenir très prochainement, car cette question m'apparaît à présent une question-clef, pour les propriétés de stabilité de la notion de catégorie accessible (en dehors d'aucune hypothèse d'existence de $\varinjlim$ ou $\varprojlim$ ). % [sic] qu'il est filtrante -> qu'il est filtrant, agreement with « il » (C/F)

% marginal handwritten mark (curved stroke) in the left margin beside Proposition 4.8.13
\textbf{Proposition 4.8.13} (\footnote[89]{\textit{Pas prouvé.} Mais OK si $\pi_0 > \card\red C$, cf. 4.15.10, p. 462.}). \textit{Soit $C$ une catégorie équivalente à une petite catégorie, et soit $\pi_0 \geqslant \card\red C$ un cardinal infini. Alors $\cM = \Ind_{\pi_0}C = \cM_{\pi_0}$ (i.e. tous les objets de $\Ind_{\pi_0}C$ sont $\pi_0$-accessibles), et $\cM$ est une enveloppe de Karoubi de $C$.}

\msp{page 213}

\dgrblock{dgrb1033}{Finite limits of a bounded ind-completion}
\begin{dcomment}{exposition}{p212}{Finite limits of a bounded ind-completion}
Grothendieck isolates the real issue in Lemma~4.8.12: the pointwise limit must still admit a presentation with the stronger filteredness condition. The statement is true. Its ordinary ind-version is SGA~4, I, 8.9.1(c), and the later manuscript returns to it in 4.11.7 \citep{SGA4I1972,Grothendieck1990XVIII}. We give the bounded argument, since it needs only a span and does not assume an indexed representation theorem for arbitrary finite categories with endomorphisms.

Suppose $C$ has finite limits, and put $\lambda=\pi_0^+$. The terminal presheaf is represented by the terminal object of $C$. For two $\lambda$-filtered presentations, the product indexing category is $\lambda$-filtered and filtered colimits commute with finite products in sets; hence their product is a $\lambda$-filtered colimit of the representing products. It remains to treat a pullback $F\times_HG$.

Choose small $\lambda$-filtered presentations $F=\varinjlim_i y(a_i)$, $G=\varinjlim_j y(b_j)$, and $H=\varinjlim_k y(c_k)$. Let $L$ have objects $(i,j,k,u,v)$, where $u:a_i\to c_k$ and $v:b_j\to c_k$ induce the given maps to $H$. Its arrows are triples of index arrows making the two squares commute. The category $L$ is $\lambda$-filtered. To see this, begin with a diagram in $L$ having fewer than $\lambda$ arrows. Choose cones on its $i$-, $j$-, and $k$-components. The maps from the resulting $a_i,b_j$ to $H$ factor at some stages of the $c$-diagram, since representables are $\lambda$-presentable. Move them and the previously chosen $k$-cone to a common stage. Each desired compatibility equation holds in $H$, so it holds after a further transition. There are fewer than $\lambda$ such equations; $\lambda$-filteredness supplies one further stage satisfying all of them. The resulting object of $L$ is the required cone. The same argument with an empty starting diagram gives nonemptiness.

There is a canonical comparison
\[
 \varinjlim_{(i,j,k,u,v)\in L}y(a_i\times_{c_k}b_j)
 \longrightarrow F\times_HG.
\]
Evaluate at $d\in C$. A pair of elements on the right is represented by maps $d\to a_i$, $d\to b_j$ whose images agree in $H(d)$. Factor the two maps to $H$ through a common $c_k$ and enlarge $k$ until their composites from $d$ agree; the pair then factors through $a_i\times_{c_k}b_j$. This proves surjectivity. If two representatives have the same image, first move their $i$- and $j$-indices forward until the two coordinate maps agree; then choose a common $k$ and impose the finitely many compatibility equations just used in the filteredness proof. The two representatives become equal in the left-hand colimit. This proves injectivity, naturally in $d$.

The displayed comparison is therefore an isomorphism. Its source is a $\lambda$-filtered colimit of representables, which proves the lemma. Together with arbitrary products it supplies all small limits. If instead all small cofiltered limits are assumed, an arbitrary product is the inverse limit of its finite subproducts, indexed by the opposite of the directed poset of finite subsets. Thus the alternative formulation using filtered projective limits in 4.8.9 is justified as well.
\end{dcomment}


Signalons tout de suite les corollaires suivants.

% marginal handwritten mark (curved stroke) in the left margin beside Corollaire 4.8.14
\textbf{Corollaire 4.8.14} (\footnote[90]{\textit{Pas prouvé.} Si finalement, cf. 4.15.10, et 4.15.11, pages 462-463.}). \textit{Soit $\cM$ une $\mathfrak{U}$-catégorie. Conditions équivalentes :}
\begin{enumerate}[leftmargin=*]
\item[\textit{a)}] \textit{$\cM$ est accessible, et est équivalent à une petite catégorie.}
\item[\textit{a$'$)}] \textit{$\cM$ est accessible, et il existe un cardinal infini $\pi_0$ tel que tous les objets de $\cM$ soient $\pi_0$-accessibles (i.e. $\cM = \cM_{\pi_0}$).}
\item[\textit{b)}] \textit{$\cM$ est équivalente à une petite catégorie, et est karoubienne, i.e. stable par facteurs directs} (\footnote[91]{\textbf{N.B.} Ceci implique que $\cM$ et $\cM^{\circ}$ sont accessible, et je conjecture que la réciproque est vraie également.}). % [sic] sont accessible -> sont accessibles, plural agreement (footnote 91)
\end{enumerate}

L'équivalence de a) et a$'$) est claire, puisque $\cM$ admet une filtration cardinale généralisée canonique $(\cM_{\pi})_{\pi\geqslant\pi_0}$, formée des catégories équivalentes à des petites catégories (condition
% foot: Version 17 décembre 2025
\printedpage{76}
% head: A. Grothendieck : Dérivateurs, Ch. XVIII, Édition des Universités de Montpellier II et Paris VII
a) des fcg), et telle que toute telle sous-catégorie soit contenue dans un des $\cM_{\pi}$ (résulte de la condition b) des fcg). Comme $\cM_{\pi_0}$ est équivalente à une petite catégorie, et est karoubienne (étant stable par $\varinjlim$ filtrantes grandes devant $\pi_0$), il est clair aussi que a$'$) $\Longrightarrow$ b). Reste à prouver b) $\Longrightarrow$ a$'$), or cela résulte aussitôt de 4.8.13 appliqué à $\cM$, qui montre de plus qu'on peut prendre pour $\pi_0$ n'importe quel cardinal infini tel que $\pi_0 \geqslant \card\red\cM$.

\msp{page 214}

\dgrblock{dgrb1034}{An accessible category can have an accessible opposite}
\begin{dcomment}{examples}{p213}{An accessible category can have an accessible opposite}
Grothendieck's proposed converse in footnote~91 would turn accessibility on both sides into a smallness criterion. Partial injections show why the existence of many filtered colimits is insufficient for this purpose. Let $\mathcal P$ have small sets as objects and partial injective functions as arrows. Thus an arrow $X\dashrightarrow Y$ consists of a subset $D\subseteq X$ and an injection $D\hookrightarrow Y$. Taking the inverse partial function defines an isomorphism $\mathcal P^{\mathrm{op}}\cong\mathcal P$. This category is finitely accessible; precise published treatments of the example include \citet[\S3.2, Proposition~3.6, Lemma~3.8, Theorem~3.9]{Heunen2013L2}. We give a direct argument, also keeping track of the larger cardinals used in this chapter.

Adjoining a basepoint identifies $\mathcal P$ with the category of pointed sets and pointed functions whose fibre over each nonbasepoint has at most one element. A partial function sends points outside its domain to the basepoint. Let $(X_i)_{i\in I}$ be a filtered diagram in $\mathcal P$, and form its colimit $L_*$ in pointed sets. The canonical map $(X_i)_*\to L_*$ has at most one element over every nonbasepoint. Indeed, if two points of $X_i$ have the same nonzero image, filteredness gives a stage at which they have the same nonzero image; the transition map is a partial injection, so the points were equal. Consequently the canonical maps define partial injections $X_i\dashrightarrow L$, where $L=L_*\setminus\{*\}$.

The same argument proves the universal property within $\mathcal P$. A cocone $X_i\dashrightarrow Y$ gives a unique pointed function $L_*\to Y_*$. If two nonzero elements of $L_*$ have the same nonzero image, represent them at a common stage. The cocone map at that stage is injective wherever it is defined; hence the representatives, and therefore the two colimit elements, coincide. The induced map is a partial injection. Thus filtered colimits exist in $\mathcal P$ and are created by this embedding into pointed sets.

Now let $\lambda$ be an infinite regular cardinal. A set $A$ of cardinal $<\lambda$ is $\lambda$-presentable in $\mathcal P$. To factor a partial injection $A\dashrightarrow L$ through a $\lambda$-filtered diagram, choose a representative for each of its fewer than $\lambda$ defined values and move them to one stage. Distinct colimit values cannot have coincident representatives, so this gives a partial injection with exactly the prescribed domain. If two such factorizations have the same composite to $L$, the equality of their values, including any values that become the basepoint, is witnessed at later stages. There are fewer than $\lambda$ witnesses; $\lambda$-filteredness puts them at one stage. This proves both surjectivity and injectivity of the presentability comparison map.

Conversely, $X$ is the colimit of its subsets of cardinal $<\lambda$, ordered by inclusion. This poset is $\lambda$-filtered by regularity. If $X$ were $\lambda$-presentable, its identity would factor through one of these inclusions $B\hookrightarrow X$. The composite being the total identity forces $X=B$, so $|X|<\lambda$. We have therefore proved the exact description
\[
 \mathcal P_\lambda^{\mathrm{pres}}
   =\{X:|X|<\lambda\},\qquad
 \mathcal P\simeq
 \operatorname{Ind}_{\lambda}\bigl(\mathcal P_\lambda^{\mathrm{pres}}\bigr),
\]
where $\operatorname{Ind}_{\lambda}$ in this display has its modern, regular-cardinal meaning. In the notation of the manuscript take $\lambda=\pi^+$: the $\pi$-accessible sets in $\mathcal P$ are exactly those of cardinal at most $\pi$. A skeleton of this subcategory is small, including its partial injections.

An isomorphism in $\mathcal P$ is an ordinary bijection: an inverse under partial composition forces both domains to be the whole sets. Thus $\mathcal P$ has objects of every small cardinality and is not essentially small relative to the chosen universe. Both it and its opposite are accessible. The conjectured converse is false; the equivalence in Corollary~4.8.14 between small accessible categories and small idempotent-complete categories is unaffected. The obstruction is visible without any higher category theory: a coproduct of two singleton sets does not exist in $\mathcal P$. The cocones in which exactly one singleton map is defined force both structural arrows to be everywhere defined and their images to be distinct. The cocone sending both singleton points to the same singleton then cannot factor through a partial injection. Completeness hypotheses cannot be dropped from the familiar duality restrictions for locally presentable categories.
\end{dcomment}


\textbf{4.8.15.} \textsc{Démonstration de} 4.8.13. Nous allons prouver :

\textbf{Lemme 4.8.15.1.} \textit{Sous les conditions de 4.8.13, soit $F \in \Ob\cM$, $x \in \Ob C$, alors}
\[
\card F(x) \ \leqslant\ \pi_0\ . \tag*{$(*)$}
\]

Ceci signifie donc, en vertu de 4.7.7, a$''$ $\Longrightarrow$ b pour $\pi \geqslant \pi_1$, compte tenu qu'ici $\pi_J = \pi_0$, $\pi = \pi_0$, qu'on a $F \in \Filt^{\pi_0}$, d'où a fortiori $F \in \cM_{\pi_0}$, ce qui implique $\cM = \cM_{\pi_0}$, i.e. la proposition 4.8.13, compte tenu de 4.8.5. Donc tout revient à prouver $(*)$.

Écrivons
\[
F \ = \ \varinjlim_{i\in I} x_i\ , \qquad\quad \text{les } x_i \in \Ob C,\ I \text{ ensemble ordonné grand devant } \pi_0,
\]
d'où
\[
F(x) \ = \ \varinjlim_{i\in I} \Hom_C(x, x_i)\ .
\]
On peut supposer
\[
\card C \ = \ \pi_0\ ,
\]
quitte à remplacer $C$ par une catégorie réduite équivalente.

Pour tout $\xi \in \Ob C$, soit
\[
I_{\xi} \ = \ \{ i \in I \mid x_i = \xi \}\ ,
\]
alors
\[
I \ = \ \bigcup_{\xi\in\Ob C} I_{\xi}\ ,
\]
\msp{page 215}
réunion d'une famille de parties de $I$ de cardinal $\card\Ob C \leqslant \pi_0$. Comme $I$ est grand devant $\pi_0$, on voit que cela implique que l'ensemble des $I_{\xi}$ est cofinal dans $I$. Quitte à remplacer $I$ par $I_{\xi}$, on peut supposer $I_{\xi} = I$, i.e. que tous les $x_i$ sont égaux à un même élément $\xi$. Quitte à remplacer $C$ par la sous-catégorie pleine $C_0 = \{\xi\}$, on peut supposer $C$ réduite à un seul élément. (Car si on prouve que $F$, en tant qu'objet de $\Ind_{\pi_0}(C_0)$, est facteur direct d'un objet de $C_0$ (i.e. de $\xi$), il en sera de même de $F$ en tant qu'objet de $\Ind C$). Donc on peut supposer
\[
C \ = \ \{\xi\}\ ,
\]
donc $C$ donné par le monoïde
\[
M \ = \ \End_C(\xi)\ , \qquad\quad \card M \ \leqslant\ \pi\ ,
\]
et la donnée du système inductif $(x_i)_{i\in I}$ revient à la donnée d'un système d'éléments
\[
\alpha_{ji} \ \in\ M \qquad\quad (i, j \in I,\ i \leqslant j)
\]
% foot: Version 17 décembre 2025
\printedpage{77}
% head: A. Grothendieck : Dérivateurs, Ch. XVIII, Édition des Universités de Montpellier II et Paris VII
satisfaisant
\[
\alpha_{ii} \ = \ 1\ , \qquad\quad \alpha_{kj}\alpha_{ji} \ = \ \alpha_{ki} \ \ \text{si } i \leqslant j \leqslant k.
\]
On a
\[
F(\xi) \ = \ \varinjlim_{i\in I} M_i \qquad\quad
\left\{
\begin{array}{l}
\text{où } M_i = M \ \forall\, i \in I, \text{ et } M_i \longrightarrow M_j \\
\text{est donné par} \\[2ex]
\qquad\qquad u \longmapsto \alpha_{ji}u\ . \\[4ex]
\end{array}
\right.
\]

\msp{page 216}

J'ai pu me ramener, par des arguments similaires (\footnote[92]{Pour ces réductions, cf. 4.15.8 (p. 453-463). L'ultime réduction se trouve dans 4.15.8.9 (p. 461), et son corollaire 4.15.10.}), au cas où le monoïde $M = \End(\xi)$ est un monoïde de monomorphismes, i.e.
\[
uv \ = \ uv' \ \ \Longrightarrow\ \ v \ = \ v' \qquad\quad \text{si } u, v, v' \in \cM\ ,
\] % [sic] u, v, v' \in \cM -> u, v, v' \in M, the elements belong to the monoid M
ce qui implique que les flèches de transition dans le système inductif des $M_i = M$ sont des monomorphismes. Il s'ensuit que l'on ne peut avoir $\card F(\xi) \leqslant \pi_0$ ($I$ étant grand devant $\pi_0$) que si $\alpha_{ji_0}$ est un isomorphisme pour $j \geqslant i_0$, et $i_0$ assez grand, i.e. si le système inductif induit pour $j \geqslant i_0$ est constant, a fortiori que le système inductif envisagé est essentiellement constant. Mais je viens de passer près de quatre heures pour essayer de prouver ce point, sans y être arrivé ! Cela semble plus délicat à prouver que je ne le prévoyais, mais je crois toujours que c'est bel et bien vrai.

\msp{page 217}

\dgrblock{dgrb1035}{The small-category theorem and its boundary}
\begin{dcomment}{status}{p216}{The small-category theorem and its boundary}
Grothendieck's intended conclusion is that enough filteredness forces a diagram in a small category to become constant up to splitting an idempotent. The reduction to a one-object category is valid. If a $\pi_0^+$-directed poset is the union of at most $\pi_0$ subsets, at least one subset is cofinal: otherwise choose one obstruction to cofinality for each subset and take a common upper bound, contradicting the union. The same subset, with its inherited order, remains $\pi_0^+$-directed. Applying this to the fibres of the object function gives a cofinal presentation whose values are one fixed object $c$; its remaining data lie in $\End_C(c)$.

The cancellation reduction and eventual-stationarity assertion on page~216 are not proved in these pages. Nor does the later annotation establish the exact endpoint printed in 4.8.13. The official continuation, 4.15.10--4.15.11, proves the conclusion under the strict inequality $\pi_0>\cardred C$ and explicitly describes it as a weaker form of 4.8.13 \citep{Grothendieck1990XVIII}. We use that strict-bound theorem for the applications here. It proves Corollary~4.8.14, because an essentially small category allows the cardinal to be enlarged. Thus an essentially small category is accessible exactly when it is idempotent complete; an accessible category with a common presentability bound on all objects is essentially small. The higher-categorical counterpart is \citet[Corollaries~5.4.3.5--5.4.3.6]{Lurie2009HTT}.

The equality case $\pi_0=\cardred C$ is deliberately kept separate in the audit. The checked forward reference and the general small-category result do not settle that fixed-cardinal assertion. We do not describe it as a currently open problem merely because its proof is absent here. What is established for this reading is the strict-bound result, the consequent small-category characterization, and the exact remaining endpoint in the printed proposition.
\end{dcomment}


\subsection*{4.9\quad Le théorème de représentation indicielle pour $\Ind_{\pi}(C)$}

\textbf{4.9.1} (\footnote[93]{\texttt{[Les numéros 4.9.1 à 4.9.6 sont utilisés deux fois, 4.9.3 trois fois.]}}). Soient $C$ une catégorie équivalente à une petite catégorie, $\pi_0$ un cardinal infini,
\[
\cM \ = \ \Ind_{\pi_0}(C) \ \subset\ C^{\wedge}\ , \tag*{(4.9.1)}
\]
$J$ une catégorie équivalente à une petite catégorie. Considérons l'inclusion pleinement fidèle $C \hookrightarrow \cM$, et l'inclusion correspondante de $\mathfrak{U}$-catégories
\[
\uHom(J, C) \ \hookrightarrow\ \uHom(J, \cM) \qquad\quad \text{(pleinement fidèle).} \tag*{(4.9.2)}
\]
On sait, comme $C \subset \cM_{\pi_0}$, que les objets de $\uHom(J, C)$ donnent des objets $\pi_0$-\textit{accessibles} de $\uHom(J, \cM)$, du moins si on suppose
\[
\card\red J \ \leqslant\ \pi_0\ , \tag*{(4.9.3)}
\]
ce que nous ferons désormais. (Cf. 4.5.8, 2°, p. 119.) Comme le deuxième membre de (4.9.2) satisfait $\mathrm{L}_{\pi_0}$ ($\cM$ y satisfait), il s'ensuit que ce foncteur (4.9.2) se prolonge en
\[
\Ind_{\pi_0}(\uHom(J, C)) \ \hookrightarrow\ \uHom(J,\ \underbrace{\cM}_{=\,\Ind_{\pi_0}(C)}\ )\ , \tag*{(4.9.3)}
\]
% foot: Version 17 décembre 2025
\printedpage{78}
% head: A. Grothendieck : Dérivateurs, Ch. XVIII, Édition des Universités de Montpellier II et Paris VII
et ce foncteur est \textit{pleinement fidèle} grâce au résultat d'accessibilité qu'on vient d'énoncer. Je m'intéresse à établir des cas où ce foncteur est essentiellement surjectif, i.e. est une \textit{équivalence} de catégories. Notons que pour que
\[
\varphi \ :\ J \longrightarrow \cM
\]
appartienne à l'image essentielle, il faut et il suffit qu'il existe un foncteur
\[
\Phi \ :\ I \ \longrightarrow\ \uHom(J, C)\ ,\ I \text{ grande devant } \pi_0 \text{ (et on peut supposer } I \text{ ordonnée)}\ , \tag*{(4.9.4)}
\]

\msp{page 218}

tel qu'on ait
\[
\varphi \ \xleftarrow{\ \sim\ }\ \varinjlim_{i\in I} \Phi(i)\ , \tag*{(4.9.5)}
\]
ou ce qui revient au même, tel qu'on ait
\[
\varphi(j) \ \xleftarrow{\ \sim\ }\ \varinjlim_{i\in I} \Phi(i, j)\ , \qquad\quad \text{pour } j \in J\ , \tag*{(4.9.6)}
\]
où on a posé
\[
\Phi(i, j) \ = \ \Phi(i)(j)\ ,
\]
interprétant donc $\Phi$ comme un foncteur
\[
\Phi \ :\ I \times J \ \longrightarrow\ C\ ; \tag*{(4.9.7)}
\]
la loi fonctorielle de $\varphi(j)$ donnée par (4.9.6) étant évidente. En d'autres termes, les $\varphi$ dans l'image essentielle de (4.4.3) sont ceux qui admettent une \textit{représentation indicielle} par un foncteur $\Phi$ sous forme (4.9.4), ou sous forme (4.9.7), avec \textit{de plus} $I$ grand devant $\pi_0$. Je rappelle a ce sujet la % [sic] a ce sujet -> à ce sujet, missing accent

\textbf{Proposition 4.9.2.} \textit{Sous les conditions précédentes sur $C$, $\pi_0$, $J$ (notamment} $\card\red J \leqslant \pi_0$\textit{), si $\varphi : J \longrightarrow \cM \defeq \Ind_{\pi_0}(C)$, pour que $\varphi$ appartienne à l'image essentielle du foncteur pleinement fidèle (4.9.3), il faut et il suffit qu'elle satisfasse les deux conditions suivantes :}

\msp{page 219}

\dgrblock{dgrb1036}{A coherent diagram needs a common index}
\begin{dcomment}{exposition}{p218}{A coherent diagram needs a common index}
Grothendieck's aim is to present an entire diagram by diagrams in the small category, so that its arrows and relations are visible before taking the colimit. Write $[J,C]$ for $\uHom(J,C)$ and assume $|\mathrm{Mor}\,J|\leq\pi_0$. For $A:J\to C$ and a $\pi_0^+$-filtered diagram $(X_i)$ in $[J,\mathcal M]$, the set of natural transformations is the equalizer
\[
 \operatorname{Nat}(A,X)=
 \operatorname{Eq}\left(
 \prod_{j\in\Ob J}\Hom(Aj,Xj)
 \rightrightarrows
 \prod_{u:j\to k}\Hom(Aj,Xk)\right).
\]
Both products have at most $\pi_0$ factors. The defining commutation theorem in sets therefore proves that $A$ is $\pi_0^+$-presentable. The usual Hom formula for ind-objects then gives full faithfulness of
\[
 \Ind_{\pi_0}([J,C])\longrightarrow[J,\Ind_{\pi_0}(C)].
\]
This argument proves uniqueness of morphisms between already indexed presentations. It does not produce the presentation of an arbitrary diagram.

For a diagram $\varphi$, let $K=[J,C]/\varphi$. If $K$ is $\pi_0^+$-filtered and each evaluation $K\to C/\varphi(j)$ is final, the tautological diagram in $[J,C]$ has colimit $\varphi$, because colimits in the target functor category are computed objectwise. Conversely, suppose $\varphi=\varinjlim_{i\in I}A_i$ is already an indexed presentation. Any family of at most $\pi_0$ objects of $K$, together with at most $\pi_0$ relations, factors through one stage $A_i$ by the presentability just proved. Any equality between those factorizations holds at a further common stage. Thus $K$ is $\pi_0^+$-filtered and $I\to K$ is final. The same factorization-and-equality argument applied at a single $j$ shows that $I\to C/\varphi(j)$ is final. Since its factorization through $K$ has final first factor, evaluation from $K$ is final as well. This proves the two directions of 4.9.2 with the size condition in its precise place.
\end{dcomment}


\begin{enumerate}[leftmargin=*]
\item[\textit{a)}] $\uHom(J, C)/\varphi$ \textit{est grande devant $\pi_0$.}
\item[\textit{b)}] \textit{Pour tout $j \in J$, le foncteur}
\[
\uHom(J, C)/\varphi \ \longrightarrow\ C/\varphi(j) \qquad\quad
\left\{
\parbox{0.35\textwidth}{\textit{foncteur entre catégories filtrantes grandes devant $\pi_0$, en vertu de a),}}
\right. \tag*{(4.9.2.1)}
\]
\textit{induit par le foncteur}
\[
\left\{
\begin{array}{rcl}
\uHom(J, C) & \longrightarrow & C \\
F & \longmapsto & F(j)\ ,
\end{array}
\right. \tag*{(4.9.2.2)}
\]
\textit{est} cofinal.
\end{enumerate}
% foot: Version 17 décembre 2025

% ---- from der-ch18_p081-085.tex
\printedpage{79}
% head: A. Grothendieck : Dérivateurs, Ch. XVIII, Édition des Universités de Montpellier II et Paris VII

Le résultat principal dans la présente section est le suivant :

\textbf{Théorème 4.9.3.} \textit{Sous les conditions générales de 4.9.1 pour $C$, $\pi_0$, $J$ (notamment} $\cardred J \leqslant \pi_0$\textit{), le foncteur pleinement fidèle (4.9.3) est une équivalence de catégories dans chacun des deux cas suivants :}
\begin{itemize}[label={}]
{\itshape
\item[a)] $C$ est stable par $\varinjlim$ filtrantes dénombrables.
\item[b)] Il existe un cardinal $\pi_0 < \pi$, et un cardinal $\pi'_0 > \pi_0$ \textup{(}\textbf{\upshape N.B.} On peut supposer que $\pi'_0$ est le successeur de $\pi_0$, donc $\pi'_0 \leqslant \pi$), tels que $C$ soit stable par limites inductives grandes devant $\pi_0$, et de cardinal $\pi'_0$.
}
\end{itemize}
(\footnote[94]{Cas commun : il existe un ordinal limite $\alpha$, tel que 1) $\card \alpha \leqslant \pi_0$. 2) Les limites inductives \texttt{[de]} type $I_\alpha = \{\beta \mid \beta < \alpha\}$ sont représentables dans $C$.}).

\textbf{Remarque 4.9.4.}
\begin{itemize}[label={}]
\item[a)] Ce théorème, cas a), contient le théorème 4.4.13, partie b) (p. 95-96), où on suppose que $C$ est stable par limites inductives de cardinal $\leqslant \pi_0$, l'hypothèse a) de 4.9.3 étant considérablement plus faible.
\item[b)] \msp{page 220}Je présume qu'en dehors d'une hypothèse restrictive du type a) ou b) ci-dessus, le foncteur pleinement fidèle (4.9.3) n'est pas toujours essentiellement surjectif, même dans le cas particulier où $J$ est une catégorie finie, p.ex. $J = B_G$, $G = \{\pm 1\}$. Pour un contre-exemple dans le cas très voisin où on prend $\cM = \Ind(C)$ (au lieu de $\Ind_{\pi_0}(C)$), et où on considère le foncteur (pleinement fidèle si $J$ est une catégorie équivalent à une catégorie \textit{finie}) % [sic] équivalent -> équivalente, agreement with « une catégorie »
\[
\Ind(\uHom(J,C)) \longrightarrow \uHom(J,\Ind(C)),
\]
je renvoie à l'exercice SGA 4 I 8.8.8.
\end{itemize}
L'application principale du théorème 4.9.3 consiste en le

\textbf{Corollaire 4.9.5.} \textit{Soient $\cM$ une catégorie} accessible \textit{(4.8.7), $J$ une catégorie équivalente à une petite catégorie. Alors la catégorie $\cM(J^\circ) = \uHom(J,\cM)$ est accessible. Plus précisément, soient $\pi_0$ un cardinal} utile \textit{pour $\cM$ (i.e. tel que $\Ind_{\pi_0}(\cM_{\pi_0}) \xrightarrow{\sim} \cM$, i.e. un cardinal tel que $\cM$ soit $\pi_0$-accessible), soit $\pi$ un cardinal infini tel que $\pi^{\pi_0} = \pi \geqslant \cardred J$} (\footnote[95]{On a donc $\pi = \pi^{\pi_0} \geqslant 2^{\pi_0} > \pi_0$.}), \textit{i.e. $\pi = c^{\pi_0}$, où $c$ est n'importe quel cardinal $c \geqslant \sup(\pi_0, \cardred J)$, alors $\uHom(J,\cM)$ est $\pi$-accessible, plus précisément on a une équivalence de catégories} % ?? uncertain: font of « alors » -- italic or upright
\[
\uHom(J,\cM) \xleftarrow{\ \sim\ } \Ind_\pi(\uHom(J,\cM_\pi)) \ . \tag*{(4.9.5.1)}
\]

\msp{page 221}On sait en effet, avec les hypothèses faites sur $\pi_0$, $\pi$, que $\pi$ est utile pour $\cM$, i.e. que
\[
\cM \xleftarrow{\ \sim\ } \Ind_\pi(\cM_\pi) \ ,
\]
d'autre part, $\cM_\pi$ est stable par $\varinjlim$ grandes devant $\pi_0$ et de cardinal $\leqslant \pi$. On peut donc appliquer 4.9.3 b), en prenant $C = \cM_\pi$, $\pi$ et $\pi_0$ comme ici (donc $\pi_0 < 2^{\pi_0} \leqslant \pi$), et $\pi'_0 = \pi$.

\dgrblock{dgrb1037}{An involution that no stage can carry}
\begin{dcomment}{development}{p220}{An involution that no stage can carry}
Grothendieck expects the indexed-representation theorem to fail for a finite category with a nontrivial automorphism. His reference to SGA~4, I, 8.8.8 gives the ordinary filtered example. The complete chapter also uses a countable signed-chain category in 4.9.6.8(2), on manuscript page~236. The following transfinite adaptation proves the failure of indexed representation for every infinite cardinal used here. We checked both antecedents and the manuscript's indexed-representation criteria; the cardinal-scaled argument below is supplied for this reading, without a claim of priority.

\textbf{A bounded-filtered involution obstruction}\ours.
For every infinite cardinal $\pi$, there is a small category $C$ and a diagram $B(\mathbb Z/2)\to\Ind_\pi(C)$ outside the essential image of
\[
 \Ind_\pi\bigl([B(\mathbb Z/2),C]\bigr)
 \longrightarrow[B(\mathbb Z/2),\Ind_\pi(C)].
\]
\textit{Proof.} Put $\lambda=\pi^+$ and let the objects of $C$ be the ordinals $\alpha<\lambda$. There are no maps $\alpha\to\beta$ when $\alpha>\beta$, just the identity when $\alpha=\beta$, and two maps, labelled $+1,-1$, when $\alpha<\beta$. Composition multiplies the labels. This is a small category, and its only endomorphisms are identities. The positive maps define a functor $A:\lambda\to C$. The ordinal $\lambda$ is $\lambda$-filtered, since every set of fewer than $\lambda$ ordinals has supremum less than $\lambda$. Therefore
\[
 F=\varinjlim_{\alpha<\lambda}y(A\alpha)
\]
belongs to $\Ind_\pi(C)$. For a fixed $\beta$, the tail $\alpha>\beta$ is cofinal, and $\Hom_C(\beta,\alpha)=\{+1,-1\}$ with identity transition maps on these two labels. Hence $F(\beta)=\{+1,-1\}$. For a morphism labelled $\epsilon$, the presheaf action is multiplication by $\epsilon$. Multiplication by $-1$ consequently defines a natural involution $\sigma:F\to F$, and it has no fixed point in any $F(\beta)$.

A functor $B(\mathbb Z/2)\to C$ must have trivial action, because every endomorphism in $C$ is an identity. A natural map from such an object $\beta$ to $(F,\sigma)$ would be an element of $F(\beta)$ fixed by $\sigma$. No such element exists. Thus
\[
 [B(\mathbb Z/2),C]/(F,\sigma)=\varnothing.
\]
In particular condition a) of 4.9.2 fails, and no indexed presentation can exist. Equivalently, every ind-colimit of the trivial actions still has trivial action, whereas $\sigma\neq1_F$. This proves the assertion.\hfill$\square$

At the ordinary filtered level the construction is the signed-chain version of the unequal two-interval example in SGA~4 \citep[I, Exercise~8.8.8]{SGA4I1972}; the signed-chain category itself also occurs later in this chapter \citep[4.9.6.8(2)]{Grothendieck1990XVIII}. The transfinite chain is essential here: replacing it by the countable chain would not produce an object of $\Ind_\pi(C)$ when $\pi$ is infinite. The obstruction is therefore a direct answer to the cardinal version posed on manuscript page~220, rather than an appeal to a counterexample at the wrong filteredness rank.
\end{dcomment}


\printedpage{80}
% head: A. Grothendieck : Dérivateurs, Ch. XVIII, Édition des Universités de Montpellier II et Paris VII

\textbf{4.9.6} (\footnote[96]{Résultats préliminaires (à la démonstration de 4.9.3).}). Le reste de la présente section 4.9 est consacré à la démonstration de 4.9.3, i.e. à la vérification, pour $J \longrightarrow \cM$, des conditions a), b) de 4.9.2. Dans toute la suite, on se fixe donc une
\[
\varphi \ : \ J \longrightarrow \cM \ = \ \Ind_\pi(C) \ . \tag*{(4.9.6.1)}
\]
On peut supposer $J$ réduite, donc petite et $\card J \leqslant \pi_0$. On va considérer le foncteur déduit du foncteur (4.9.2.1)
\[
\uHom(J,C)/\varphi \longrightarrow \prod_{j \in \Ob J} C/\varphi(j) \ , \tag*{(4.9.6.2)}
\]
dont les foncteurs en question se déduisent en composant avec les foncteurs projection
\[
\prod C/\varphi(j) \longrightarrow C/\varphi(j_0) \ . \tag*{(4.9.6.3)}
\]
Comme les catégories facteurs $C/\varphi(j)$ sont filtrantes (et même grandes devant $\pi$), leur produit est également filtrant, et même grand devant $\pi$ (immédiat), et de même
\msp{page 222}
pour les produits partiels. Il s'ensuit aussitôt que les projections (4.9.6.3) sont cofinales, en vertu du

\dgrblock{dgrb1038}{Why colimits inside the small category repair coherence}
\begin{dcomment}{exposition}{p221}{Why colimits inside the small category repair coherence}
Grothendieck's sufficient hypothesis makes the successive approximations to a diagram converge to an actual diagram in $C$. The countable case has a particularly transparent proof. Fix $\varphi:J\to\Ind_\pi(C)$, with $|\mathrm{Mor}\,J|\leq\pi$, and assume that $C$ has colimits of countable chains. Begin with arbitrary maps $X_0(j)\to\varphi(j)$ from objects of $C$. Since each $C/\varphi(j)$ is $\pi^+$-filtered, one can choose $X_1(j)$ and, for every arrow $u:j\to k$, a lift
\[
 d_0(u):X_0(j)\longrightarrow X_1(k)
\]
of its composite to $\varphi(k)$. The identity arrows give the transition maps $t_0(j)=d_0(1_j)$. Repeat this construction, imposing at the next stage all relations
\[
 d_{n+1}(v)d_n(u)
 =t_{n+1}(\ell)d_n(vu)
 \qquad(j\xrightarrow{u}k\xrightarrow{v}\ell).
\]
Before refinement the two sides have the same image in $\varphi(\ell)$, so they become equal at a further stage of its category of elements. There are at most $\pi$ equations and objects to dominate, and a common refinement imposes all of them. Taking $u=1_j$ gives compatibility of the $d_n(v)$ with the transition maps. Thus their colimits define arrows of a functor $X:J\to C$, with $X(j)=\varinjlim_nX_n(j)$; the displayed relation proves composition and the identity relation proves the identity law.

There is a subtle but necessary final step. Yoneda need not preserve these colimits in $C$. Instead, every $\varphi(j)$, being a $\pi^+$-filtered colimit of representables, carries any existing colimit in $C$ indexed by at most $\pi$ arrows to a limit of sets. Indeed each representable has that property and $\pi^+$-filtered colimits of sets commute with those limits. Consequently the compatible elements of $\varphi(j)(X_n(j))$ define an element of $\varphi(j)(X(j))$, hence a map $X(j)\to\varphi(j)$. These maps form a natural transformation. This proves F~1 of page~226.

To prove F~$3_\pi$, start with $X_0(j)=G(j)$ and the given pointwise maps $u_\alpha(j):F_\alpha(j)\to G(j)$. Besides the composition equations, impose the equations making $G\to X$ natural and making each composite $F_\alpha\to G\to X$ natural. For example the two maps $d_0(v)u_\alpha(j)$ and $t_0(k)u_\alpha(k)F_\alpha(v)$ already agree after mapping to $\varphi(k)$, so they can be equalized at the next refinement. Their number is still at most $\pi$. The resulting $G\to X$ and $F_\alpha\to X$ have exactly the compatibility requested by F~$3_\pi$. The faithful-functor lemma below now proves 4.9.3(a).

The same construction works along any infinite limit ordinal $\delta$ of cardinal at most $\pi$ for which the relevant chain colimits exist in $C$. At a limit stage $\beta<\delta$, choose a cone in each $C/\varphi(j)$ on the preceding chain. At a successor stage, impose the compatibility of the new arrow lifts with every earlier transition and all composition equations. At most $\pi$ data are involved, so the same filteredness argument applies. At the final colimit the equations make the arrow maps independent of the chosen stage, and the preceding preservation argument supplies the map to $\varphi$. This proves the intended two-cardinal case: if $\mu<\pi$ and $C$ has $\mu^+$-filtered colimits of cardinal $\mu^+\leq\pi$, use $\delta=\mu^+$. The theorem's printed b) reuses $\pi_0$ with a different role from its headline; writing the ind-rank as $\pi$ and the smaller rank as $\mu$ removes that ambiguity.

For Corollary~4.9.5 take $C=\mathcal M_\pi$, with $\pi^{\pi_0}=\pi$ and $|\mathrm{Mor}\,J|\leq\pi$. This $C$ is closed under the existing $\pi_0^+$-filtered colimits of cardinal at most $\pi$, and the argument applies. Thus the theorem gives the claimed accessibility at the stated larger cardinal. The later statement that a small diagram category in an accessible category is accessible is therefore a continuation of this precise coherence argument; its higher-categorical counterpart is \citet[Proposition~5.4.4.3]{Lurie2009HTT}.
\end{dcomment}


\textbf{Lemme 4.9.6.1.} \textit{Soient $I$, $I'$ deux catégories avec $I'$ \textup{0}-connexe (p.ex. filtrante), alors le foncteur projection $pr_1 : I \times I' \longrightarrow I$ est cofinal.}

En effet, pour $i' \in \Ob I'$ \texttt{[plutôt} $i \in \Ob I$\texttt{]}, on a
\[
i\backslash(I \times I') \ \simeq \ (i\backslash I) \times I' \ ,
\]
or $i\backslash I$ et $I'$ étant 0-connexe, il en est de même de leur produit.

\textbf{Lemme 4.9.6.2} (\footnote[97]{Attention, il n'est pas vrai qu'un produit infini de catégories 0-connexes est 0-connexe !}). \textit{Soit $(I_\alpha)_{\alpha \in A}$ une famille de catégories} filtrantes \textit{(resp. grandes devant $\pi$). Alors $I_A \stackrel{\text{déf}}{=} \prod_{\alpha \in A} I_\alpha$ est filtrante (resp. grande devant $\pi$), et pour toute partie $B \subset A$, la projection canonique $I_A \longrightarrow I_B$ est cofinal.} % [sic] cofinal -> cofinale, agreement with « la projection »

(\footnote[98]{Cf. démonstration élégante ci-dessous, 4.9.6.6.}). La stabilité de la propriété “filtrante” ou “grande devant $\pi$” par AQT \texttt{[âne qui trotte]}. On a, si $B' = A\backslash B$, $I_A \simeq I_B \times I_{B'} \longrightarrow I_B$, et on applique le lemme 4.9.6.1 pour conclure que ce foncteur est cofinal, $I_{B'}$ étant filtrante, donc 0-connexe.

\msp{page 223}Le foncteur (4.9.6.3) étant cofinal, pour établir que (4.9.2.1) l'est, il suffit de prouver qu'il en est ainsi du foncteur (4.9.6.2), ce que nous ferons. Notons de plus que ledit foncteur est \textit{fidèle}, car il s'insère dans le carré commutatif
\[
\begin{tikzcd}[column sep=large, row sep=large]
\uHom(J,C)/\varphi \arrow[r] \arrow[d] & \prod C/\varphi(j) \arrow[d] \\
\uHom(J,C) \arrow[r] & C^{\Ob J}\ ,
\end{tikzcd}
\]
\printedpage{81}
% head: A. Grothendieck : Dérivateurs, Ch. XVIII, Édition des Universités de Montpellier II et Paris VII
où les flèches verticales et la deuxième flèche horizontale sont fidèles. En définitive, pour établir le théorème 4.9.3, nous allons étudier le foncteur \textit{fidèle} (4.9.6.2)
\[
\begin{aligned}
&f \ : \ \uHom(J,C)/\Phi \ \longrightarrow \ \textstyle\prod_{j \in J_0} C/\varphi(j) \\ % [sic] \Phi -> \varphi, the functor is written \varphi everywhere else
&\text{(où on a posé } J_0 = \Ob J,\ J_1 = \Fl J) \ ,
\end{aligned}
\tag*{$(*)$}
\]
et nous allons prouver :
\begin{itemize}[label={}]
\item[\textcircled{a}] Ce foncteur fidèle $(*)$ est \textit{cofinal.}
\item[\textcircled{b}] La catégorie source $\uHom(J,C)/\varphi$ est grande devant $\pi_0$.
\end{itemize}
Nous allons pour ceci dégager un lemme auxiliaire sur la cofinalité des foncteurs \textit{fidèles} :

\dgrblock{dgrb1039}{Connected factors can require unbounded zigzags}
\begin{dcomment}{development}{p222}{Connected factors can require unbounded zigzags}
Grothendieck's warning in footnote~97 explains why the product argument must use filteredness, not merely connectedness. We make the obstruction quantitative; this also gives a concrete test for an infinite-product argument elsewhere in the reader. We checked the product assertions in 4.9.6 and the standard filtered-category criteria in SGA~4, I, 8.1. The following exact component calculation is supplied here as a small development, without a priority claim.

\textbf{Components of a product of fences}\ours.
Let $P_n$ have objects $0,1,\ldots,n$ and, besides identities, one arrow along each adjacent edge, always directed from the even vertex to the odd vertex. Then every $P_n$ is connected, whereas two objects $a=(a_n)$ and $b=(b_n)$ of $\prod_{n\geq1}P_n$ are connected if and only if
\[
 \sup_{n\geq1}|a_n-b_n|<\infty.
\]
Consequently the product has exactly $2^{\aleph_0}$ connected components.

\textit{Proof.} There are no composable nonidentity arrows in $P_n$, so its underlying undirected graph is the path with $n$ edges. A zigzag of length $r$ in the product projects to a zigzag of length $r$ in every factor; therefore $|a_n-b_n|\leq r$ for all $n$. Conversely, suppose all these distances are at most $r$. Choose the unique shortest path in each factor. Turn it into a zigzag with the fixed orientation pattern forward, backward, forward, backward, and so on, by inserting identity arrows when the next edge has the opposite direction. At most two positions are needed for each edge, and further identities pad every path to length $2r$. The resulting coordinatewise arrows are arrows of the product, so they give one finite zigzag from $a$ to $b$.

There are at most continuum many objects in the product. For every real $t\in[0,1]$, take $a_n(t)=\lfloor tn\rfloor$. If $s\neq t$, the distances $|a_n(s)-a_n(t)|$ are unbounded, hence these objects lie in distinct components. This gives the matching lower bound.\hfill$\square$

Filtered categories avoid this defect because any finite family of objects has a common target, and any parallel pair can be equalized. For a product of $\pi^+$-filtered categories, choose a cocone coordinatewise for each diagram of at most $\pi$ arrows. Choice assembles these cones into a cocone in the product, including the empty-diagram case. Every projection to a subproduct is final: its comma categories are products of a category with an initial object and the complementary filtered product. This proves 4.9.6.1--4.9.6.2 while accounting for the warning.
\end{dcomment}


\msp{page 224}\textbf{Lemme 4.9.6.3} (\footnote[99]{Ce lemme fait partie du sorite général sur la cofinalité et les catégories filtrantes.}). \textit{Soit $f : I \longrightarrow I'$ un foncteur, et considérons les conditions :}
{\itshape
\begin{itemize}[label={}]
\item[a)] $I$ filtrante, $f$ cofinal, i.e. (en présence de $I$ filtrante) $f$ satisfait les conditions familières \rbox{F 1}, \rbox{F 2}.
\item[b)] $I'$ est filtrante, et les $i'\backslash I$ (pour $i' \in \Ob I'$) sont filtrantes.
\item[b$'$)] $I'$ \textup{0}-connexe, satisfait PF (2) \textup{(}\footnote[100]{p.ex. PF 2 est satisfaite si $I'$ est \textit{ordonnée}, sans avoir à supposer d'avance $I'$ filtrante, seulement 0-connexe.}\textup{)} (égalisation des double-flèches).
\item[b$''$)] $I'$ est filtrante, les $i'\backslash I$ pseudo-filtrantes, i.e. satisfont à la condition PF $1'$.
\item[c)] $I'$ filtrante, et $f$ satisfait à F 1, ainsi que la condition suivante.
\begin{itemize}[label={}]
\item[\rbox{F 3}] Pour tout couple $(i_0,i_1)$ d'objets de $I$ et toute flèche $f(i_0) \xrightarrow{u'} f(i_1)$, il existe un objet $i_2$ de $I$, et une flèche $v' : f(i_1) \longrightarrow f(i_2)$, tels que $v'$ et $v'u'$ soient relevables en des flèches $i_1 \xrightarrow{v} i_2$ et $i_0 \xrightarrow{w} i_2$ (i.e. il existe $i_2$ et $v : i_1 \longrightarrow i_2$, $w : i_0 \longrightarrow i_2$ tels qu'on ait $f(v)u' = f(w)$).
\end{itemize}
\end{itemize}
\[
\begin{tikzcd}
i_0 \arrow[rr, bend left=20, "w"] & i_1 \arrow[r, "v"] & i_2 % ?? uncertain: line style of the arrow v -- solid or densely dotted
\end{tikzcd}
\]
\[
\begin{tikzcd}
f(i_0) \arrow[r, "u'"] \arrow[rr, dotted, bend right=20, "f(w)"'] & f(i_1) \arrow[r, "f(v) = v'"] & f(i_2)
\end{tikzcd}
\]
\begin{itemize}[label={}]
\item[1)] On a les implications
\[
a \iff b \iff b' \implies b'' \iff c \ ,
\]
et si $f$ est \textup{fidèle} (p.ex. $I$ ordonnée), les cinq conditions a), b), b$'$), b$''$), c) sont équivalentes.
\end{itemize}
}

\printedpage{82}
% head: A. Grothendieck : Dérivateurs, Ch. XVIII, Édition des Universités de Montpellier II et Paris VII
{\itshape
\begin{itemize}[label={}]
\item[2)] \msp{page 225}Soit $\pi$ un cardinal infini, et considérons les conditions :
\begin{itemize}[label={}]
\item[$a_\pi$)] $I$ grande devant $\pi$, $f$ cofinal, i.e. F 1, F 2.
\item[$b_\pi$)] $I'$ grande devant $\pi$, les $i'\backslash I$ ($i' \in \Ob I'$) aussi.
\item[$b_\pi{}''$)] $I'$ grande devant $\pi$, les $i'\backslash I$ ($i' \in \Ob I'$) sont $\pi$-pseudo-filtrantes, i.e. satisfont PS $1_\pi{}'$. % [sic] PS 1_π' -> PF 1_π', the conditions are elsewhere called PF
\item[$c_\pi$)] $I'$ grande devant $\pi$, et $f$ satisfait F 1 et
\begin{itemize}[label={}]
\item[\rbox{F $3_\pi$}] \textup{(}\footnote[101]{\textbf{N.B.} La condition F $3_\pi$ est tautologiquement satisfaite si $f$ est \textit{pleinement fidèle.}}\textup{)}. Pour toute famille $(i_\alpha)_{\alpha \in A}$ d'objets de $I$, avec \textup{card} $A \leqslant \pi$, et $i \in \Ob I$, et une famille de flèches $u'_\alpha : f(i_\alpha) \longrightarrow f(i)$, existe $j \in \Ob I$ et une flèche $v' : f(i) \longrightarrow f(j)$, telle que $v'$ \textup{et} les $v'u'_\alpha$ proviennent de flèches de $I$, i.e. il existe des flèches $v : i \longrightarrow j$, $w_\alpha : i_\alpha \longrightarrow j$ telles que $f(v)u'_\alpha = f(w_\alpha)$ $\forall\ \alpha \in A$.
\end{itemize}
\end{itemize}
\[
\begin{tikzcd}[row sep=huge]
i_\alpha \arrow[d, dash] \arrow[rr, dotted, bend left=20, "w_\alpha"] & i \arrow[r, dotted, "v"] \arrow[d, dash] & j \arrow[d, dash] \\
f(i_\alpha) \arrow[r, "u'_\alpha"] & f(i) \arrow[r, "v' = f(v)"] & f(j)
\end{tikzcd}
\]
Ceci posé, on a
\[
a_\pi \iff b_\pi \implies c_\pi \implies b_\pi{}'' \ ,
\]
et si $f$ est fidèle, ou $I$ ordonné, les quatre conditions $a_\pi$, $b_\pi$, $b_\pi{}''$, $c_\pi$ sont équivalentes.
\end{itemize}
}
(\footnote[102]{Démonstration de la proposition donnée dans 4.9.6.7, page 227.}).

\textbf{4.9.6.4. N.B.} Dans le cas qui nous intéresse, où $f$ est le foncteur $(*)$ (p. 223), on sait que $I'$ est grande devant $\pi$. On voudrait établir que $I$ l'est et que $f$ est cofinal, i.e. la condition $\mathrm{a}_\pi$. On le prouvera en prenant $\mathrm{b}_\pi$ (qui implique $\mathrm{a}_\pi$ puisque $f$ est fidèle). Il reste donc à
\msp{page 226}
prouver, sur le foncteur $f$ dans $(*)$, p. 223, les deux conditions suivantes :
\begin{itemize}[label={}]
\item[F 1)] Pour tout système de flèches dans $\cM$
\[
X(j) \xrightarrow{p(j)} \varphi(j) \ , \qquad j \in J_0 = \Ob J, \text{ les } X(j) \text{ dans } \Ob C \ ,
\]
il existe un foncteur $F \ : \ J \longrightarrow C$ et $F \xrightarrow{q} \varphi$, enfin un système de flèches $X(j) \xrightarrow{u(j)} F(j)$ dans les $C/\varphi(j)$, i.e. rendant commutatif les diagrammes % [sic] commutatif -> commutatifs, agreement with « les diagrammes »
\[
\begin{tikzcd}[column sep=small]
X(j) \arrow[rr, "u(j)"] \arrow[dr, "p(j)"'] & & F(j) \arrow[dl, "q(j)"] \\
& \varphi(j)\ . &
\end{tikzcd}
\]
\item[F $3_\pi$)] Pour tout système $(F_\alpha)_{\alpha \in A}$ ($\card A \leqslant \pi$) d'éléments dans $\uHom(C,J)$, et des flèches % [sic] \uHom(C,J) -> \uHom(J,C), the F_α are functors J -> C
\[
F_\alpha \xrightarrow{p_\alpha} \varphi \qquad \text{dans } \uHom(J,\cM),
\]
et pour $G$ dans $\uHom(J,C)$ et $G \xrightarrow{q} \varphi$ dans $\uHom(J,\cM)$, étant donné un système de flèches $u_\alpha(j)$,
\printedpage{83}
% head: A. Grothendieck : Dérivateurs, Ch. XVIII, Édition des Universités de Montpellier II et Paris VII
\[
\begin{tikzcd}[column sep=small]
F_\alpha(j) \arrow[rr, "u_\alpha(j)", "\text{comm.}"'] \arrow[dr, "p_\alpha(j)"'] & & G(j) \arrow[dl, "q(j)"] \\
& \varphi(j)\ , &
\end{tikzcd}
\qquad u_\alpha(j) \text{ flèche dans } C/\varphi(j) \ ,
\]
il existe $H \xrightarrow{q} \varphi$ dans $\uHom(J,C)$ et une flèche $G \xrightarrow{v} H$ dans $\uHom(J,C)/\varphi$, telle que pour tout $\alpha \in A$, le système de flèches composés
\msp{page 227}
\[
\begin{tikzcd}[row sep=huge]
F_\alpha(j) \arrow[r, "u_\alpha(j)", "\text{comm.}"'] \arrow[dr, dotted, "p_\alpha(j)"'] & G(j) \arrow[r, "v(j)", "\text{comm.}"'] \arrow[d, dotted, "q(j)"] & H(j) \arrow[dl, dotted, "q(j)"] \\
& \varphi(j) &
\end{tikzcd}
\]
(pour $j \in J_0 \in \Ob J$) provienne d'un \textit{foncteur} $w_\alpha : F_\alpha \longrightarrow H$. % [sic] J_0 \in \Ob J -> J_0 = \Ob J, J_0 was defined as Ob J
\end{itemize}
Pour prouver 4.9.3, nous sommes donc ramené à prouver les deux propositions F 1), F $3_\pi$) précédentes, sous l'une des deux conditions a) ou b) de 4.9.3.

\dgrblock{dgrb1040}{The faithful cofinality test, with its hypotheses restored}
\begin{dcomment}{exposition}{p225}{The faithful cofinality test, with its hypotheses restored}
Grothendieck reduces the construction of a coherent presentation to the ability to lift a bounded family of arrows after one common refinement. Faithfulness then turns equations seen downstairs into actual equations upstairs. This is the precise role of F~$3_\pi$; the evaluation functor into the product of the categories of elements is faithful because a natural transformation is determined by its components.

The printed condition b$'$ in 4.9.6.3 is incomplete. It only says that $I'$ is connected and equalizes parallel pairs, so it cannot imply an assertion about $I$ or $f$. The functor $\varnothing\to1$ satisfies it and is not final. The proof in the continuation uses, in addition, that every $(i'\downarrow f)$ is filtered; inserting that clause gives the intended b$'$. The assertions involving b) and the faithful test c) are correct. Here is a proof of the version needed in these pages.

Suppose $f:I\to I'$ is faithful, $I'$ is $\pi^+$-filtered, F~1 holds, and F~$3_\pi$ holds. Given a diagram $D:K\to I$ with at most $\pi$ arrows, choose a cone $u_k:fDk\to z$ in $I'$, and then a map $z\to f(i)$ by F~1. F~$3_\pi$ gives $v:i\to h$ and $w_k:Dk\to h$ whose images are the composites of those cone maps with $f(v)$. For an arrow $a:k\to\ell$ of $K$ we have
\[
 f(w_\ell D(a))=f(w_k).
\]
Faithfulness gives $w_\ell D(a)=w_k$, so $I$ is $\pi^+$-filtered. Its nonemptiness follows from F~1 and nonemptiness of $I'$.

Apply the same construction to a bounded diagram of objects $(i_k,p_k:z\to f(i_k))$ in $(z\downarrow f)$. In $I'$ take a cone on its image together with $z$ and all $p_k$. After F~1 and F~$3_\pi$, the maps $w_k$ obtained above satisfy
\[
 f(w_k)p_k=p
\]
for one map $p:z\to f(h)$. They are therefore a cone in the comma category. The empty case is F~1. Every comma category is $\pi^+$-filtered, hence nonempty and connected, so $f$ is final. This proves c$_\pi\Rightarrow$b$_\pi\Rightarrow$a$_\pi$ under faithfulness, including all quantifiers used in 4.9.6.4.

Faithfulness cannot simply be discarded. For a nontrivial group $G$, the functor $BG\to1$ satisfies F~1 and every F~$3_\pi$, since all equations in the target are automatic and arbitrary arrows upstairs can be chosen. But $BG$ is not filtered: distinct automorphisms cannot be equalized. This example identifies the exact step that fails when equality of images is substituted for equality of arrows.
\end{dcomment}


\textbf{Corollaire 4.9.6.5} (de 4.9.6.3). \textit{Soit $(f_\alpha : I_\alpha \longrightarrow I'_\alpha)$ une famille de foncteurs cofinaux, les catégories source étant filtrantes. Alors le foncteur produit $f = \prod f_\alpha : I = \prod_\alpha I_\alpha \longrightarrow I' = \prod_\alpha I'_\alpha$ est cofinal.}

En effet, il faut vérifier que les catégories $i'\backslash I$, pour $i' \in \Ob I'$, sont 0-connexes. Or si $i' = (i'_\alpha)_{\alpha \in A}$, on a
\[
i'\backslash I \ \simeq \ \prod_\alpha (i'_\alpha\backslash I_\alpha) \ ,
\]
et par 4.9.6.3, 1°) a $\Longrightarrow$ c, les $i'_\alpha\backslash I_\alpha$ sont filtrantes, donc leur produit est filtrant
\msp{page 228}
(cf. 4.4.6.2), donc 0-connexe, q.e.d. % [sic] 4.4.6.2 -> 4.9.6.2, the lemma on products of filtrantes is 4.9.6.2

\textbf{Corollaire 4.4.6.6.} \textit{Soit $\pi$ un cardinal infini. Tout produit de catégories grandes devant $\pi$ est grand devant $\pi$.} % [sic] 4.4.6.6 -> 4.9.6.6, numbering of the section and footnote 98

Soit $I = \prod_{\alpha \in A} I_\alpha$, soit $B$ un ensemble de cardinal $\leqslant \pi$, il faut prouver que $I \xrightarrow{\mathrm{diag}_I} I^B$ est cofinal. Or ce foncteur s'identifie au produit des foncteurs similaires $I_\alpha \longrightarrow I_\alpha^B$, lesquels sont cofinaux, les catégories source étant de plus filtrantes. Donc on peut conclure par le corollaire 4.9.6.5.

\textbf{4.9.6.7.} \textsc{Démonstration de} 4.9.6.3. (Partie 1°.)

a $\Longrightarrow$ b. On sait déjà que a) implique que $I'$ est filtrante (SGA 4 I 8.1.3 b). Reste à prouver que pour tout $i' \in \Ob I'$, $i'\backslash I$ est filtrante (\footnote[103]{Court-circuit : on sait déjà que $i'\backslash I$ est 0-connexe, d'autre part elle est \textit{localement} filtrante, car coétale sur $I$ qui est filtrante, donc elle est filtrante.}). Deux conditions PF $1'$ et PF 2 à vérifier (en plus du fait que $i'\backslash I \neq \emptyset$, qui est connu par cofinalité de $f$, impliquant que $i'\backslash I$ est 0-connexe).
\begin{itemize}[label={}]
\item[1)] Soient $(i,u')$, $(j,v')$ deux éléments de $i'\backslash I$, il faut voir qu'ils sont majorés par un troisième élément. Or $I$ étant filtrante, $i$, $j$ sont majorés par $k$. On aurait fini donc si on savait que $f(u_1)u' = f(v_1)v'$, soit $w'$, car alors $(k,w')$ ferait l'affaire.
\end{itemize}


\dgrblock{dgrb1041}{Completing the argument of manuscript page 228}
\begin{dcomment}{exposition}{p228}{Completing the argument of manuscript page 228}
On manuscript page~228 Grothendieck's proof has reached the point just before the two candidate arrows from $i'$ are equalized. The remaining step is short once finality is used in the correct comma category. If $I$ is filtered and $f$ is final, every $(i'\downarrow f)$ is connected. Given two of its objects, join them by a finite zigzag. Project that zigzag to $I$ and take a cone there. Commutativity along each arrow of the zigzag makes the composites from $i'$ agree at its vertex, so that vertex is an object receiving maps from the two original comma objects. Parallel arrows in the comma category can be equalized by filteredness of $I$; the resulting arrow is still over $i'$. Hence the comma category is filtered. This finishes the particular step begun in 4.9.6.7.

For the stronger cardinal statement, begin with a diagram of at most $\pi$ arrows in $(i'\downarrow f)$. Choose one object and a finite connecting zigzag to each of the others. The union of their projections and of the original diagram still has at most $\pi$ arrows, because $\pi$ is infinite. A cone in $I$ simultaneously imposes all relations and all equalities of maps from $i'$. Thus every comma category is $\pi^+$-filtered. Likewise a bounded diagram in $I'$ can be lifted object by object using F~1; connect, inside the appropriate comma categories, the two choices associated with each arrow. A cone on the resulting bounded diagram in $I$ yields a cone in $I'$. This proves a$_\pi\Rightarrow$b$_\pi$.

Conversely, suppose $I'$ and all $(i'\downarrow f)$ are $\pi^+$-filtered. Finality is immediate. First derive F~$3_\pi$. Given $u'_\alpha:f(i_\alpha)\to f(i)$, filteredness of $(f(i_\alpha)\downarrow f)$ supplies arrows $w_\alpha:i_\alpha\to k_\alpha$ and $v_\alpha:i\to k_\alpha$ with $f(w_\alpha)=f(v_\alpha)u'_\alpha$. Regard all $v_\alpha$ as arrows from $(i,1)$ in $(f(i)\downarrow f)$. A cone on this bounded diagram supplies maps $h_\alpha:k_\alpha\to h$ with all $h_\alpha v_\alpha$ equal to one $v:i\to h$. The arrows $h_\alpha w_\alpha$ prove F~$3_\pi$.

Now choose a cone on $fD$ in $I'$, follow its vertex by a map to some $f(i)$, and apply F~$3_\pi$. This gives arrows $w_k:Dk\to h$ which satisfy all the cone equations after applying $f$. For each arrow $a:k\to\ell$, the two arrows $w_\ell D(a),w_k$ are parallel arrows in $(fDk\downarrow f)$. Filteredness there gives $e_a:h\to h_a$ equalizing them. In $(f(h)\downarrow f)$ take a cone on the bounded diagram consisting of $(h,1)$ and these arrows $e_a$. Its map $e:h\to H$ factors through every $e_a$, so the maps $ew_k$ satisfy every cone equation in $I$. This proves $\pi^+$-filteredness of $I$ without faithfulness, and establishes b$_\pi\Rightarrow$a$_\pi$.

Finally, for final functors $f_\alpha:I_\alpha\to I'_\alpha$ with filtered sources, the comma category of the product functor is
\[
 ((i'_\alpha)_\alpha\downarrow\prod_\alpha f_\alpha)
 \cong\prod_\alpha(i'_\alpha\downarrow f_\alpha).
\]
Each factor is filtered by the argument just given; their product is filtered and therefore connected. This proves 4.9.6.5. Applying it to the diagonal functors $I_\alpha\to I_\alpha^B$ gives the diagonal proof of the next corollary, whose heading is misnumbered 4.4.6.6 in the edition. The direct coordinatewise cone proof gives the same result and avoids treating infinite products of merely connected categories as connected.
\end{dcomment}
\begin{otherlanguage}{english}
\dgroverview{XVIII-b}{Chapter XVIII \textperiodcentered\ \textup{XVIII-b} \textperiodcentered\ Overview}
\begin{otherlanguage*}{english}
These pages belong to the ordinary categorical foundations of the theory of derivators. Grothendieck is asking how a large category, and then its categories of diagrams, can be recovered from controlled small data. The recurring difficulty is that a bound on the size of a presentation is not yet an invariant of the object presented. The passage from one to the other requires both a factorization argument and an account of the relations which that argument introduces. Retracts, categories of elements, and changes of indexing category are therefore the central constructions of this part of Chapter~XVIII.

There is a cardinal convention to fix at the outset. A category called \emph{grande devant $\pi$} admits cocones on diagrams with at most $\pi$ arrows. In the customary notation for regular cardinals it is $\pi^+$-filtered. An object called $\pi$-accessible in the manuscript is consequently $\pi^+$-presentable. The parameter $\pi$ itself need not be regular. We retain the manuscript's indices, including its later adjective ``regular'' for a useful cardinal satisfying a particular presentation bound. That adjective is a separate definition, not a statement about the set-theoretic cofinality of $\pi$.

The opening, which continues Proposition~4.5.13 from the preceding pages, compares two bounded-presentation filtrations with the intrinsic filtration $\mathcal M_\pi$. In a cocomplete category the basic calculation is straightforward: a colimit indexed by a category with at most $\pi$ arrows, of $\pi$-accessible objects, is again $\pi$-accessible, because a limit of that size in sets commutes with $\pi^+$-filtered colimits. The converse requires generation. When every object is a sufficiently filtered colimit of objects of a small subcategory, the identity of a presentable object factors through one stage. What follows is a retract of a stage. Lemma~4.5.14 suppresses this last distinction. The commentary gives the corrected implication and an explicit counterexample in sets to its unrestricted wording. Its later application to a subcategory closed under bounded colimits remains valid, since such a subcategory also contains the required idempotent splittings.

The discussion of universal strict epimorphic families has a different purpose. It seeks hypotheses under which one bounded presentation already contains the relations needed for descent. The printed direct-product example in groups is defective even after the accompanying footnote: the two coordinate subgroups of $C_2\times C_2$ do not form a strict epimorphic family. Replacing the product by the free product gives the intended example, including failure of epimorphy after pullback. We distinguish this correction from the unresolved proof transition in Corollary~4.5.16, where a universal epimorphic hypothesis is used as though it were universally strictly epimorphic. The strengthened version has a complete proof. The argument given in these pages does not establish the weaker printed version or the entire deduction of case~C of Proposition~4.5.13.

Section~4.6 builds a small category closed under the colimits needed for reconstruction. The qualitative transfinite construction is sound: at each stage there is a set of diagrams and hence a set of chosen colimits; regularity of the stopping ordinal ensures that every bounded diagram appears early enough to be treated. The numerical estimate for the number of new morphisms is a separate assertion. It is false for an arbitrary small starting category, as the source's own later footnotes suspect. A finitary variety with many binary idempotent operations provides a direct counterexample. Thus the existence of a small reconstruction category must be retained while the claimed sharp bound requires another argument. The normalization of a finite initial cardinal in Corollary~4.6.2 repairs an additional elementary edge case, but does not repair that proof.

The resulting cocomplete recognition theorem is the locally presentable theory: a $\pi_0$-paratopos is, with the stated hypotheses restored, locally $\pi_0^+$-presentable. The presheaf reflection and the small-category presentation explain why suitable filtered colimits commute with bounded limits. This places the discussion alongside Gabriel--Ulmer duality and the systematic accounts of accessible categories, without asserting that Grothendieck was drawing on those accounts here \citep{GabrielUlmer1971,MakkaiPare1989,AdamekRosicky1994}. His immediate reference is SGA~4. One of its assertions, that full faithfulness of a filtered-completion functor forces the original objects to be accessible, fails at the stated generality; the rigid real ring supplies the manuscript's counterexample. The intended Ind application instead follows by evaluation of presheaves.

Exactness does not fix the generating cardinal. To answer footnote~59, we construct, for every prescribed infinite $\pi_0$, a sheaf topos whose only $\pi_0^+$-presentable object is initial, although filtered colimits are exact and a small dense category exists. The construction uses a set of cardinal $\pi_0^{++}$ with the topology of complements of subsets of smaller cardinality. Every nonempty support has an increasing $\pi_0^+$-indexed cover by proper opens; restricting a sheaf to this cover rules out presentability at the proposed rank. This is a supplied development of the question, with a complete proof and no claim of priority. The closing examples of \S4.6 also separate the conditions on images, unions and pushouts of monomorphisms. Pointed sets satisfy all six conditions under discussion and give an example beyond abelian categories and topoi.

Section~4.7 removes the assumption of cocompleteness. Now $\mathcal M=\Ind_{\pi_0}C$, and the primed filtration denotes closure under repeated bounded filtered colimits; its meaning must not be imported from the preceding section. The cardinal equation $\pi^{\pi_0}=\pi$ pays for adjoining all the small witnesses needed to make a subcategory of an indexing category sufficiently filtered. We spell out this closure operation, including the old arrows and cone legs omitted from the printed construction. The resulting comparison of presentation size and presentability is the ordinary version of the modern change-of-cardinal arguments \citep[\S3]{Vasey2020Accessible}.

The retract question which Grothendieck leaves at the centre of this comparison has a positive answer in the literature. Chen's bounded closure theorem, and his formulation of the Makkai--Par\'e retract-removal result, show that iterating the specified bounded filtered colimits does not enlarge the resulting class \citep[Corollaries~4.13 and~4.15]{Chen2021Sifted}. We also give an elementary retract argument which preserves the number of arrows. Consequently $\Filt^\pi=\Filt'^\pi$ at every permitted cardinal. At the initial cardinal the common value is $\Kar(C)$, not necessarily $C$: a two-arrow idempotent category is sufficiently filtered at every rank and realizes a nonrepresentable retract. This endpoint explains several otherwise misleading assertions in the text.

The generalized filtration axioms and their useful cardinals require an additional distinction. A filtered slice is a presentation only when its canonical cocone has the correct colimit. Density supplies that condition, but the axioms printed in 4.7.7 do not force it at every level. We exhibit a filtration of sets whose initial slices are terminal and whose initial realization nevertheless misses every nonempty set. This disproves the bare converse in 4.7.12 and identifies the missing hypothesis. Once usefulness is understood through the realization equivalence, every useful cardinal satisfies the quantitative ``regularity'' requirement of 4.7.13. The commentary proves this with a size-controlled cofinal directed poset, so no example of a useful but nonregular cardinal exists under that interpretation. Conversely, the category of nonempty $\aleph_1$-directed posets has unboundedly many nonuseful manuscript cardinals: every singular $\theta>\aleph_1$ of countable cofinality gives failure of $\theta^+$-accessibility. We provide the complete argument, refining the rank choices in \citet[Example~3.8]{Vasey2020Accessible}. This concerns the accessibility spectrum, rather than the presentability rank of one object \citep{LiebermanRosickyVasey2019Sizes}.

Section~4.8 returns to density and gives the conceptual recognition of accessible categories. The author's initial doubt about enlargement of a dense subcategory is eventually withdrawn, correctly: naturality against the smaller dense category determines the missing components. The characterization of the smallest presentable part as an idempotent completion is likewise valid. Two other assertions need more substantial correction. First, Lemma~4.8.11 omits ``pro'' from the conclusion supplied by SGA~4, I, 8.13.2(b). An accessible finite-limit-preserving functor need only be prorepresentable. Locally constant functions from the Cantor space give an explicit counterexample to representability; preservation of all small limits gives the intended representability theorem. Second, the conjecture in footnote~91 that accessibility of a category and its opposite forces essential smallness is false. Sets and partial injections form a self-dual accessible category with objects of every small cardinality \citep[\S3.2]{Heunen2013L2}; we prove the precise presentability criterion directly.

The small-category theorem requires a more restrained verdict. The argument in these pages does not establish the equality endpoint in 4.8.13. The later manuscript proves a strict cardinal inequality, sufficient for Corollary~4.8.14, and explicitly calls it a weaker result. The continuity result checked in \citet{Tendas2022Continuity} also has a strict threshold. We therefore retain the theorem on essentially small accessible categories while leaving that exact endpoint uncertified; this is not a claim that the endpoint is an established open problem today.

Section~4.9 explains why a diagram of ind-objects need not come from diagrams at individual stages. Coherence can impose a genuine obstruction, even for a one-object category with an involution. We extend the SGA~4 counterexample to every infinite filtering bound using a transfinite signed-chain category. The resulting ind-object has an involution which no stage admits. Under the appropriate colimit hypotheses on $C$, Grothendieck's common-index construction repairs the obstruction. The final cofinality lemmas express that construction through simultaneous factorization and equalization. We restore the missing hypotheses in their printed alternatives and complete the particular proof under way on manuscript page~228.

The later homotopical setting preserves this architecture while replacing sets of maps by mapping spaces and strict diagrams by coherent diagrams. Lurie's recognition theorem characterizes presentable infinity-categories through accessibility, cocompleteness and accessible presheaf localization; his diagram-category theorem supplies the corresponding stability under small diagrams \citep[Theorem~5.5.1.1 and Proposition~5.4.4.3]{Lurie2009HTT}. These are structural counterparts of the constructions here. They explain the relevance of this long ordinary-categorical preparation to the work on derivators without identifying strict diagrams in a homotopy category with coherent diagrams in a homotopy theory.
\end{otherlanguage*}
\end{otherlanguage}
\dgrchunkend
